\documentclass[11pt,a4paper]{article}

\usepackage[utf8]{inputenc}
\usepackage[
  left=1.05in,
  right=1.05in,
  top=1in,
  bottom=1in,
  headheight=14pt
]{geometry}
\usepackage{amsmath,amssymb,amsthm,mathtools}
\usepackage{bm}
\usepackage{enumitem}
\usepackage{booktabs,array}
\usepackage{microtype}
\usepackage{xcolor}
\usepackage[breakable,skins]{tcolorbox}
\usepackage[hidelinks]{hyperref}
\usepackage{fancyhdr}
% Start a new page unless #1 of vertical space remains.
\makeatletter
\newcommand{\Needspace}[1]{\par\penalty-100\begingroup
  \dimen@=#1\relax
  \dimen@ii=\pagegoal\advance\dimen@ii by -\pagetotal
  \ifdim\dimen@>\dimen@ii\ifdim\dimen@ii>\z@\vfil\fi\break\fi
  \endgroup}
\makeatother

\linespread{1.25}
\setlength{\parindent}{0pt}
\setlength{\parskip}{1em plus 0.25em minus 0.1em}
\setlength{\abovedisplayskip}{1.2em plus 0.3em minus 0.2em}
\setlength{\belowdisplayskip}{1.2em plus 0.3em minus 0.2em}
\setlength{\abovedisplayshortskip}{0.9em plus 0.2em}
\setlength{\belowdisplayshortskip}{1em plus 0.2em}
\setlength{\jot}{9pt}
\setlength{\emergencystretch}{2em}
\raggedbottom

\setlist[itemize]{leftmargin=*,itemsep=6pt,topsep=6pt,parsep=3pt}
\setlist[enumerate]{leftmargin=*,itemsep=6pt,topsep=6pt,parsep=3pt}

% Plain black frame for the question statement.
\newtcolorbox{qbox}{
  enhanced,
  colback=white,colframe=black,
  boxrule=0.6pt,sharp corners,
  left=10pt,right=10pt,top=4pt,bottom=6pt,
  before skip=14pt,after skip=10pt
}

% Revision block: the facts needed for the question, marked by a rule
% on the left only.
\newtcolorbox{remember}{
  enhanced,
  colback=white,frame hidden,
  borderline west={1.4pt}{0pt}{black},
  sharp corners,
  left=10pt,right=6pt,top=2pt,bottom=2pt,
  before skip=6pt,after skip=12pt
}
\newenvironment{rem}
  {\begin{remember}\noindent\textbf{Remember}\par\nobreak\vspace{-0.3em}}
  {\end{remember}}

\newcounter{qn}
% \Q{title}{type}{marks}
\newcommand{\Q}[3]{\Needspace{14\baselineskip}\refstepcounter{qn}\subsection*{Q\theqn.\ #1}%
  \addcontentsline{toc}{subsection}{Q\theqn.\ #1}%
  \noindent\textit{#2}\hfill\textit{[#3 marks]}\par}

\newenvironment{solution}
  {\par\noindent\textbf{Solution.}\par\nobreak}
  {\par\nobreak\noindent\hfill\(\square\)\par}

\newcommand{\step}[1]{\par\addvspace{0.5em}\noindent\textbf{#1}\par\nobreak}
\newcommand{\pd}[2]{\frac{\partial #1}{\partial #2}}
\newcommand{\pdd}[2]{\frac{\partial^2 #1}{\partial #2^2}}
\newcommand{\hb}{\bar h}
\newcommand{\lam}{\lambda_{n+1}}

\pagestyle{fancy}
\fancyhf{}
\fancyhead[L]{MTL712}
\fancyhead[R]{Practice: Notes 4, 5 and 6}
\fancyfoot[C]{\thepage}
\renewcommand{\headrulewidth}{0.4pt}
\renewcommand{\footrulewidth}{0pt}

\title{\textbf{MTL712: Computational Methods for Differential\\Equations}\\[6pt]
\large Practice Questions with Complete Solutions\\
on Notes 4, 5 and 6}
\author{\normalsize Multistep methods, BVPs, parabolic, elliptic and
hyperbolic PDEs}
\date{}

\begin{document}

\maketitle
\thispagestyle{plain}

These questions cover every topic of Notes 4, 5 and 6, in the order in
which the topics appear in the notes.  They are written in the style
of the major examinations, including the questions of recent papers:
proofs, derivations, convergence and stability analyses, numerical
calculations and short conceptual answers.

Each question has three parts.
\begin{itemize}
  \item The \emph{question}, in a box, with its type and suggested
    marks.
  \item A \textbf{Remember} block, marked by a rule on the left: the
    definitions, theorems and formulas needed for the question.  Read
    these blocks on their own for a quick revision.
  \item A complete \emph{solution}, written as a full-marks answer.
\end{itemize}
Try each question before reading its solution.  All numerical answers
have been checked by computer.

\begin{center}
\renewcommand{\arraystretch}{1.25}
\begin{tabular}{@{}lll@{}}
\toprule
Section & Topic & Notes\\
\midrule
1 & Linear multistep methods for \(y'=f\) & Note 4\\
2 & Methods for \(y''=f\) & Note 4\\
3 & Boundary value problems & Note 4\\
4 & Classification and difference operators & Note 4\\
5 & Parabolic equations: schemes, LTE, convergence, stability & Note 4\\
6 & Lecture 35: ADI, three-level and Du Fort--Frankel schemes & Note 5\\
7 & Block matrices and Du Fort--Frankel by the matrix method & Note 6\\
8 & Elliptic equations and the five-point formula & Note 6\\
9 & ADI iterations: Dirichlet, Neumann and Robin problems & Note 6\\
10 & Hyperbolic equations: the wave equation & Note 6\\
11 & The heat equation: comparison of methods & Note 6\\
\bottomrule
\end{tabular}
\end{center}

{\linespread{1.0}\selectfont\setlength{\parskip}{0pt}\tableofcontents}
\clearpage

%#####################################################################
\Needspace{32\baselineskip}
\section{Linear Multistep Methods for \(y'=f(x,y)\)}
%#####################################################################

Notation as in Note 4: \(\rho(\zeta)=\sum\alpha_j\zeta^j\),
\(\sigma(\zeta)=\sum\beta_j\zeta^j\), \(w=\zeta-1\), and the error
constant of an order-\(p\) method is \(C_{p+1}\).

%---------------------------------------------------------------------
\Q{Order via \(\rho(\zeta)/\log\zeta-\sigma(\zeta)\)}{Proof}{5}\label{q:orderlemma}
\begin{qbox}
Let \(f\) be sufficiently smooth and \(p\) a positive integer.  Prove
that the linear \(k\)-step method
\(\sum_{j=0}^k\alpha_jy_{n+j}=h\sum_{j=0}^k\beta_jf_{n+j}\) is of order
\(p\) if and only if
\[
  \varphi(\zeta)=\frac{\rho(\zeta)}{\log\zeta}-\sigma(\zeta)
\]
has a zero of multiplicity \(p\) at \(\zeta=1\).  Show also that then
\[
  \varphi(\zeta)=C_{p+1}(\zeta-1)^p+O\bigl((\zeta-1)^{p+1}\bigr).
\]
\end{qbox}

\begin{rem}
\begin{itemize}
  \item \(L[y(x);h]=\sum_j\bigl[\alpha_jy(x+jh)-h\beta_jy'(x+jh)\bigr]
    =\sum_{q\ge0}C_qh^qy^{(q)}(x)\), with
    \(C_0=\sum\alpha_j\), \(C_1=\sum j\alpha_j-\sum\beta_j\),
    \(C_q=\frac1{q!}\sum j^q\alpha_j-\frac1{(q-1)!}\sum j^{q-1}\beta_j\).
  \item Order \(p\): \(C_0=\dots=C_p=0\), \(C_{p+1}\ne0\) (error
    constant).
  \item Key trick: apply \(L\) to \(y=e^x\) at \(x=0\), then put
    \(\zeta=e^h\).
\end{itemize}
\end{rem}

\begin{solution}
\step{Step 1: apply \(L\) to \(e^x\)}
Every derivative of \(e^x\) equals \(1\) at \(x=0\).  Hence
\[
  g(h):=L[e^x;h]\big|_{x=0}
  =\sum_j\bigl(\alpha_je^{jh}-h\beta_je^{jh}\bigr)
  =\rho(e^h)-h\,\sigma(e^h)
  =\sum_{q\ge0}C_qh^q .
\]

\step{Step 2: change of variable}
Put \(\zeta=e^h\), so \(h=\log\zeta\) and \(\zeta\to1\iff h\to0\)
(\(\log\) is single valued after cutting the plane along the negative
real axis \((-\infty,0]\)).  Then
\[
  g(\log\zeta)=\rho(\zeta)-\log\zeta\,\sigma(\zeta)
  =\log\zeta\cdot\varphi(\zeta).
\]

\step{Step 3 (\(\Leftarrow\))}
Suppose \(\varphi(\zeta)=c(\zeta-1)^p+O\bigl((\zeta-1)^{p+1}\bigr)\)
with \(c\ne0\).  Since \(\zeta-1=e^h-1=h+O(h^2)\),
\((\zeta-1)^p=h^p+O(h^{p+1})\), so
\[
  g(h)=h\,\varphi(e^h)=h\bigl[c\,h^p+O(h^{p+1})\bigr]
  =c\,h^{p+1}+O(h^{p+2}).
\]
Comparing with \(g(h)=\sum C_qh^q\): \(C_0=\dots=C_p=0\) and
\(C_{p+1}=c\ne0\).  So the method has order \(p\), with
\(C_{p+1}=c\).

\step{Step 4 (\(\Rightarrow\))}
If the order is \(p\), then \(g(h)=C_{p+1}h^{p+1}+O(h^{p+2})\) and
\[
  \varphi(\zeta)=\frac{g(\log\zeta)}{\log\zeta}
  =C_{p+1}(\log\zeta)^p+O\bigl((\log\zeta)^{p+1}\bigr)
  =C_{p+1}(\zeta-1)^p+O\bigl((\zeta-1)^{p+1}\bigr),
\]
using \(\log\zeta=(\zeta-1)+O\bigl((\zeta-1)^2\bigr)\).  Since
\(C_{p+1}\ne0\), the zero at \(\zeta=1\) has multiplicity exactly
\(p\).
\end{solution}

%---------------------------------------------------------------------
\Q{A one-parameter family of two-step methods}{Analysis}{8}\label{q:family}
\begin{qbox}
Consider the family
\[
  y_{n+2}-(1+a)y_{n+1}+a\,y_n
  =\frac h2\bigl[(3-a)f_{n+1}-(1+a)f_n\bigr],
\]
where \(a\) is a real parameter.
\begin{enumerate}[label=(\roman*)]
  \item For which \(a\) is the method consistent?
  \item What is the maximum order, and the error constant in that case?
  \item For which \(a\) is it zero-stable?
  \item For which \(a\) is it convergent?  Is the method of maximal
    order convergent?
\end{enumerate}
\end{qbox}

\begin{rem}
\begin{itemize}
  \item Consistent \(\iff\) \(\rho(1)=0\) and \(\rho'(1)=\sigma(1)\)
    \(\iff C_0=C_1=0\).
  \item Zero-stable (root condition): all roots of \(\rho\) satisfy
    \(|\zeta|\le1\), and those with \(|\zeta|=1\) are simple.
  \item Dahlquist: convergent \(\iff\) consistent and zero-stable.
  \item First Dahlquist barrier: zero-stable \(k\)-step order
    \(\le k+1\) (\(k\) odd), \(\le k+2\) (\(k\) even).
\end{itemize}
\end{rem}

\begin{solution}
Here \(\alpha_0=a\), \(\alpha_1=-(1+a)\), \(\alpha_2=1\);
\(\beta_0=-\frac{1+a}2\), \(\beta_1=\frac{3-a}2\), \(\beta_2=0\).

\step{(i) Consistency}
\(C_0=a-(1+a)+1=0\) and
\(C_1=\sum j\alpha_j-\sum\beta_j=\bigl(-(1+a)+2\bigr)-\bigl(\frac{3-a}2-\frac{1+a}2\bigr)
=(1-a)-(1-a)=0\).  So the method is consistent for \emph{every} \(a\).

\step{(ii) Order}
\begin{align*}
  C_2&=\tfrac12\bigl(-(1+a)+4\bigr)-\tfrac{3-a}2=\tfrac{3-a}2-\tfrac{3-a}2=0,\\
  C_3&=\tfrac16\bigl(-(1+a)+8\bigr)-\tfrac12\cdot\tfrac{3-a}2
      =\tfrac{7-a}6-\tfrac{3-a}4=\tfrac{a+5}{12},\\
  C_4&=\tfrac1{24}\bigl(-(1+a)+16\bigr)-\tfrac16\cdot\tfrac{3-a}2
      =\tfrac{15-a}{24}-\tfrac{3-a}{12}=\tfrac{a+9}{24}.
\end{align*}
For \(a\ne-5\) the order is \(2\) with \(C_3=\frac{a+5}{12}\).  For
\(a=-5\), \(C_3=0\) and \(C_4=\frac4{24}=\frac16\ne0\): the maximum
order is \(3\), with error constant \(C_4=\frac16\).  (For \(a=-5\) the
method is \(y_{n+2}+4y_{n+1}-5y_n=h(4f_{n+1}+2f_n)\).)

\step{(iii) Zero-stability}
\(\rho(\zeta)=\zeta^2-(1+a)\zeta+a=(\zeta-1)(\zeta-a)\), with roots
\(1\) and \(a\).  The root condition holds if and only if \(|a|\le1\)
and \(a\ne1\) (for \(a=1\), \(\zeta=1\) is a double root on the unit
circle).  So the method is zero-stable for \(-1\le a<1\).

\step{(iv) Convergence}
By Dahlquist's theorem, convergent \(\iff\) consistent and
zero-stable, i.e.\ for \(-1\le a<1\).  The method of maximal order
(\(a=-5\)) has the root \(\zeta=-5\), so it is \emph{not} zero-stable
and not convergent.  This agrees with the Dahlquist barrier: an
explicit zero-stable two-step method has order at most \(2\).
\end{solution}

%---------------------------------------------------------------------
\Q{Two methods from \(\rho(\zeta)=\zeta^2-1\)}{Numerical / Henrici's theorem}{10}\label{q:old1}
\begin{qbox}
Let \(\rho(\zeta)=\zeta^2-1\).
\begin{enumerate}[label=(\alph*)]
  \item Use Henrici's construction to find the \emph{explicit} method of
    maximal order.  State its order and error constant.
  \item Find the \emph{implicit} method of maximal order, its order and
    error constant.
  \item Show that the method in (a) has no interval of absolute
    stability.
\end{enumerate}
\end{qbox}
\begin{rem}
\begin{itemize}
  \item Henrici: given \(\rho\) (\(\rho(1)=0\), \(\rho'(1)\ne0\)),
    expand \(\rho(\zeta)/\log\zeta\) in powers of \(w=\zeta-1\) and
    truncate at \(w^{k'}\) to get \(\sigma\).  Explicit: \(k'=k-1\);
    implicit: \(k'=k\).
  \item \(\dfrac{w}{\log(1+w)}=1+\dfrac w2-\dfrac{w^2}{12}+\dfrac{w^3}{24}-\dfrac{19w^4}{720}+\cdots\).
  \item The first omitted coefficient is \(C_{p+1}\).
  \item Stability polynomial \(\pi(z,\hb)=\rho(z)-\hb\sigma(z)\);
    absolutely stable at \(\hb\) if all its roots satisfy \(|z|<1\).
\end{itemize}
\end{rem}
\begin{solution}
\step{Expansion}
With \(w=\zeta-1\), \(\rho=w(2+w)\), and using
\(\frac{w}{\log(1+w)}=1+\frac w2-\frac{w^2}{12}+\frac{w^3}{24}-\frac{19w^4}{720}+\cdots\),
\[
  \frac{\rho(\zeta)}{\log\zeta}
  =(2+w)\Bigl[1+\frac w2-\frac{w^2}{12}+\frac{w^3}{24}-\frac{19w^4}{720}+\cdots\Bigr]
  =2+2w+\frac13w^2+0\cdot w^3-\frac1{90}w^4+\cdots .
\]

\step{(a) Explicit (\(k'=k-1=1\))}
Keep terms up to \(w\): \(\sigma=2+2w=2\zeta\).  The method is the
\emph{mid-point rule}
\[
  y_{n+2}-y_n=2hf_{n+1}.
\]
The first omitted term is \(\frac13w^2\), so the order is \(2\) and the
error constant is \(C_3=\frac13\).

\step{(b) Implicit (\(k'=k=2\))}
Keep terms up to \(w^2\):
\(\sigma=2+2w+\frac13w^2=\frac13(\zeta^2+4\zeta+1)\).  This is
\emph{Milne--Simpson},
\[
  y_{n+2}-y_n=\frac h3\bigl(f_{n+2}+4f_{n+1}+f_n\bigr).
\]
The \(w^3\) coefficient vanishes, so the order is \(4\), with
\(C_5=-\frac1{90}\).

\step{(c) Absolute stability of the mid-point rule}
The stability polynomial is \(\pi(z,\hb)=z^2-2\hb z-1\), with roots
\(z_{1,2}=\hb\pm\sqrt{\hb^2+1}\).  Their product is \(-1\), so
\(|z_1||z_2|=1\), and both cannot be less than \(1\) in modulus.  Indeed
for \(\hb>0\), \(z_1=\hb+\sqrt{\hb^2+1}>1\); for \(\hb<0\),
\(z_2=\hb-\sqrt{\hb^2+1}<-1\); and for \(\hb=0\) the roots are \(\pm1\).
Hence the method is absolutely stable for no real \(\hb\): the interval
of absolute stability is empty.  (It is zero-stable, since the roots
\(\pm1\) of \(\rho\) are simple, but only weakly stable.)
\end{solution}

%---------------------------------------------------------------------
\Q{Henrici's theorem}{Proof}{8}\label{q:old7}
\begin{qbox}
Let \(\rho\) have degree \(k\) with \(\rho(1)=0\), \(\rho'(1)\ne0\), and
let \(0\le k'\le k\).  Prove that there is a unique \(\sigma\) of degree
\(\le k'\) such that the method \((\rho,\sigma)\) has order
\(\ge k'+1\).  (You may use: the order is \(p\) iff
\(\varphi=\rho/\log\zeta-\sigma\) has a zero of order \(p\) at
\(\zeta=1\).)
\end{qbox}
\begin{rem}
\begin{itemize}
  \item Order lemma (Q\ref{q:orderlemma}): order \(p\) \(\iff\)
    \(\varphi=\rho/\log\zeta-\sigma\) has a zero of order \(p\) at
    \(\zeta=1\).
  \item \(\rho\) and \(\log\zeta\) both have simple zeros at \(1\), so
    \(\rho/\log\zeta\) is holomorphic there, with value \(\rho'(1)\).
  \item Taylor coefficients are unique.
\end{itemize}
\end{rem}
\begin{solution}
\(\rho\) and \(\log\zeta\) both have a simple zero at \(\zeta=1\), so
\(\rho/\log\zeta\) is holomorphic near \(1\):
\[
  \frac{\rho(\zeta)}{\log\zeta}=\sum_{m\ge0}c_m(\zeta-1)^m .
\]
From \(\rho=\log\zeta\sum c_m(\zeta-1)^m\), differentiating and setting
\(\zeta=1\) gives \(c_0=\rho'(1)\ne0\).  Define
\(\sigma(\zeta)=\sum_{m=0}^{k'}c_m(\zeta-1)^m\).  Then
\[
  \varphi(\zeta)=\sum_{m>k'}c_m(\zeta-1)^m
\]
has a zero of order at least \(k'+1\) at \(\zeta=1\), so the order is at
least \(k'+1\).  Also \(\sigma(1)=c_0=\rho'(1)\), so the method is
consistent.  Conversely, any \(\sigma\) of degree \(\le k'\) giving
order \(\ge k'+1\) must agree with the Taylor series of
\(\rho/\log\zeta\) up to \((\zeta-1)^{k'}\); since Taylor coefficients
are unique, \(\sigma\) is unique.
\end{solution}

%---------------------------------------------------------------------
\Q{Simpson's three-eighths method}{Numerical}{8}\label{q:old2}
\begin{qbox}
For the three-step method
\[
  y_{n+3}-y_n=\frac{3h}8\bigl(f_{n+3}+3f_{n+2}+3f_{n+1}+f_n\bigr),
\]
\begin{enumerate}[label=(\alph*)]
  \item show that it is consistent and zero-stable;
  \item find its order and error constant;
  \item explain why its order is the largest possible for a zero-stable
    three-step method.
\end{enumerate}
\end{qbox}
\begin{rem}
\begin{itemize}
  \item \(C_q=\frac1{q!}\sum j^q\alpha_j-\frac1{(q-1)!}\sum j^{q-1}\beta_j\),
    \(q\ge2\).
  \item Weakly stable: zero-stable with roots other than \(1\) on the
    unit circle.
  \item First Dahlquist barrier: order of a zero-stable \(k\)-step
    method \(\le k+1\) (\(k\) odd), \(\le k+2\) (\(k\) even), and
    \(\le k\) if it is explicit.
\end{itemize}
\end{rem}
\begin{solution}
Here \(\alpha=(-1,0,0,1)\), \(\beta=\frac38(1,3,3,1)\),
\(\rho(\zeta)=\zeta^3-1\), \(\sigma(\zeta)=\frac38(\zeta^3+3\zeta^2+3\zeta+1)\).

\step{(a)}
\(\rho(1)=0\) and \(\rho'(1)=3=\sigma(1)=\frac38\cdot8\), so the method
is consistent.  The roots of \(\rho\) are the cube roots of unity
\(1,e^{\pm2\pi i/3}\): all have modulus \(1\) and are simple, so the
root condition holds and the method is zero-stable (weakly stable,
since roots other than \(1\) lie on the unit circle).

\step{(b)}
Use \(C_q=\frac1{q!}\sum j^q\alpha_j-\frac1{(q-1)!}\sum j^{q-1}\beta_j\).
First \(C_0=\rho(1)=0\) and \(C_1=\rho'(1)-\sigma(1)=3-3=0\).  For
\(q\ge2\), computing each sum directly,
\[
  \sum_jj^q\alpha_j=3^q,
  \qquad
  \sum_jj^{q-1}\beta_j=\tfrac38\bigl(3^{q-1}+3\cdot2^{q-1}+3\bigr)\quad(q\ge2),
\]
and
\begin{align*}
  C_2&=\tfrac92-\tfrac38(3+6+3)=\tfrac92-\tfrac92=0,\\
  C_3&=\tfrac{27}6-\tfrac38\cdot\tfrac{9+12+3}{2}=\tfrac92-\tfrac92=0,\\
  C_4&=\tfrac{81}{24}-\tfrac38\cdot\tfrac{27+24+3}{6}=\tfrac{27}8-\tfrac{27}8=0,\\
  C_5&=\tfrac{243}{120}-\tfrac38\cdot\tfrac{81+48+3}{24}
      =\tfrac{81}{40}-\tfrac{33}{16}=-\tfrac3{80}\ne0 .
\end{align*}
So the order is \(4\) and the error constant is \(C_5=-\frac3{80}\).

\step{(c)}
By the first Dahlquist barrier, a zero-stable \(k\)-step method has
order at most \(k+1\) when \(k\) is odd.  For \(k=3\) the maximum is
\(4\), which this method attains.
\end{solution}

%---------------------------------------------------------------------
\Q{BDF2 from \(\sigma(\zeta)=\zeta^2\)}{Derivation}{8}\label{q:old3}
\begin{qbox}
Find the polynomial \(\rho\) of degree \(2\) such that the method with
\(\sigma(\zeta)=\zeta^2\) has maximal order.  Write the method, find its
order and error constant, and check zero-stability.  Also give the
method normalised so that \(\alpha_2=1\).
\end{qbox}
\begin{rem}
\begin{itemize}
  \item Given \(\sigma\), choose \(\rho\): expand
    \(\log\zeta\cdot\sigma(\zeta)\) in powers of \(w=\zeta-1\) and keep
    terms up to \(w^k\).
  \item \(\log(1+w)=w-\frac{w^2}2+\frac{w^3}3-\cdots\).
  \item Strongly stable: \(\zeta=1\) is the only root of \(\rho\) on
    the unit circle.
\end{itemize}
\end{rem}
\begin{solution}
\step{Construction}
For given \(\sigma\), expand \(\log\zeta\cdot\sigma(\zeta)\) in powers of
\(w\) and truncate at degree \(k=2\):
\[
  \log(1+w)\,(1+w)^2
  =\Bigl(w-\frac{w^2}2+\frac{w^3}3-\cdots\Bigr)(1+2w+w^2)
  =w+\frac32w^2+\frac13w^3+\cdots .
\]
So
\[
  \rho=w+\frac32w^2=(\zeta-1)+\frac32(\zeta-1)^2
  =\frac32\zeta^2-2\zeta+\frac12,
\]
and the method is
\[
  \frac32y_{n+2}-2y_{n+1}+\frac12y_n=hf_{n+2}
  \qquad\text{(BDF2)}.
\]

\step{Order and error constant}
\[
  \varphi(\zeta)=\frac{\rho-\log\zeta\cdot\sigma}{\log\zeta}
  =\frac{-\frac13w^3+\cdots}{w+\cdots}=-\frac13w^2+\cdots,
\]
so by the order lemma the order is \(2\) and \(C_3=-\frac13\).
Normalised (divide by \(\frac32\)):
\[
  y_{n+2}-\frac43y_{n+1}+\frac13y_n=\frac23hf_{n+2},
  \qquad C_3=-\frac29 .
\]

\step{Zero-stability}
\(\rho(\zeta)=\frac12(3\zeta-1)(\zeta-1)\) has roots \(1\) and
\(\frac13\): the root condition holds, and \(1\) is the only root on the
unit circle, so the method is strongly stable.  Consistency:
\(\rho(1)=0\), \(\rho'(1)=3-2=1=\sigma(1)\).
\end{solution}

%---------------------------------------------------------------------
\Q{Intervals of absolute stability}{Stability analysis}{8}\label{q:absint}
\begin{qbox}
\begin{enumerate}[label=(\alph*)]
  \item Find the interval of absolute stability of the two-step
    Adams--Bashforth method
    \(y_{n+2}-y_{n+1}=\frac h2\bigl(3f_{n+1}-f_n\bigr)\).
  \item Find the region of absolute stability of the trapezoidal rule
    \(y_{n+1}-y_n=\frac h2(f_{n+1}+f_n)\) for \(y'=\lambda y\).  Is it
    A-stable?
\end{enumerate}
\end{qbox}

\begin{rem}
\begin{itemize}
  \item Test equation \(y'=\lambda y\), \(\hb=\lambda h\);
    stability polynomial \(\pi(z,\hb)=\rho(z)-\hb\,\sigma(z)\).
  \item Absolutely stable at \(\hb\) \(\iff\) every root of
    \(\pi(\cdot,\hb)\) has \(|z|<1\).
  \item A real quadratic \(z^2+a_1z+a_0\) has both roots in \(|z|<1\)
    if and only if \(|a_0|<1\) and \(|a_1|<1+a_0\) (Schur--Cohn /
    Jury test).
  \item A-stable: the region contains the whole left half-plane
    \(\operatorname{Re}\hb<0\).
\end{itemize}
\end{rem}

\begin{solution}
\step{(a) Adams--Bashforth 2}
\(\rho(z)=z^2-z\), \(\sigma(z)=\frac12(3z-1)\), so
\[
  \pi(z,\hb)=z^2-\Bigl(1+\frac{3\hb}2\Bigr)z+\frac{\hb}2 .
\]
Here \(a_1=-\bigl(1+\frac32\hb\bigr)\), \(a_0=\frac12\hb\).  The
conditions are:
\begin{itemize}
  \item \(|a_0|<1\): \(-2<\hb<2\).
  \item \(|a_1|<1+a_0\), i.e.\
    \(-\bigl(1+\frac{\hb}2\bigr)<1+\frac{3\hb}2<1+\frac{\hb}2\).
    The right inequality gives \(\hb<0\); the left one gives
    \(-2<2\hb\), i.e.\ \(\hb>-1\).
\end{itemize}
Hence the interval of absolute stability is
\[
  \boxed{-1<\hb<0}.
\]
(Check at the end points: \(\hb=0\) gives roots \(0,1\);
\(\hb=-1\) gives \(z^2+\frac12z-\frac12=(z+1)(z-\frac12)\), a root
\(-1\).)

\step{(b) Trapezoidal rule}
Applied to \(y'=\lambda y\):
\(\bigl(1-\frac{\hb}2\bigr)y_{n+1}=\bigl(1+\frac{\hb}2\bigr)y_n\), so
the single root is
\[
  z=\frac{1+\hb/2}{1-\hb/2}=\frac{2+\hb}{2-\hb}.
\]
\(|z|<1\iff|2+\hb|<|2-\hb|\iff\) \(\hb\) is closer to \(-2\) than to
\(2\) \(\iff\operatorname{Re}\hb<0\).  So the region of absolute
stability is the open left half-plane, and the interval is
\((-\infty,0)\).  The method is A-stable.  In particular, for
\(y'=-y\) (\(\lambda=-1\)) it is absolutely stable for every \(h>0\).
\end{solution}

%---------------------------------------------------------------------
\Q{Absolute instability for small positive \(\hb\)}{Proof}{10}\label{q:old5}
\begin{qbox}
Prove that every consistent zero-stable LMM is absolutely unstable for
all sufficiently small \(\hb>0\).  Explain why this is not a defect.
\end{qbox}
\begin{rem}
\begin{itemize}
  \item \(\pi(e^{\hb},\hb)=\sum_qC_q\hb^{\,q}=O(\hb^{\,p+1})\) for a
    method of order \(p\).
  \item \(\pi(z,\hb)=(\alpha_k-\hb\beta_k)\prod_s(z-r_s(\hb))\), and
    \(r_s(\hb)\to\zeta_s\) (roots of \(\rho\)) as \(\hb\to0\).
  \item Zero-stability makes \(\zeta=1\) a simple root.
\end{itemize}
\end{rem}
\begin{solution}
Let the order be \(p\ge1\).  Expanding the exponentials,
\[
  \pi(e^{\hb},\hb)=\rho(e^{\hb})-\hb\,\sigma(e^{\hb})
  =\sum_{q\ge0}C_q\hb^{\,q}=O(\hb^{\,p+1}).
\]
Factorise \(\pi(z,\hb)=(\alpha_k-\hb\beta_k)\prod_{s=1}^k(z-r_s(\hb))\).
As \(\hb\to0\), \(r_s(\hb)\to\zeta_s\), the roots of \(\rho\).
Consistency gives \(\rho(1)=0\) and zero-stability makes this root
simple, say \(\zeta_1=1\), \(\zeta_s\ne1\) for \(s\ge2\).  So in
\[
  \pi(e^{\hb},\hb)=(\alpha_k-\hb\beta_k)\prod_s\bigl(e^{\hb}-r_s(\hb)\bigr)
\]
every factor except \(e^{\hb}-r_1(\hb)\) tends to a non-zero limit.
Hence \(|e^{\hb}-r_1(\hb)|\le K\hb^{\,p+1}\), and
\[
  |r_1(\hb)|\ge e^{\hb}-K\hb^{\,p+1}\ge1+\hb-K\hb^{\,p+1}>1
\]
once \(K\hb^{\,p}<1\).  So some root has modulus greater than \(1\): the
method is absolutely unstable for small \(\hb>0\).

This is not a defect: for \(\lambda>0\) the exact solution
\(e^{\lambda x}\) grows, and \(r_1\approx e^{\hb}\) is the root that
correctly reproduces this growth.
\end{solution}

%---------------------------------------------------------------------
\Q{A convergent method is consistent}{Proof}{8}\label{q:old4}
\begin{qbox}
Prove that a convergent linear multistep method satisfies
\(\rho(1)=0\) and \(\rho'(1)=\sigma(1)\).  You may assume that a
convergent method is zero-stable.
\end{qbox}
\begin{rem}
\begin{itemize}
  \item Convergence: \(y_n\to y(x)\) as \(h\to0\), \(x_n=x\) fixed,
    for all consistent starting values.
  \item Test problems: \(y'=0,\ y(0)=1\) gives \(\rho(1)=0\);
    \(y'=1,\ y(0)=0\) gives \(\rho'(1)=\sigma(1)\).
  \item A convergent method is zero-stable (so \(\zeta=1\) is a simple
    root of \(\rho\)).
\end{itemize}
\end{rem}
\begin{solution}
\step{\(\rho(1)=0\)}
Apply the method to \(y'=0\), \(y(0)=1\) (solution \(y\equiv1\)), with
exact starting values \(y_0=\dots=y_{k-1}=1\).  The difference equation
\(\sum_j\alpha_jy_{n+j}=0\) does not involve \(h\), so \(y_n\) does not
depend on \(h\), and convergence (\(y_n\to1\) as \(h\to0\) with
\(nh=x\) fixed, i.e.\ \(n\to\infty\)) means \(\lim_{n\to\infty}y_n=1\).
Letting \(n\to\infty\) in the difference equation gives
\(\sum_j\alpha_j=\rho(1)=0\).

\step{\(\rho'(1)=\sigma(1)\)}
Apply the method to \(y'=1\), \(y(0)=0\) (solution \(y=x\)):
\[
  \sum_j\alpha_jy_{n+j}=h\,\sigma(1).
\]
Since \(\rho(1)=0\) and the method is zero-stable, \(1\) is a simple root
of \(\rho\), so \(\rho'(1)\ne0\).  Put \(K=\sigma(1)/\rho'(1)\) and
\(y_n=Knh\).  Then
\[
  \sum_j\alpha_jK(n+j)h=Kh\bigl(n\rho(1)+\rho'(1)\bigr)=Kh\rho'(1)=h\sigma(1),
\]
so \(y_n\) solves the difference equation, with starting values
\(Ksh\to0\).  Convergence forces \(y_n=Knh=Kx\to x\) for every \(x>0\),
hence \(K=1\), i.e.\ \(\rho'(1)=\sigma(1)\).
\end{solution}

%---------------------------------------------------------------------
\Q{Short answers on multistep methods}{Short answer}{8}\label{q:old6}
\begin{qbox}
Answer briefly, with reasons.
\begin{enumerate}[label=(\alph*)]
  \item Can a zero-stable three-step method have order \(5\)?
  \item Is the method
    \(y_{n+2}-3y_{n+1}+2y_n=h(\beta_2f_{n+2}+\beta_1f_{n+1}+\beta_0f_n)\)
    convergent for some choice of the \(\beta_j\)?
  \item A \(k\)-step method has \(2k+2\) coefficients.  Why can it reach
    order \(2k\) at most, and what does zero-stability cost?
  \item State Dahlquist's equivalence theorem.
\end{enumerate}
\end{qbox}
\begin{rem}
\begin{itemize}
  \item Dahlquist equivalence: convergent \(\iff\) consistent \(+\)
    zero-stable.
  \item Barrier: zero-stable \(k\)-step order \(\le k+1\) (odd \(k\)),
    \(\le k+2\) (even \(k\)).
  \item Without zero-stability, \(2k\) is the maximum order (\(2k+1\)
    free coefficients).
\end{itemize}
\end{rem}
\begin{solution}
\step{(a)}
No.  By the first Dahlquist barrier the order of a zero-stable
\(k\)-step method is at most \(k+1\) for odd \(k\); for \(k=3\) that is
\(4\).

\step{(b)}
No.  \(\rho(\zeta)=(\zeta-1)(\zeta-2)\) has the root \(2\) outside the
unit disc, so the method is not zero-stable for any \(\sigma\), and
zero-stability is necessary for convergence.

\step{(c)}
After normalising (\(\alpha_k=1\)) there are \(2k+1\) free coefficients,
so we can impose \(C_0=\dots=C_{2k}=0\), giving order \(2k\).
Zero-stability reduces the attainable order to \(k+1\) (odd \(k\)) or
\(k+2\) (even \(k\)): roughly half.

\step{(d)}
An LMM applied to \(y'=f(x,y)\), with \(f\) Lipschitz in \(y\) and
consistent starting values, is convergent if and only if it is
consistent and zero-stable.
\end{solution}

%---------------------------------------------------------------------
\Q{Two steps of Adams--Bashforth}{Numerical}{6}\label{q:ab2num}
\begin{qbox}
Solve \(y'=x+y\), \(y(0)=1\), by the two-step Adams--Bashforth method
\(y_{n+2}=y_{n+1}+\frac h2\bigl(3f_{n+1}-f_n\bigr)\) with \(h=0.1\),
taking \(y_1\) from the exact solution \(y=2e^x-x-1\).  Find \(y_2\)
and \(y_3\) and their errors.  Is \(\hb=\lambda h\) inside the interval
of absolute stability?
\end{qbox}

\begin{rem}
\begin{itemize}
  \item A \(k\)-step method needs \(k\) starting values; \(y_1\) must
    come from a one-step method (or the exact solution) of adequate
    order.
  \item AB2 is explicit, order \(2\), error constant
    \(C_3=\frac5{12}\); interval of absolute stability \((-1,0)\)
    (Q\ref{q:absint}).
\end{itemize}
\end{rem}

\begin{solution}
\(f(x,y)=x+y\), \(y_0=1\), \(y_1=2e^{0.1}-1.1=1.110342\).
Then \(f_0=0+1=1\) and \(f_1=0.1+1.110342=1.210342\).

\step{\(y_2\)}
\[
  y_2=y_1+\frac{0.1}2\bigl(3f_1-f_0\bigr)
  =1.110342+0.05\,(3.631026-1)=1.110342+0.131551=1.241893 .
\]
\step{\(y_3\)}
\(f_2=0.2+1.241893=1.441893\), so
\begin{align*}
  y_3&=y_2+0.05\,(3f_2-f_1)=1.241893+0.05\,(4.325679-1.210342)\\
  &=1.241893+0.155767=1.397660 .
\end{align*}
\step{Errors}
\begin{center}
\begin{tabular}{@{}cccc@{}}
\toprule
\(x\) & AB2 & exact \(2e^x-x-1\) & error\\
\midrule
0.2 & 1.241893 & 1.242806 & \(9.1\times10^{-4}\)\\
0.3 & 1.397660 & 1.399718 & \(2.1\times10^{-3}\)\\
\bottomrule
\end{tabular}
\end{center}
The local error per step is about \(C_3h^3y'''=\frac5{12}(0.001)(2e^x)
\approx9\times10^{-4}\), as observed for the first step; errors then
accumulate.

\step{Absolute stability}
Here \(\partial f/\partial y=\lambda=1>0\), so \(\hb=0.1\) is
\emph{outside} \((-1,0)\): the computed solution grows.  This is
correct behaviour, since the exact solution grows like \(e^x\)
(Q\ref{q:old5}).
\end{solution}

%#####################################################################
\Needspace{32\baselineskip}
\section{Methods for \(y''=f(x,y)\)}
%#####################################################################

%---------------------------------------------------------------------
\Q{Consistency and zero-stability for \(y''=f\)}{Short answer}{5}\label{q:old10}
\begin{qbox}
Verify that Störmer's method satisfies the consistency conditions for
\(y''=f\), and explain why the root condition for \(y''=f\) allows a
double root of \(\rho\) on the unit circle.
\end{qbox}
\begin{rem}
\begin{itemize}
  \item Method \(\sum\alpha_jy_{n+j}=h^2\sum\beta_jf_{n+j}\);
    \(C_0=\sum\alpha_j\), \(C_1=\sum j\alpha_j\),
    \(C_q=\frac1{q!}\sum j^q\alpha_j-\frac1{(q-2)!}\sum j^{q-2}\beta_j\).
  \item Consistency: \(\rho(1)=\rho'(1)=0\), \(\rho''(1)=2\sigma(1)\).
  \item Zero-stability for \(y''=f\): roots of \(\rho\) in
    \(|\zeta|\le1\), roots on \(|\zeta|=1\) of multiplicity \(\le2\).
  \item Convergent \(\iff\) consistent and zero-stable.
\end{itemize}
\end{rem}
\begin{solution}
Störmer: \(\rho(\zeta)=(\zeta-1)^2\), \(\sigma(\zeta)=\zeta\).  Then
\(\rho(1)=0\), \(\rho'(1)=0\), \(\rho''(1)=2=2\sigma(1)\): the
consistency conditions \(C_0=C_1=C_2=0\) hold.

Consistency for \(y''=f\) \emph{forces} \(\zeta=1\) to be a double root
of \(\rho\) (since \(\rho(1)=\rho'(1)=0\)).  This mirrors the exact
problem: \(y''=0\) has the solutions \(1\) and \(x\), and the difference
equation \(\rho(E)y_n=0\) reproduces them as \(1\) and \(n=(x-a)/h\).
Because the starting values must also approximate the initial slope
(condition (ii) in the definition of convergence), the coefficient of
the component \(n\) is \(O(h)\), so this component stays bounded: it is
just the linear part \(x\,y'(a)\) of the solution.  A double root of
modulus one is therefore harmless, whereas multiplicity \(3\) would add
components like \(n^2\), which are not controlled in this way.
\end{solution}

%---------------------------------------------------------------------
\Q{Störmer and Numerov from \(\rho(\zeta)=(\zeta-1)^2\)}{Derivation}{6}\label{q:numhenrici}
\begin{qbox}
For methods \(\rho(E)y_n=h^2\sigma(E)f_n\) for \(y''=f(x,y)\), the order
is \(p\) if and only if
\(\varphi(\zeta)=\rho(\zeta)/(\log\zeta)^2-\sigma(\zeta)\) has a zero of
order \(p\) at \(\zeta=1\).  Taking \(\rho(\zeta)=(\zeta-1)^2\), use this
to construct the explicit and the implicit methods of maximal order.
Give their orders and error constants.
\end{qbox}

\begin{rem}
\begin{itemize}
  \item Henrici-type theorem for \(y''=f\): if \(\rho(1)=\rho'(1)=0\),
    \(\rho''(1)\ne0\), then for each \(0\le k'\le k\) there is a unique
    \(\sigma\) of degree \(\le k'\) giving order \(\ge k'+1\): truncate
    the expansion of \(\rho/(\log\zeta)^2\) at \(w^{k'}\).
  \item For an order-\(p\) method,
    \(\varphi(\zeta)=C_{p+2}(\zeta-1)^p+\cdots\).
  \item \(\dfrac{w^2}{\log^2(1+w)}=1+w+\dfrac{w^2}{12}+0\cdot w^3-\dfrac{w^4}{240}+\cdots\).
\end{itemize}
\end{rem}

\begin{solution}
With \(w=\zeta-1\), \(\rho=w^2\), and
\[
  \frac{\rho(\zeta)}{(\log\zeta)^2}=\frac{w^2}{\log^2(1+w)}
  =1+w+\frac{w^2}{12}+0\cdot w^3-\frac{w^4}{240}+\cdots .
\]
(From \(\log(1+w)=w\bigl(1-\frac w2+\frac{w^2}3-\cdots\bigr)\), so
\(\frac{w^2}{\log^2(1+w)}=\bigl(1-\frac w2+\frac{w^2}3-\frac{w^3}4+\cdots\bigr)^{-2}\),
expanded term by term.)

\step{Explicit (\(k'=k-1=1\))}
\(\sigma=1+w=\zeta\).  The method is \emph{Störmer's method}
\[
  y_{n+2}-2y_{n+1}+y_n=h^2f_{n+1}.
\]
The first omitted term is \(\frac1{12}w^2\), so \(p=2\) and
\(C_4=\frac1{12}\); since \(\sigma(1)=1\), the normalised error
constant is also \(\frac1{12}\).

\step{Implicit (\(k'=k=2\))}
\(\sigma=1+w+\frac{w^2}{12}=\frac1{12}\bigl(\zeta^2+10\zeta+1\bigr)\).
The method is \emph{Numerov's method}
\[
  y_{n+2}-2y_{n+1}+y_n=\frac{h^2}{12}\bigl(f_{n+2}+10f_{n+1}+f_n\bigr).
\]
The \(w^3\) coefficient vanishes, so \(\varphi=-\frac1{240}w^4+\cdots\):
the order is \(4\) and \(C_6=-\frac1{240}\) (again \(\sigma(1)=1\)).

Both have \(\rho=(\zeta-1)^2\), a double root on the unit circle, which
is allowed for \(y''=f\): they are zero-stable, consistent, and hence
convergent.
\end{solution}

%---------------------------------------------------------------------
\Q{A one-parameter family containing Störmer and Numerov}{Derivation}{8}\label{q:old8}
\begin{qbox}
Consider
\[
  y_{n+1}-2y_n+y_{n-1}=h^2\bigl[\beta f_{n+1}+(1-2\beta)f_n+\beta f_{n-1}\bigr].
\]
\begin{enumerate}[label=(\alph*)]
  \item Show that it has order at least \(2\) for every \(\beta\), and
    find the error constant.
  \item For which \(\beta\) is the order higher?  Find that order and
    the error constant.
\end{enumerate}
\end{qbox}
\begin{rem}
\begin{itemize}
  \item Symmetric methods (\(\alpha_j=\alpha_{k-j}\),
    \(\beta_j=\beta_{k-j}\)) have all odd \(C_q=0\) when expanded about
    the centre.
  \item Order \(p\) for \(y''=f\): \(C_0=\dots=C_{p+1}=0\),
    \(C_{p+2}\ne0\); error constant \(C_{p+2}/\sigma(1)\).
  \item Störmer: \(\beta=0\), order 2; Numerov: \(\beta=\frac1{12}\),
    order 4.
\end{itemize}
\end{rem}
\begin{solution}
Expand about \(x_n\).  With
\(y(x_{n\pm1})=\sum\frac{(\pm h)^q}{q!}y^{(q)}\) and the right-hand side
weights \(\beta,1-2\beta,\beta\) at \(x_{n+1},x_n,x_{n-1}\), the method
is symmetric, so all odd \(C_q\) vanish.  Using the two-step
coefficients \(C_0,\dots,C_6\) derived for Störmer and Numerov in
Note 4, with \(a_1=-2\), \(a_2=1\), \(b_0=b_2=\beta\),
\(b_1=1-2\beta\):
\begin{align*}
  C_0&=1-2+1=0,\qquad C_1=1-1=0,\\
  C_2&=\tfrac12(1+1)-(\beta+1-2\beta+\beta)=1-1=0,\\
  C_4&=\tfrac1{24}\cdot2-\tfrac12\cdot2\beta=\tfrac1{12}-\beta,\\
  C_6&=\tfrac1{720}\cdot2-\tfrac1{24}\cdot2\beta=\tfrac1{360}-\tfrac\beta{12}.
\end{align*}

\step{(a)}
\(C_0=\dots=C_3=0\) for every \(\beta\), so the order is at least \(2\).
If \(\beta\ne\frac1{12}\), the order is exactly \(2\) with
\(C_4=\frac1{12}-\beta\).  (\(\beta=0\) gives Störmer with
\(C_4=\frac1{12}\).)

\step{(b)}
For \(\beta=\frac1{12}\), \(C_4=C_5=0\) and
\(C_6=\frac1{360}-\frac1{144}=-\frac1{240}\ne0\), so the order is \(4\).
This is Numerov's method,
\[
  y_{n+1}-2y_n+y_{n-1}=\frac{h^2}{12}\bigl(f_{n+1}+10f_n+f_{n-1}\bigr),
  \qquad C_6=-\frac1{240}.
\]
Since \(\sigma(1)=1\), the normalised error constants are the same.
\end{solution}

%---------------------------------------------------------------------
\Q{Numerov's method for \(y''=y\)}{Numerical}{8}\label{q:old9}
\begin{qbox}
Use Numerov's method with \(h=0.1\) to approximate \(y(0.2)\) and
\(y(0.3)\) for \(y''=y\), \(y(0)=1\), \(y'(0)=0\).  Obtain \(y_1\) from
a fourth-order Taylor series.  Compare with the exact solution
\(\cosh x\).
\end{qbox}
\begin{rem}
\begin{itemize}
  \item Numerov:
    \(y_{n+1}-2y_n+y_{n-1}=\frac{h^2}{12}(f_{n+1}+10f_n+f_{n-1})\),
    order 4, error \(-\frac1{240}h^6y^{(6)}\) per step.
  \item For linear \(f=g(x)y\) the implicit equation is solved
    directly:
    \(\bigl(1-\frac{h^2}{12}g_{n+1}\bigr)y_{n+1}
    =\bigl(2+\frac{10h^2}{12}g_n\bigr)y_n
    -\bigl(1-\frac{h^2}{12}g_{n-1}\bigr)y_{n-1}\).
  \item Starting value from a Taylor series of matching order.
\end{itemize}
\end{rem}
\begin{solution}
\step{Starting value}
\(y''=y\) gives \(y''(0)=1\), \(y'''(0)=y'(0)=0\), \(y^{(4)}(0)=1\), so
\[
  y_1=1+\frac{h^2}2+\frac{h^4}{24}=1.0050041667 .
\]

\step{Recurrence}
With \(f=y\), Numerov is linear in \(y_{n+1}\):
\[
  \Bigl(1-\frac{h^2}{12}\Bigr)y_{n+1}
  =\Bigl(2+\frac{10h^2}{12}\Bigr)y_n-\Bigl(1-\frac{h^2}{12}\Bigr)y_{n-1}.
\]
For \(h=0.1\): \(1-\frac{h^2}{12}=0.9991666667\) and
\(2+\frac{10h^2}{12}=2.0166666667\).

\step{Steps}
\begin{align*}
  y_2&=\frac{2.0166666667\times1.0050041667-0.9991666667\times1}{0.9991666667}
     =1.0200667570,\\
  y_3&=\frac{2.0166666667\times1.0200667570-0.9991666667\times1.0050041667}{0.9991666667}
     =1.0453385226 .
\end{align*}

\step{Comparison}
\begin{center}
\begin{tabular}{@{}cccc@{}}
\toprule
\(x\) & Numerov & \(\cosh x\) & error\\
\midrule
0.2 & 1.0200667570 & 1.0200667556 & \(1.4\times10^{-9}\)\\
0.3 & 1.0453385226 & 1.0453385141 & \(8.5\times10^{-9}\)\\
\bottomrule
\end{tabular}
\end{center}
The tiny errors reflect the order \(4\) of the method (local error
\(-\frac1{240}h^6y^{(6)}\approx-4\times10^{-9}\) per step).
\end{solution}

%---------------------------------------------------------------------
\Q{A tabulated solution and Numerov's method}{Numerical}{6}\label{q:numtable}
\begin{qbox}
A table of a function \(y=f(x)\) is given.
\begin{center}
\begin{tabular}{@{}cccccc@{}}
\toprule
\(x\) & 1 & 2 & 3 & 4 & 5\\
\midrule
\(y\) & 2.00000 & 8.25000 & 27.11111 & 64.06250 & 125.04000\\
\bottomrule
\end{tabular}
\end{center}
It is known that \(f\) satisfies
\(\dfrac{d^2y}{dx^2}=\dfrac{n(n+1)}{x^2}\,y\), where \(n\) is a positive
integer.
\begin{enumerate}[label=(\alph*)]
  \item Find \(n\), using Numerov's formula centred at \(x=4\).
  \item Compute \(f(6)\) by Numerov's method with step size \(1\).
\end{enumerate}
\end{qbox}

\begin{rem}
\begin{itemize}
  \item Numerov with \(f=g(x)y\), \(g=c/x^2\), \(c=n(n+1)\):
    \(y_{n+1}-2y_n+y_{n-1}=\frac{h^2}{12}\bigl(g_{n+1}y_{n+1}+10g_ny_n+g_{n-1}y_{n-1}\bigr)\).
  \item The formula is linear in \(c\): one centre gives one equation
    for \(c\).  It is only approximate (error \(O(h^6y^{(6)})\)), so
    round to the nearest value of the form \(n(n+1)\).
  \item To march: solve the same formula for \(y_{n+1}\).
\end{itemize}
\end{rem}

\begin{solution}
Let \(c=n(n+1)\) and \(h=1\).

\step{(a) Finding \(n\)}
Centred at \(x=4\) (points \(3,4,5\)):
\[
  y_5-2y_4+y_3=\frac c{12}\Bigl(\frac{y_5}{25}+10\,\frac{y_4}{16}+\frac{y_3}{9}\Bigr).
\]
Left side: \(125.04-128.125+27.11111=24.02611\).  Bracket:
\(5.00160+40.03906+3.01235=48.05301\).  Hence
\[
  c=\frac{12\times24.02611}{48.05301}=5.99990\approx6 .
\]
So \(n(n+1)=6\), and since \(n>0\), \(\boxed{n=2}\).  (Centring at
\(x=2\) gives \(c\approx5.90\); the error of Numerov's formula is
larger there because the solution varies faster relative to \(h=1\).
Either way the nearest value \(n(n+1)\) is \(6\).)

\step{(b) Computing \(f(6)\)}
With \(c=6\), Numerov at \(x=5\) solved for \(y_6\):
\[
  \Bigl(1-\frac6{12\cdot36}\Bigr)y_6
  =\Bigl(2+\frac{10\cdot6}{12\cdot25}\Bigr)y_5
  -\Bigl(1-\frac6{12\cdot16}\Bigr)y_4,
\]
i.e.
\[
  0.986111\,y_6=2.2\times125.04-0.96875\times64.0625
  =275.08800-62.06055=213.02745,
\]
so
\[
  \boxed{f(6)\approx216.0278}.
\]
(The data are \(y=x^3+x^{-2}\), the solution for \(n=2\); the exact
value is \(216.02778\), so the error is about \(6\times10^{-5}\).)
\end{solution}

%---------------------------------------------------------------------
\Q{Störmer and Numerov for a non-linear equation}{Numerical}{8}\label{q:numnonlin}
\begin{qbox}
For \(y''=6y^2\), \(y(0)=1\), \(y'(0)=-2\) (exact solution
\(y=1/(1+x)^2\)), with \(h=0.1\):
\begin{enumerate}[label=(\alph*)]
  \item find \(y_1\) from the Taylor series up to \(h^4\);
  \item find \(y_2\) by Störmer's method;
  \item find \(y_2\) by Numerov's method, solving the non-linear
    equation by Newton's method.
\end{enumerate}
Compare with the exact values.
\end{qbox}

\begin{rem}
\begin{itemize}
  \item Störmer \(y_{n+1}=2y_n-y_{n-1}+h^2f_n\) is explicit: no
    iteration.
  \item Numerov is implicit; for non-linear \(f\), solve
    \(G(y_{n+1})=y_{n+1}-\frac{h^2}{12}f(x_{n+1},y_{n+1})-R_n=0\) by
    Newton:
    \(y^{(s+1)}=y^{(s)}-G(y^{(s)})/G'(y^{(s)})\), starting from the
    Störmer-like predictor \(2y_n-y_{n-1}\).
\end{itemize}
\end{rem}

\begin{solution}
\step{(a) Starting value}
From \(y''=6y^2\): \(y''(0)=6\), \(y'''=12yy'\Rightarrow y'''(0)=-24\),
\(y^{(4)}=12(y'^2+yy'')\Rightarrow y^{(4)}(0)=12(4+6)=120\).  So
\[
  y_1=1-0.2+\frac{6(0.01)}2-\frac{24(0.001)}6+\frac{120(0.0001)}{24}
  =0.8265,
\]
against the exact \(1/1.21=0.826446\).

\step{(b) Störmer}
\[
  y_2=2y_1-y_0+h^2\cdot6y_1^2=0.653+0.01\times6\times0.683102=0.693986 .
\]

\step{(c) Numerov}
\(y_2-2y_1+y_0=\frac{h^2}{12}\bigl(6y_2^2+60y_1^2+6y_0^2\bigr)\), with
\(60y_1^2=40.986135\), gives
\[
  G(z)=z-0.653-\frac{0.01}{12}\bigl(6z^2+46.986135\bigr)
      =-0.005z^2+z-0.692155=0,
  \qquad G'(z)=1-0.01z .
\]
Newton from \(z_0=2y_1-y_0=0.653\):
\[
  G(z_0)=-0.041287,\quad G'(z_0)=0.99347,\quad
  z_1=0.653+\frac{0.041287}{0.99347}=0.694558,
\]
and \(z_2=0.694567\), unchanged to six decimals afterwards.  (The
quadratic can also be solved directly:
\(z=\bigl(1-\sqrt{1-0.02\times0.692155}\bigr)/0.01=0.694567\).)

\step{Comparison}
\begin{center}
\begin{tabular}{@{}lcc@{}}
\toprule
 & \(y_2\) & error\\
\midrule
Störmer & 0.693986 & \(4.6\times10^{-4}\)\\
Numerov & 0.694567 & \(1.2\times10^{-4}\)\\
exact \(1/1.44\) & 0.694444 & --\\
\bottomrule
\end{tabular}
\end{center}
Numerov (order 4) is more accurate than Störmer (order 2); part of
its remaining error comes from the starting value \(y_1\).
\end{solution}

%#####################################################################
\Needspace{32\baselineskip}
\section{Boundary Value Problems}
%#####################################################################

%---------------------------------------------------------------------
\Q{A mixed boundary condition and tridiagonal form}{Derivation}{6}\label{q:old12}
\begin{qbox}
For \(y''=f(x)\) on \([0,1]\) with \(y'(0)-y(0)=1\), write the boundary
equation at \(x_0=0\) using the second-order one-sided formula for
\(y'(0)\).  Then eliminate \(y_2\) using the equation at \(x_1\) so that
the system becomes tridiagonal.
\end{qbox}
\begin{rem}
\begin{itemize}
  \item Second-order formulas:
    \(y''_i\approx\frac{y_{i-1}-2y_i+y_{i+1}}{h^2}\),
    \(y'_i\approx\frac{y_{i+1}-y_{i-1}}{2h}\),
    \(y'_0\approx\frac{-3y_0+4y_1-y_2}{2h}\),
    \(y'_N\approx\frac{3y_N-4y_{N-1}+y_{N-2}}{2h}\), all \(O(h^2)\).
  \item Row \(0\) then has an extra entry \(R_0y_2\).  If \(B_1\ne0\),
    replace row \(0\) by row\(_0-\frac{R_0}{B_1}\)row\(_1\) to make the
    matrix tridiagonal.
\end{itemize}
\end{rem}
\begin{solution}
With \(y'(x_0)\approx\frac{-3y_0+4y_1-y_2}{2h}\), the boundary row is
\[
  -y_0+\frac{-3y_0+4y_1-y_2}{2h}=1 .
\]
It contains \(R_0y_2\) with \(R_0=-\frac1{2h}\).  The row at \(x_1\) is
\(\frac{y_0-2y_1+y_2}{h^2}=f_1\), with \(B_1=\frac1{h^2}\).  Form
row\(_0-\frac{R_0}{B_1}\)row\(_1\)\(=\)row\(_0+\frac h2\)row\(_1\):
\[
  -y_0+\frac{-3y_0+4y_1-y_2}{2h}+\frac{y_0-2y_1+y_2}{2h}=1+\frac h2f_1 .
\]
The \(y_2\) terms cancel, leaving
\[
  \Bigl(-1-\frac1h\Bigr)y_0+\frac1h\,y_1=1+\frac h2f_1,
\]
which has non-zero entries only in columns \(0\) and \(1\).
\end{solution}

%---------------------------------------------------------------------
\Q{The Thomas algorithm and its stability}{Derivation / proof}{8}\label{q:thomas}
\begin{qbox}
For the tridiagonal system \(A_iy_{i-1}+C_iy_i+B_iy_{i+1}=F_i\),
\(i=0,\dots,N\) (\(A_0=B_N=0\)):
\begin{enumerate}[label=(\alph*)]
  \item derive the recurrences
    \(y_i=\alpha_{i+1}y_{i+1}+\beta_{i+1}\) of the Thomas algorithm;
  \item show that if the matrix is strictly diagonally dominant,
    \(|C_i|>|A_i|+|B_i|\), then no denominator vanishes and
    \(|\alpha_i|<1\) for all \(i\);
  \item explain why errors do not grow in the back substitution;
  \item for the Crank--Nicolson scheme for \(u_t=u_{xx}\) with
    Dirichlet conditions (Mitchell--Griffiths, Ch.~2, Example 3),
    verify that diagonal dominance holds for every \(r>0\).
\end{enumerate}
\end{qbox}

\begin{rem}
\begin{itemize}
  \item \(\alpha_1=-\frac{B_0}{C_0}\), \(\beta_1=\frac{F_0}{C_0}\);
    \(\alpha_{i+1}=\dfrac{-B_i}{C_i+\alpha_iA_i}\),
    \(\beta_{i+1}=\dfrac{F_i-\beta_iA_i}{C_i+\alpha_iA_i}\);
    \(y_N=\beta_{N+1}\), then \(y_i=\alpha_{i+1}y_{i+1}+\beta_{i+1}\)
    backwards.
  \item Equivalent to \(M=LU\) with \(L\) unit lower bidiagonal and
    \(U\) upper bidiagonal.  Cost \(O(N)\).
\end{itemize}
\end{rem}

\begin{solution}
\step{(a) Derivation}
Row \(0\): \(C_0y_0+B_0y_1=F_0\), so \(y_0=\alpha_1y_1+\beta_1\) with
\(\alpha_1=-B_0/C_0\), \(\beta_1=F_0/C_0\).  Suppose
\(y_{i-1}=\alpha_iy_i+\beta_i\).  Substituting into row \(i\),
\[
  A_i(\alpha_iy_i+\beta_i)+C_iy_i+B_iy_{i+1}=F_i
  \ \Longrightarrow\
  y_i=\frac{-B_i}{C_i+\alpha_iA_i}\,y_{i+1}+\frac{F_i-\beta_iA_i}{C_i+\alpha_iA_i}
  =\alpha_{i+1}y_{i+1}+\beta_{i+1}.
\]
In the last row \(B_N=0\), so \(\alpha_{N+1}=0\) and
\(y_N=\beta_{N+1}\).  The forward sweep computes all
\(\alpha_i,\beta_i\); back substitution gives
\(y_{N-1},\dots,y_0\).  Each step costs a fixed number of operations:
\(O(N)\) in total.

\step{(b) Diagonal dominance}
Induction.  \(|\alpha_1|=|B_0|/|C_0|<1\) since \(|C_0|>|B_0|\).  If
\(|\alpha_i|<1\), then
\[
  |C_i+\alpha_iA_i|\ge|C_i|-|\alpha_i||A_i|\ge|C_i|-|A_i|>|B_i|\ge0,
\]
so the denominator is non-zero, and
\(|\alpha_{i+1}|=|B_i|/|C_i+\alpha_iA_i|<1\).  Hence the algorithm
never divides by zero and \(|\alpha_i|<1\) for all \(i\).

\step{(c) Error propagation}
If \(y_{i+1}\) carries an error \(e_{i+1}\), back substitution gives
\(e_i=\alpha_{i+1}e_{i+1}\), so \(|e_i|\le|e_{i+1}|\): errors are not
amplified.  For the finite difference scheme of
\(y''-q(x)y=f\) with \(q>0\), the rows
\(y_{i-1}-(2+h^2q_i)y_i+y_{i+1}=h^2f_i\) are strictly diagonally
dominant, so the Thomas algorithm is safe for every \(h\).

\step{(d) Crank--Nicolson}
Each time step solves
\(-\frac r2U_{m-1}^{n+1}+(1+r)U_m^{n+1}-\frac r2U_{m+1}^{n+1}=d_m\),
\(m=1,\dots,M-1\), with known right-hand side \(d_m\).  Here
\(|C_m|=1+r\) and \(|A_m|+|B_m|=\frac r2+\frac r2=r\) (one of them is
absent in the first and last rows).  Since \(1+r>r\) for every
\(r>0\), the system is strictly diagonally dominant for all \(r\): the
Thomas algorithm is always stable for Crank--Nicolson.
\end{solution}

%---------------------------------------------------------------------
\Q{A tridiagonal system solved as a difference equation}{Derivation / numerical}{4}\label{q:tridde}
\begin{qbox}
Solve \(Ax=b\), where
\[
  A=\begin{bmatrix}
    -\frac52&1&&\\
    1&-\frac52&1&\\
    &1&-\frac52&1\\
    &&1&-\frac52
  \end{bmatrix},
  \qquad
  b=\begin{bmatrix}-1\\0\\0\\0\end{bmatrix}.
\]
State the difference equation which replaces the matrix formulation,
and solve the system by solving the difference equation.
\end{qbox}

\begin{rem}
\begin{itemize}
  \item A constant tridiagonal system
    \(cx_{i-1}+ax_i+bx_{i+1}=0\) is a linear difference equation; the
    first and last rows become boundary conditions by introducing
    \(x_0\) and \(x_{N+1}\).
  \item Try \(x_i=m^i\): auxiliary equation \(bm^2+am+c=0\).  Distinct
    roots: \(x_i=P\,m_1^i+Q\,m_2^i\); fix \(P,Q\) from the two boundary
    conditions.
  \item The same idea proves the eigenvalue formula for tridiagonal
    matrices (Q\ref{q:tridproof}).
\end{itemize}
\end{rem}

\begin{solution}
\step{The difference equation}
Introduce \(x_0=1\) and \(x_5=0\).  Then all four rows read
\[
  x_{i-1}-\tfrac52x_i+x_{i+1}=0,\qquad i=1,2,3,4,
  \qquad x_0=1,\ x_5=0 .
\]
(Row \(1\): \(x_0-\frac52x_1+x_2=0\iff-\frac52x_1+x_2=-1\); row \(4\):
\(x_3-\frac52x_4+x_5=0\iff x_3-\frac52x_4=0\).)

\step{General solution}
With \(x_i=m^i\): \(m^2-\frac52m+1=0\), so \(m=2\) or \(m=\frac12\), and
\[
  x_i=P\,2^i+Q\,2^{-i}.
\]

\step{Boundary conditions}
\(x_0=P+Q=1\) and \(x_5=32P+\frac Q{32}=0\Rightarrow Q=-1024P\).  Then
\(P(1-1024)=1\), so
\[
  P=-\frac1{1023},\qquad Q=\frac{1024}{1023},
  \qquad
  \boxed{x_i=\frac{2^{10-i}-2^i}{1023}}.
\]

\step{The solution}
\[
  x_1=\frac{510}{1023}=0.498534,\qquad
  x_2=\frac{252}{1023}=0.246334,
\]
\[
  x_3=\frac{120}{1023}=0.117302,\qquad
  x_4=\frac{48}{1023}=0.046921 .
\]
Check, row 1: \(-\frac52(0.498534)+0.246334=-1.000001\approx-1\); row 4:
\(0.117302-\frac52(0.046921)=0.000000\).
\end{solution}

%---------------------------------------------------------------------
\Q{A Dirichlet problem by finite differences}{Numerical}{10}\label{q:old11}
\begin{qbox}
Solve \(y''-y=0\), \(y(0)=0\), \(y(1)=1\), with the central difference
scheme and \(h=\frac14\).  Solve the tridiagonal system by the Thomas
algorithm and compare with the exact solution
\(y=\sinh x/\sinh1\).
\end{qbox}
\begin{rem}
\begin{itemize}
  \item \(y''-q(x)y=f\): rows
    \(y_{i-1}-(2+h^2q_i)y_i+y_{i+1}=h^2f_i\); boundary values move to
    the right-hand side.
  \item Thomas algorithm: express each unknown through the next
    (forward sweep), then back-substitute.
  \item Error of the central scheme: \(O(h^2)\).
\end{itemize}
\end{rem}
\begin{solution}
\step{The system}
\(\frac{y_{i-1}-2y_i+y_{i+1}}{h^2}-y_i=0\) gives
\(y_{i-1}-d\,y_i+y_{i+1}=0\) with \(d=2+h^2=2.0625\), for \(i=1,2,3\),
with \(y_0=0\), \(y_4=1\):
\[
  \begin{pmatrix}-d&1&0\\1&-d&1\\0&1&-d\end{pmatrix}
  \begin{pmatrix}y_1\\y_2\\y_3\end{pmatrix}
  =\begin{pmatrix}0\\0\\-1\end{pmatrix}.
\]
The matrix is strictly diagonally dominant (\(2.0625>2\)), so the
solution is unique.

\step{Forward sweep (express each unknown through the next)}
\begin{align*}
  \text{row 1:}&\quad y_1=\frac1d\,y_2=0.484848\,y_2,\\
  \text{row 2:}&\quad y_2=\frac1{d-0.484848}\,y_3=0.633854\,y_3,\\
  \text{row 3:}&\quad y_2-d\,y_3=-1\ \Rightarrow\
    y_3=\frac1{d-0.633854}=0.699963 .
\end{align*}

\step{Back substitution}
\(y_2=0.633854\times0.699963=0.443674\) and
\(y_1=0.484848\times0.443674=0.215115\).

\step{Comparison}
\begin{center}
\begin{tabular}{@{}cccc@{}}
\toprule
\(x\) & FD & exact & error\\
\midrule
0.25 & 0.215115 & 0.214952 & \(1.6\times10^{-4}\)\\
0.50 & 0.443674 & 0.443409 & \(2.6\times10^{-4}\)\\
0.75 & 0.699963 & 0.699724 & \(2.4\times10^{-4}\)\\
\bottomrule
\end{tabular}
\end{center}
The errors are \(O(h^2)\), as expected from the \(O(h^2)\) central
difference.
\end{solution}

%---------------------------------------------------------------------
\Q{Diagonal dominance and uniqueness}{Proof}{6}\label{q:old14}
\begin{qbox}
Prove that a strictly diagonally dominant matrix is non-singular.
Deduce that the central difference scheme for \(y''-q(x)y=f(x)\) with
Dirichlet conditions and \(q>0\) has a unique solution for every \(h\).
\end{qbox}
\begin{rem}
\begin{itemize}
  \item Strictly diagonally dominant: \(|a_{ii}|>\sum_{j\ne i}|a_{ij}|\)
    for every row.
  \item Proof idea: if \(Ax=0\), look at the row where \(|x_i|\) is
    largest.
\end{itemize}
\end{rem}
\begin{solution}
Suppose \(Ax=0\), \(x\ne0\), and pick \(i\) with
\(|x_i|=\max_j|x_j|>0\).  Row \(i\) gives
\(a_{ii}x_i=-\sum_{j\ne i}a_{ij}x_j\), so
\[
  |a_{ii}||x_i|\le\sum_{j\ne i}|a_{ij}||x_j|\le\Bigl(\sum_{j\ne i}|a_{ij}|\Bigr)|x_i|<|a_{ii}||x_i|,
\]
a contradiction.  So \(A\) is non-singular.

The scheme gives rows \(y_{i-1}-(2+h^2q_i)y_i+y_{i+1}=h^2f_i\).  The
diagonal entry has modulus \(2+h^2q_i>2\), and the off-diagonal entries
sum to at most \(2\) in modulus.  Hence the matrix is strictly
diagonally dominant and the solution is unique, for every \(h>0\).
\end{solution}

%---------------------------------------------------------------------
\Q{Newton's method for a non-linear BVP}{Numerical}{6}\label{q:old13}
\begin{qbox}
For \(y''=y^2\), \(y(0)=y(1)=1\), with \(h=\frac12\), write the
discrete equation for the single unknown \(y_1\approx y(\frac12)\).
Perform one Newton step from \(y_1^{(0)}=1\) and compare with the exact
root of the discrete equation.
\end{qbox}
\begin{rem}
\begin{itemize}
  \item Non-linear BVP \(y''=f(x,y,y')\): the difference equations
    \(G(Y)=0\) are solved by Newton,
    \(Y^{(s+1)}=Y^{(s)}-J^{-1}G(Y^{(s)})\), with tridiagonal Jacobian
    \(J\).
  \item Newton converges quadratically near a simple root.
\end{itemize}
\end{rem}
\begin{solution}
The scheme \(y_0-2y_1+y_2=h^2y_1^2\) with \(y_0=y_2=1\) gives
\[
  G(y_1)=2-2y_1-\tfrac14y_1^2=0,
  \qquad G'(y_1)=-2-\tfrac12y_1 .
\]
From \(y_1^{(0)}=1\): \(G(1)=-0.25\), \(G'(1)=-2.5\), so
\[
  y_1^{(1)}=1-\frac{G(1)}{G'(1)}=1-0.1=0.9 .
\]
The exact positive root is \(y_1=-4+2\sqrt6=0.898979\); one step already
has error \(1.0\times10^{-3}\), and the next step would square it
(quadratic convergence).  For larger \(N\) the Jacobian
\(J=A-h^2\,\mathrm{diag}(2y_i)\) is tridiagonal and each step is a
Thomas solve.
\end{solution}

%---------------------------------------------------------------------
\Q{Reduction to homogeneous boundary conditions}{Short derivation}{5}\label{q:old15}
\begin{qbox}
For collocation applied to \(y''+p(x)y'+q(x)y=f(x)\) on \([0,1]\) with
\(y(0)=1\), \(y'(1)=2\), find the linear function \(d(x)\) used to make
the boundary conditions homogeneous, and the new right-hand side.  Why
does the reduction fail for the pure Neumann conditions
\(y'(0)=c_0\), \(y'(1)=c_1\)?
\end{qbox}
\begin{rem}
\begin{itemize}
  \item BCs \(a_0y(a)+b_0y'(a)=c_0\), \(a_1y(b)+b_1y'(b)=c_1\) with
    \((a_0b_1-a_1b_0)+a_0a_1(b-a)\ne0\).
  \item Homogenise: \(\tilde y=y-d\), with \(d\) linear satisfying the
    BCs; then \(L\tilde y=f-L(d)\).
\end{itemize}
\end{rem}
\begin{solution}
Here \(a_0=1,b_0=0\) at \(x=0\) and \(a_1=0,b_1=1\) at \(x=1\).  With
\(d=\alpha x+\beta\): \(d(0)=\beta=1\) and \(d'(1)=\alpha=2\), so
\(d(x)=2x+1\).  (The determinant condition
\((a_0b_1-a_1b_0)+a_0a_1(b-a)=1\ne0\) holds.)  Then \(\tilde y=y-d\)
satisfies \(\tilde y(0)=0\), \(\tilde y'(1)=0\) and
\[
  L\tilde y=f-L(d)=f(x)-2p(x)-(2x+1)q(x).
\]
For pure Neumann conditions \(a_0=a_1=0\), so the determinant
\((a_0b_1-a_1b_0)+a_0a_1(b-a)=0\): a linear \(d\) cannot match two
slope conditions unless \(c_0=c_1\), and then not uniquely.
\end{solution}

%---------------------------------------------------------------------
\Q{Spectral collocation with a sine basis}{Numerical}{6}\label{q:speccol}
\begin{qbox}
Solve \(-y''+y=1\), \(y(0)=y(\pi)=0\), by collocation with
\(y_N(x)=\sum_{i=1}^N\xi_i\sin(ix)\) at the points
\(x_j=\frac{j\pi}{N+1}\), for \(N=2\) and \(N=3\).  Compare \(y_N(\pi/2)\)
with the exact value, given
\(y(x)=1-\dfrac{\cosh(x-\pi/2)}{\cosh(\pi/2)}\).
\end{qbox}

\begin{rem}
\begin{itemize}
  \item For \(-y''+q(x)y=f\) on \((0,\pi)\) with \(\psi_i=\sin(ix)\)
    (satisfying the BCs, \(-\psi_i''=i^2\psi_i\)):
    \[
      \sum_{i=1}^N\xi_i\bigl[i^2+q(x_j)\bigr]\sin(ix_j)=f(x_j),
      \qquad x_j=\frac{j\pi}{N+1},\ j=1,\dots,N .
    \]
  \item Symmetry: if \(q,f\) are symmetric about \(\pi/2\), the
    solution is too, and the coefficients of \(\sin(ix)\) with \(i\)
    even vanish.
\end{itemize}
\end{rem}

\begin{solution}
Here \(q=1\), \(f=1\), so the coefficient of \(\xi_i\) is
\((i^2+1)\sin(ix_j)\).

\step{\(N=2\): points \(\frac\pi3,\frac{2\pi}3\)}
\begin{align*}
  2\xi_1\sin\tfrac\pi3+5\xi_2\sin\tfrac{2\pi}3&=1,\\
  2\xi_1\sin\tfrac{2\pi}3+5\xi_2\sin\tfrac{4\pi}3&=1 .
\end{align*}
Since \(\sin\frac\pi3=\sin\frac{2\pi}3=\frac{\sqrt3}2\) and
\(\sin\frac{4\pi}3=-\frac{\sqrt3}2\), subtracting gives \(\xi_2=0\),
and then \(\sqrt3\,\xi_1=1\):
\[
  \xi_1=\frac1{\sqrt3}=0.577350,\qquad
  y_2\bigl(\tfrac\pi2\bigr)=\xi_1=0.577350 .
\]

\step{\(N=3\): points \(\frac\pi4,\frac\pi2,\frac{3\pi}4\)}
By symmetry \(\xi_2=0\).  The equations at \(\frac\pi4\) and
\(\frac\pi2\) are
\begin{align*}
  2\xi_1\tfrac{\sqrt2}2+10\xi_3\tfrac{\sqrt2}2&=1
    \iff 2\xi_1+10\xi_3=\sqrt2,\\
  2\xi_1\cdot1+10\xi_3\cdot(-1)&=1 .
\end{align*}
Adding and subtracting:
\[
  \xi_1=\frac{1+\sqrt2}4=0.603553,\qquad
  \xi_3=\frac{\sqrt2-1}{20}=0.020711,
\]
\[
  y_3\bigl(\tfrac\pi2\bigr)=\xi_1\sin\tfrac\pi2+\xi_3\sin\tfrac{3\pi}2
  =\xi_1-\xi_3=0.582843 .
\]

\step{Comparison}
Exact: \(y(\pi/2)=1-1/\cosh(\pi/2)=0.601463\).  The errors are
\(0.0241\) for \(N=2\) and \(0.0186\) for \(N=3\): adding basis
functions improves the approximation.  (The right-hand side \(f=1\) is
not zero at the end points, so the sine series converges slowly; for
smooth data vanishing at the ends the convergence would be much
faster.)
\end{solution}

%#####################################################################
\Needspace{32\baselineskip}
\section{Classification and Difference Operators}
%#####################################################################

%---------------------------------------------------------------------
\Q{Why \(B^2-AC\) decides the type}{Derivation}{6}\label{q:charderiv}
\begin{qbox}
For \(Au_{xx}+2Bu_{xy}+Cu_{yy}+Du_x+Eu_y+Fu=G\), suppose \(u\), \(u_x\),
\(u_y\) are known along a curve.  Show that the second derivatives are
not determined along the curve exactly when
\(A\,dy^2-2B\,dx\,dy+C\,dx^2=0\), and deduce the classification.
Apply it to \(t\,u_{xx}+2x\,u_{xt}+u_{tt}=0\).
\end{qbox}

\begin{rem}
\begin{itemize}
  \item Characteristics: curves along which \(u_{xx},u_{xy},u_{yy}\) are
    not determined by the data (and may be discontinuous).
  \item \(\frac{dy}{dx}=\frac{B\pm\sqrt{B^2-AC}}A\): hyperbolic
    (\(B^2-AC>0\), two real families), parabolic (\(=0\), one),
    elliptic (\(<0\), none).
  \item Only \(A,B,C\) decide the type; with variable coefficients
    the type can change from point to point.
\end{itemize}
\end{rem}

\begin{solution}
\step{Three equations for three unknowns}
Along the curve, \(u_x\) and \(u_y\) are known functions, so
\[
  d(u_x)=u_{xx}\,dx+u_{xy}\,dy,\qquad
  d(u_y)=u_{xy}\,dx+u_{yy}\,dy .
\]
Together with the PDE (all lower-order terms are known along the
curve), this is a linear system for \(u_{xx},u_{xy},u_{yy}\):
\[
  \begin{pmatrix}A&2B&C\\ dx&dy&0\\ 0&dx&dy\end{pmatrix}
  \begin{pmatrix}u_{xx}\\u_{xy}\\u_{yy}\end{pmatrix}
  =\begin{pmatrix}G-Du_x-Eu_y-Fu\\ d(u_x)\\ d(u_y)\end{pmatrix}.
\]
It fails to determine the second derivatives exactly when the
determinant vanishes:
\[
  A(dy^2-0)-2B(dx\,dy-0)+C(dx^2)=0
  \iff A\,dy^2-2B\,dx\,dy+C\,dx^2=0 .
\]

\step{Classification}
Dividing by \(dx^2\) (\(A\ne0\)):
\(A\bigl(\frac{dy}{dx}\bigr)^2-2B\frac{dy}{dx}+C=0\), so
\[
  \frac{dy}{dx}=\frac{B\pm\sqrt{B^2-AC}}A .
\]
Two real characteristic directions if \(B^2-AC>0\) (hyperbolic), one
if \(B^2-AC=0\) (parabolic), none if \(B^2-AC<0\) (elliptic).

\step{Example}
\(t\,u_{xx}+2x\,u_{xt}+u_{tt}=0\): \(A=t\), \(B=x\), \(C=1\), so
\(B^2-AC=x^2-t\).  The equation is hyperbolic for \(t<x^2\) (below the
parabola), parabolic on \(t=x^2\), and elliptic for \(t>x^2\).  In the
hyperbolic region the characteristics satisfy
\(\frac{dt}{dx}=\frac{x\pm\sqrt{x^2-t}}{t}\) (for \(t\ne0\)).
\end{solution}

%---------------------------------------------------------------------
\Q{Classify and find characteristics}{Short answer}{6}\label{q:old16}
\begin{qbox}
Classify each equation and, where real, find the characteristic
directions \(dy/dx\).
\begin{enumerate}[label=(\alph*)]
  \item \(u_{xx}+4u_{xy}+4u_{yy}=0\);
  \item \(u_{xx}-2u_{xy}-3u_{yy}+u_x=0\);
  \item \(x\,u_{xx}+u_{yy}=0\).
\end{enumerate}
\end{qbox}
\begin{rem}
\begin{itemize}
  \item Compare with \(Au_{xx}+2Bu_{xy}+Cu_{yy}+\cdots\): note the
    factor \(2\) on \(B\).
  \item \(B^2-AC>0\) hyperbolic, \(=0\) parabolic, \(<0\) elliptic;
    \(\frac{dy}{dx}=\frac{B\pm\sqrt{B^2-AC}}A\).
\end{itemize}
\end{rem}
\begin{solution}
Compare with \(Au_{xx}+2Bu_{xy}+Cu_{yy}+\cdots\) and use
\(\frac{dy}{dx}=\frac{B\pm\sqrt{B^2-AC}}A\).
\step{(a)}
\(A=1\), \(B=2\), \(C=4\): \(B^2-AC=0\), \emph{parabolic}; one
(double) characteristic direction \(dy/dx=2\).
\step{(b)}
\(A=1\), \(B=-1\), \(C=-3\): \(B^2-AC=1+3=4>0\), \emph{hyperbolic};
\(dy/dx=-1\pm2\), i.e.\ \(1\) and \(-3\).  (The term \(u_x\) does not
affect the type.)
\step{(c)}
\(A=x\), \(B=0\), \(C=1\): \(B^2-AC=-x\).  Elliptic for \(x>0\),
parabolic on \(x=0\), hyperbolic for \(x<0\), where
\(dy/dx=\pm\sqrt{-x}/x=\mp1/\sqrt{-x}\).
\end{solution}

%---------------------------------------------------------------------
\Q{The central second difference}{Derivation}{6}\label{q:old17}
\begin{qbox}
Using \(E=e^{hD}\) and \(\delta=E^{1/2}-E^{-1/2}\), show that
\[
  \delta^2=h^2D^2+\frac{h^4}{12}D^4+\frac{h^6}{360}D^6+\cdots,
\]
and deduce the leading truncation error of
\(\frac{u_{m+1}-2u_m+u_{m-1}}{h^2}\approx u_{xx}\).  Then invert the
series to get the first two terms of \(D^2\) in powers of \(\delta\).
\end{qbox}
\begin{rem}
\begin{itemize}
  \item \(E=e^{hD}\), \(\delta=E^{1/2}-E^{-1/2}=2\sinh\frac{hD}2\),
    so \(D=\frac2h\sinh^{-1}\frac\delta2\).
  \item \(D^2=\frac1{h^2}\bigl(\delta^2-\frac1{12}\delta^4+\frac1{90}\delta^6-\cdots\bigr)\).
  \item \(\frac{\delta^2u_m}{h^2}=u_{xx}+\frac{h^2}{12}u_{xxxx}+O(h^4)\).
\end{itemize}
\end{rem}
\begin{solution}
\(\delta=2\sinh\frac{hD}2\), so
\(\delta^2=4\sinh^2\frac{hD}2=2(\cosh hD-1)\).  With
\(\cosh t=1+\frac{t^2}2+\frac{t^4}{24}+\frac{t^6}{720}+\cdots\),
\[
  \delta^2=h^2D^2+\frac{h^4D^4}{12}+\frac{h^6D^6}{360}+\cdots .
\]
Hence
\(\frac{\delta^2u_m}{h^2}=u_{xx}+\frac{h^2}{12}u_{xxxx}+O(h^4)\): the
truncation error is \(\frac{h^2}{12}u_{xxxx}\).  Inverting,
\(h^2D^2=\delta^2-\frac{h^4D^4}{12}+\cdots=\delta^2-\frac{\delta^4}{12}+\cdots\)
(to this order \(h^4D^4=\delta^4\)), so
\(D^2=\frac1{h^2}\bigl(\delta^2-\frac1{12}\delta^4+\cdots\bigr)\), in
agreement with Note 4.
\end{solution}

%#####################################################################
\Needspace{32\baselineskip}
\section{Parabolic Equations: Schemes, Truncation Error, Convergence, Stability}
%#####################################################################

Throughout, \(r=k/h^2\) and \(s=\sin\frac{\beta h}2\).

%---------------------------------------------------------------------
\Q{Two steps of the explicit scheme}{Numerical}{8}\label{q:old18}
\begin{qbox}
Solve \(u_t=u_{xx}\), \(u(x,0)=\sin\pi x\), \(u(0,t)=u(1,t)=0\), by the
explicit scheme with \(h=\frac14\), \(r=\frac14\).  Find \(U\) at
\(x=\frac12\) after one and two time steps and compare with the exact
solution \(e^{-\pi^2t}\sin\pi x\).
\end{qbox}
\begin{rem}
\begin{itemize}
  \item Explicit (Schmidt) scheme:
    \(U_m^{n+1}=(1-2r)U_m^n+r\bigl(U_{m+1}^n+U_{m-1}^n\bigr)\), i.e.\
    \(U^{n+1}=(1+r\delta^2)U^n\); stable for \(r\le\frac12\).
  \item \(\delta^2\sin(\pi x_m)=-4\sin^2\frac{\pi h}2\,\sin(\pi x_m)\):
    \(\sin\pi x\) data stay proportional to \(\sin\pi x\).
  \item Exact solution for \(u(x,0)=\sin\pi x\):
    \(u=e^{-\pi^2t}\sin\pi x\).
\end{itemize}
\end{rem}
\begin{solution}
\(k=rh^2=\frac1{64}\).  Initially \(U^0=(0,\ \sin\frac\pi4,\ 1,\ \sin\frac{3\pi}4,\ 0)=(0,0.707107,1,0.707107,0)\).
\step{Step 1}
\(U_m^1=(1-2r)U_m^0+r(U_{m+1}^0+U_{m-1}^0)=\frac12U_m^0+\frac14(U_{m+1}^0+U_{m-1}^0)\):
\[
  U_2^1=\tfrac12(1)+\tfrac14(0.707107+0.707107)=0.853553,
  \quad
  U_1^1=\tfrac12(0.707107)+\tfrac14(1)=0.603553 .
\]
\step{Step 2}
\(U_2^2=\frac12(0.853553)+\frac14(2\times0.603553)=0.728553\).

\step{Why the numbers are so clean}
\(\sin\pi x_m\) is an eigenvector of the scheme: \(U_m^n=\xi^n\sin\pi x_m\)
with \(\xi=1-4r\sin^2\frac{\pi h}2=1-\sin^2\frac\pi8=\cos^2\frac\pi8=0.853553\).

\step{Comparison at \(x=\frac12\)}
\begin{center}
\begin{tabular}{@{}cccc@{}}
\toprule
\(t\) & explicit & exact & error\\
\midrule
\(1/64\) & 0.853553 & 0.857090 & \(3.5\times10^{-3}\)\\
\(2/64\) & 0.728553 & 0.734603 & \(6.0\times10^{-3}\)\\
\bottomrule
\end{tabular}
\end{center}
\end{solution}

%---------------------------------------------------------------------
\Q{Higher-order explicit schemes and the choice \(r=\frac16\)}{Derivation}{8}\label{q:r16}
\begin{qbox}
\begin{enumerate}[label=(\alph*)]
  \item Show that the local truncation error of the explicit scheme is
    \(T_m^n=\frac12k^2\bigl(u_{tt}-\frac1{6r}u_{xxxx}\bigr)+\cdots\), and
    deduce that \(r=\frac16\) gives a higher-order scheme.
  \item From
    \(u_m^{n+1}=\bigl[1+r\delta^2+\frac12r(r-\frac16)\delta^4+\cdots\bigr]u_m^n\),
    derive the five-point explicit scheme obtained by keeping the
    \(\delta^4\) term, and check that its coefficients add up to \(1\).
\end{enumerate}
\end{qbox}

\begin{rem}
\begin{itemize}
  \item \(u_m^{n+1}=\exp(kD^2)u_m^n\) with
    \(D^2=\frac1{h^2}\bigl(\delta^2-\frac1{12}\delta^4+\cdots\bigr)\).
  \item LTE: substitute the exact solution into the scheme;
    consistent if \(T/k\to0\).
  \item For the heat equation \(u_{tt}=u_{xxxx}\) (differentiate
    \(u_t=u_{xx}\)).
  \item \(\delta^4U_m=U_{m+2}-4U_{m+1}+6U_m-4U_{m-1}+U_{m-2}\).
\end{itemize}
\end{rem}

\begin{solution}
\step{(a) Truncation error}
\(T_m^n=u_m^{n+1}-(1-2r)u_m^n-r(u_{m+1}^n+u_{m-1}^n)\).  Expanding
about \((mh,nk)\),
\[
  u_m^{n+1}=u+ku_t+\tfrac12k^2u_{tt}+\cdots,\qquad
  r\bigl(u_{m+1}^n+u_{m-1}^n\bigr)=2ru+ku_{xx}+\tfrac1{12}kh^2u_{xxxx}+\cdots,
\]
using \(rh^2=k\).  The \(u\) terms cancel (\(1-(1-2r)-2r=0\)), so
\[
  T_m^n=k(u_t-u_{xx})+\tfrac12k^2u_{tt}-\tfrac1{12}kh^2u_{xxxx}+\cdots
  =\tfrac12k^2\Bigl(u_{tt}-\frac1{6r}u_{xxxx}\Bigr)+\cdots,
\]
since \(u_t=u_{xx}\) and \(kh^2=k^2/r\).  So \(T=O(k^2+kh^2)\) and
\(T/k=O(k+h^2)\): consistent.  Because \(u_{tt}=u_{xxxx}\), the
principal part is \(\frac12k^2\bigl(1-\frac1{6r}\bigr)u_{xxxx}\), which
vanishes when \(r=\frac16\).  The scheme
\(U_m^{n+1}=\frac23U_m^n+\frac16\bigl(U_{m+1}^n+U_{m-1}^n\bigr)\) then
has \(T=O(k^3)\) (equivalently \(O(k^3+kh^4)\)).  This agrees with
the series: the \(\delta^4\) coefficient \(\frac12r(r-\frac16)\) is zero
at \(r=\frac16\).

\step{(b) Five-point explicit scheme}
Put \(c=\frac12r(r-\frac16)\) and keep
\(U_m^{n+1}=(1+r\delta^2+c\,\delta^4)U_m^n\).  Expanding \(\delta^2\)
and \(\delta^4\) and collecting:
\begin{align*}
  U_m:&\quad 1-2r+6c=\tfrac12\bigl(2-5r+6r^2\bigr),\\
  U_{m\pm1}:&\quad r-4c=\tfrac23r(2-3r),\\
  U_{m\pm2}:&\quad c=-\tfrac1{12}r(1-6r).
\end{align*}
So
\[
  U_m^{n+1}=\tfrac12(2-5r+6r^2)U_m^n+\tfrac23r(2-3r)\bigl(U_{m+1}^n+U_{m-1}^n\bigr)
  -\tfrac1{12}r(1-6r)\bigl(U_{m+2}^n+U_{m-2}^n\bigr).
\]
Sum of coefficients: \((1-2r+6c)+2(r-4c)+2c=1\).  This must hold
because a constant is an exact solution of the heat equation.
\end{solution}

%---------------------------------------------------------------------
\Q{The fourth-order explicit formula: truncation error and stability}{Derivation / stability analysis}{8}\label{q:fourth}
\begin{qbox}
(Mitchell--Griffiths, Ch.~2, Exercises 1 and 6.)  For the explicit
formula
\[
  U_m^{n+1}=\tfrac12(2-5r+6r^2)U_m^n+\tfrac23r(2-3r)\bigl(U_{m+1}^n+U_{m-1}^n\bigr)
  -\tfrac1{12}r(1-6r)\bigl(U_{m+2}^n+U_{m-2}^n\bigr),
\]
\begin{enumerate}[label=(\alph*)]
  \item find the principal part of the truncation error;
  \item find the von Neumann stability condition.
\end{enumerate}
\end{qbox}

\begin{rem}
\begin{itemize}
  \item This formula is \((1+r\delta^2+c\,\delta^4)U^n\) with
    \(c=\frac12r(r-\frac16)\) (Q\ref{q:r16}).
  \item Exact:
    \(u^{n+1}=\bigl[1+r\delta^2+c\,\delta^4+d\,\delta^6+\cdots\bigr]u^n\),
    \(d=\frac16r\bigl(r^2-\frac12r+\frac1{15}\bigr)\).
  \item \(\delta^6u=h^6\partial_x^6u+O(h^8)\);
    \(\delta^2\to-4S\), \(\delta^4\to16S^2\) on a Fourier mode,
    \(S=\sin^2\frac{\beta h}2\in[0,1]\).
\end{itemize}
\end{rem}

\begin{solution}
\step{(a) Truncation error}
The exact solution satisfies
\(u_m^{n+1}=\exp(kD^2)u_m^n=\bigl[1+r\delta^2+c\,\delta^4+d\,\delta^6+\cdots\bigr]u_m^n\),
and the formula keeps only the terms up to \(\delta^4\).  So the
principal part of the truncation error is
\[
  T=d\,\delta^6u_m^n+\cdots
  =\frac r{180}\bigl(30r^2-15r+2\bigr)h^6\frac{\partial^6u}{\partial x^6}+\cdots
  =\frac{kh^4}{180}\bigl(30r^2-15r+2\bigr)\frac{\partial^6u}{\partial x^6}+\cdots .
\]
For fixed \(r\), \(T/k=O(h^4)\): the formula is fourth order in space
(hence its name).  The factor \(30r^2-15r+2\) has no real root
(discriminant \(225-240<0\)), so no choice of \(r\) removes this term.

\step{(b) Stability}
With \(U_m^n=\xi^ne^{i\beta mh}\) and \(S=\sin^2\frac{\beta h}2\),
\[
  \xi=1-4rS+16cS^2=1-4rS+8r\bigl(r-\tfrac16\bigr)S^2 .
\]
\emph{\(\xi\le1\):} \(-4rS+8r(r-\frac16)S^2\le0\iff2(r-\frac16)S\le1\);
the worst case \(S=1\) gives \(r\le\frac23\).

\emph{\(\xi\ge-1\):} put \(g(S)=\frac12(\xi+1)=1-2rS+4r(r-\frac16)S^2\).
If \(r\le\frac16\), \(g\) is concave, so its minimum is at an end point:
\(g(0)=1\) and \(g(1)=4r^2-\frac83r+1>0\) (negative discriminant).  If
\(r>\frac16\), the minimum of \(g\) over all \(S\) is
\(1-\frac{r}{4(r-1/6)}\), which is \(\ge0\) since \(4r-\frac23\ge r\) for
\(r\ge\frac29\); for \(\frac16<r<\frac29\) the vertex lies beyond
\(S=1\), and again \(g(1)>0\).  So \(\xi\ge-1\) always holds.

Hence the formula is stable if and only if
\[
  \boxed{0<r\le\tfrac23}.
\]
\end{solution}

%---------------------------------------------------------------------
\Q{Explicit scheme in self-adjoint form}{Derivation}{6}\label{q:selfadj}
\begin{qbox}
For \(u_t=\partial_x\bigl(a(x)u_x\bigr)\) with \(a>0\), derive the
explicit scheme
\[
  U_m^{n+1}=(1-2ra)U_m^n+r\bigl(a+\tfrac12ha'\bigr)U_{m+1}^n
  +r\bigl(a-\tfrac12ha'\bigr)U_{m-1}^n,
\]
with \(a,a'\) evaluated at \(x_m\).  Give the equivalent half-point
form, and a condition under which all coefficients are non-negative.
\end{qbox}

\begin{rem}
\begin{itemize}
  \item \(u_t=a'u_x+au_{xx}\), so \(L=a'D+aD^2\) and
    \(u^{n+1}=[1+kL+O(k^2)]u^n\).
  \item \(DU_m\approx\frac{U_{m+1}-U_{m-1}}{2h}\),
    \(D^2U_m\approx\frac{\delta^2U_m}{h^2}\).
  \item Half-point form:
    \(D(aD)U_m\approx\frac1{h^2}\bigl[a_{m+\frac12}(U_{m+1}-U_m)-a_{m-\frac12}(U_m-U_{m-1})\bigr]\).
  \item Non-negative coefficients summing to \(1\) give stability by
    the maximum principle.
\end{itemize}
\end{rem}

\begin{solution}
\step{Derivation}
\(u_t=\partial_x(au_x)=a'u_x+au_{xx}\).  Neglecting \(O(k^2)\),
\(u_m^{n+1}=\bigl[1+k(a'D+aD^2)\bigr]u_m^n\).  With central
differences,
\[
  U_m^{n+1}=U_m^n+k\Bigl[\frac{a'}{2h}\bigl(U_{m+1}^n-U_{m-1}^n\bigr)
  +\frac a{h^2}\bigl(U_{m+1}^n-2U_m^n+U_{m-1}^n\bigr)\Bigr].
\]
With \(k=rh^2\), the coefficient of \(U_{m+1}^n\) is
\(ra+\frac12rha'\), that of \(U_{m-1}^n\) is \(ra-\frac12rha'\), and
that of \(U_m^n\) is \(1-2ra\): the stated scheme.

\step{Half-point form}
Applying \(\delta\) twice with \(a\) at half points,
\[
  U_m^{n+1}=U_m^n+r\bigl[a_{m+\frac12}\bigl(U_{m+1}^n-U_m^n\bigr)
  -a_{m-\frac12}\bigl(U_m^n-U_{m-1}^n\bigr)\bigr].
\]
Since \(a_{m\pm\frac12}=a\pm\frac12ha'+O(h^2)\), the two forms agree
to this order.  For \(a\equiv1\) both reduce to the classical explicit
scheme.

\step{Non-negative coefficients}
All coefficients are non-negative if
\[
  ra\le\tfrac12\quad\text{and}\quad h|a'|\le2a,
\]
which hold for \(r\le\frac1{2\max a}\) and \(h\) small enough.  They add
up to \(1\), so then
\(\max_m|U_m^{n+1}|\le\max_m|U_m^n|\): the scheme is stable (maximum
principle).
\end{solution}

%---------------------------------------------------------------------
\Q{Crank--Nicolson: derivation and truncation error}{Derivation}{8}\label{q:cnlte}
\begin{qbox}
Derive the Crank--Nicolson scheme for \(u_t=u_{xx}\) from
\(u_m^{n+1}=\exp(kD^2)u_m^n\), and show that its local truncation error
is
\[
  T=-\frac{k^3}{12}u_{ttt}-\frac{kh^2}{12}u_{xxxx}+\cdots=O(k^3+kh^2).
\]
\end{qbox}

\begin{rem}
\begin{itemize}
  \item Split symmetrically about \(t_{n+1/2}\):
    \(\exp(-\frac12kD^2)u^{n+1}=\exp(\frac12kD^2)u^n\).
  \item Crank--Nicolson:
    \((1-\frac r2\delta^2)U_m^{n+1}=(1+\frac r2\delta^2)U_m^n\), i.e.\
    \((1+r)U_m^{n+1}-\frac r2(U_{m+1}^{n+1}+U_{m-1}^{n+1})
    =(1-r)U_m^n+\frac r2(U_{m+1}^n+U_{m-1}^n)\).
  \item Expand about the half level: odd powers of \(\frac k2\) cancel.
\end{itemize}
\end{rem}

\begin{solution}
\step{Derivation}
\(u^{n+1}=\exp(kD^2)u^n\iff\exp(-\frac12kD^2)u_m^{n+1}=\exp(\frac12kD^2)u_m^n\).
Replace \(D^2\) by \(\delta^2/h^2\) and keep two terms of each
exponential:
\[
  \bigl(1-\tfrac12r\delta^2\bigr)U_m^{n+1}=\bigl(1+\tfrac12r\delta^2\bigr)U_m^n .
\]

\step{Truncation error}
\(T=u_m^{n+1}-u_m^n-\frac r2\bigl(\delta^2u_m^{n+1}+\delta^2u_m^n\bigr)\).
Expand about \((mh,(n+\frac12)k)\), writing \(u\) for the value
there:
\[
  u_m^{n+1}-u_m^n=ku_t+\frac{k^3}{24}u_{ttt}+\cdots,
  \qquad
  \tfrac12\bigl(u^{n+1}+u^n\bigr)=u+\frac{k^2}8u_{tt}+\cdots .
\]
Since \(\delta^2v=h^2v_{xx}+\frac{h^4}{12}v_{xxxx}+\cdots\),
\[
  \frac r2\bigl(\delta^2u^{n+1}+\delta^2u^n\bigr)
  =r\,\delta^2\Bigl(u+\frac{k^2}8u_{tt}+\cdots\Bigr)
  =ku_{xx}+\frac{kh^2}{12}u_{xxxx}+\frac{k^3}8u_{xxtt}+\cdots .
\]
Hence
\[
  T=k(u_t-u_{xx})+\frac{k^3}{24}u_{ttt}-\frac{k^3}8u_{xxtt}-\frac{kh^2}{12}u_{xxxx}+\cdots .
\]
With \(u_t=u_{xx}\) we have \(u_{xxtt}=u_{ttt}\), so
\[
  T=\Bigl(\frac1{24}-\frac18\Bigr)k^3u_{ttt}-\frac{kh^2}{12}u_{xxxx}
  =-\frac{k^3}{12}u_{ttt}-\frac{kh^2}{12}u_{xxxx}+\cdots .
\]
\(T/k=O(k^2+h^2)\): Crank--Nicolson is second order in both time and
space, compared with \(O(k+h^2)\) for the explicit scheme.
\end{solution}

%---------------------------------------------------------------------
\Q{One Crank--Nicolson step}{Numerical}{8}\label{q:old23}
\begin{qbox}
For \(u_t=u_{xx}\), \(u(x,0)=\sin\pi x\), \(u(0,t)=u(1,t)=0\), take one
Crank--Nicolson step with \(h=\frac14\), \(r=1\).  Solve the system
and compare with the exact solution.
\end{qbox}
\begin{rem}
\begin{itemize}
  \item Crank--Nicolson amplification factor
    \(\xi=\frac{1-2rs^2}{1+2rs^2}\), \(|\xi|<1\) for all \(r>0\).
  \item Each step: a tridiagonal system (Thomas); use symmetry of the
    data to reduce the unknowns.
\end{itemize}
\end{rem}
\begin{solution}
Crank--Nicolson,
\((1+r)U_m^{1}-\frac r2(U_{m+1}^1+U_{m-1}^1)=(1-r)U_m^0+\frac r2(U_{m+1}^0+U_{m-1}^0)\),
with \(r=1\) reads
\(2U_m^{1}-\frac12(U_{m+1}^1+U_{m-1}^1)=\frac12(U_{m+1}^0+U_{m-1}^0)\).
Right-hand sides: \(m=1\): \(\frac12(1+0)=0.5\); \(m=2\):
\(\frac12(0.707107\cdot2)=0.707107\); \(m=3\): \(0.5\).
\[
  \begin{pmatrix}2&-\frac12&0\\-\frac12&2&-\frac12\\0&-\frac12&2\end{pmatrix}
  \begin{pmatrix}U_1\\U_2\\U_3\end{pmatrix}
  =\begin{pmatrix}0.5\\0.707107\\0.5\end{pmatrix}.
\]
By symmetry \(U_1=U_3\).  Row 2: \(U_2=\frac12(0.707107+U_1)\).  Row 1:
\(2U_1-\frac14(0.707107+U_1)=0.5\Rightarrow\frac74U_1=0.676777\), so
\[
  U_1=U_3=0.386730,\qquad U_2=0.546918 .
\]
Check: \(U_2=\xi\) with
\(\xi=\frac{1-2r\sin^2(\pi/8)}{1+2r\sin^2(\pi/8)}=0.546918\).  Exact at
\(t=k=\frac1{16}\), \(x=\frac12\): \(e^{-\pi^2/16}=0.539641\); error
\(7.3\times10^{-3}\), even though \(r=1\) is twice the explicit limit.
\end{solution}

%---------------------------------------------------------------------
\Q{Truncation error of the fully implicit scheme}{Derivation}{6}\label{q:old22}
\begin{qbox}
Find the principal part of the local truncation error of
\(U_m^{n+1}-U_m^n=r\,\delta^2U_m^{n+1}\) for \(u_t=u_{xx}\), and state
its order of consistency.
\end{qbox}
\begin{rem}
\begin{itemize}
  \item Fully implicit: \((1-r\delta^2)U_m^{n+1}=U_m^n\); expand about
    the new level \((mh,(n+1)k)\).
  \item Consistent if \(T/k\to0\); order \((p,q)\) if
    \(T/k=O(k^p+h^q)\).
\end{itemize}
\end{rem}
\begin{solution}
\(T=u_m^{n+1}-u_m^n-r\delta^2u_m^{n+1}\).  Expand about
\((mh,(n+1)k)\):
\[
  u_m^n=u-ku_t+\tfrac12k^2u_{tt}-\cdots,
  \qquad
  r\delta^2u_m^{n+1}=k\,u_{xx}+\tfrac1{12}kh^2u_{xxxx}+\cdots .
\]
Hence
\[
  T=k(u_t-u_{xx})-\tfrac12k^2u_{tt}-\tfrac1{12}kh^2u_{xxxx}+\cdots
   =-\tfrac12k^2u_{tt}-\tfrac1{12}kh^2u_{xxxx}+\cdots,
\]
so \(T=O(k^2+kh^2)\) and \(T/k=O(k+h^2)\to0\): consistent, first order
in time and second order in space.  (Unlike the explicit scheme, no
choice of \(r\) cancels the principal part, because both terms have the
same sign when \(u_{tt}=u_{xxxx}\).)
\end{solution}

%---------------------------------------------------------------------
\Q{Consistency with a derivative boundary condition}{Analysis}{7}\label{q:consbc}
\begin{qbox}
Discuss the consistency of the scheme
\[
  u_k^{n+1}=(1-2r)u_k^n+r\bigl(u_{k+1}^n+u_{k-1}^n\bigr),\quad
  k=1,\dots,M-1,\qquad r=\frac{\nu\Delta t}{\Delta x^2},
\]
\[
  u_M^{n+1}=0,\qquad u_0^n=u_1^n,
\]
with the problem \(v_t=\nu v_{xx}\) on \((0,1)\), \(v(x,0)=f(x)\),
\(v(1,t)=0\), \(v_x(0,t)=0\).  How can the boundary treatment be
improved?
\end{qbox}

\begin{rem}
\begin{itemize}
  \item A scheme for an IBVP is consistent if the difference equation
    \emph{and} each difference boundary condition are consistent with
    the PDE and the corresponding boundary condition.
  \item One-sided \(\frac{u_1-u_0}{\Delta x}\) approximates \(v_x(0)\)
    with error \(\frac{\Delta x}2v_{xx}=O(\Delta x)\); the central
    \(\frac{u_1-u_{-1}}{2\Delta x}\) with error \(O(\Delta x^2)\).
\end{itemize}
\end{rem}

\begin{solution}
Let \(v\) be a smooth exact solution and \(v_k^n=v(k\Delta x,n\Delta t)\).

\step{Interior points}
In the form approximating the PDE,
\begin{align*}
  \tau_k^n&=\frac{v_k^{n+1}-v_k^n}{\Delta t}
  -\nu\frac{v_{k+1}^n-2v_k^n+v_{k-1}^n}{\Delta x^2}\\
  &=(v_t-\nu v_{xx})+\frac{\Delta t}2v_{tt}-\frac{\nu\Delta x^2}{12}v_{xxxx}+\cdots
  =O(\Delta t)+O(\Delta x^2),
\end{align*}
which tends to zero: the interior scheme is consistent, of order
\((1,2)\).

\step{Boundary \(x=1\)}
\(u_M^{n+1}=0\) imposes \(v(1,t)=0\) exactly: zero truncation error.

\step{Boundary \(x=0\)}
\(u_0^n=u_1^n\) is the approximation \(\frac{u_1^n-u_0^n}{\Delta x}=0\)
of \(v_x(0,t)=0\).  Its truncation error is
\[
  \frac{v(\Delta x,t)-v(0,t)}{\Delta x}-v_x(0,t)
  =\frac{\Delta x}2v_{xx}(0,t)+O(\Delta x^2)=O(\Delta x)\to0 .
\]
So the boundary condition is consistent, but only first order.

\step{Conclusion and improvement}
The scheme is consistent with the IBVP.  The first-order boundary
condition in general reduces the overall accuracy to
\(O(\Delta x)\).  A second-order treatment uses a fictitious point:
\(\frac{u_1^n-u_{-1}^n}{2\Delta x}=0\Rightarrow u_{-1}^n=u_1^n\), and the
scheme at \(k=0\) becomes \(u_0^{n+1}=(1-2r)u_0^n+2ru_1^n\).
\end{solution}

%---------------------------------------------------------------------
\Q{Convergence of the explicit scheme}{Proof}{10}\label{q:old26}
\begin{qbox}
Prove that for \(u_t=u_{xx}\) on \([0,1]\) with Dirichlet conditions,
the explicit scheme converges for \(0<r\le\frac12\), with
\(\max_m|u_m^n-U_m^n|\le T\,A\,(k+h^2)\).
\end{qbox}
\begin{rem}
\begin{itemize}
  \item Error \(z_m^n=u_m^n-U_m^n\) satisfies the scheme plus the
    truncation error \(T_m^n=O(k^2+kh^2)\).
  \item For \(r\le\frac12\) the coefficients \(1-2r,\ r,\ r\) are
    non-negative and add to \(1\) (maximum principle).
  \item Convergence: \(\max|z|\to0\) as \(h,k\to0\) with \(nk=T\) fixed.
\end{itemize}
\end{rem}
\begin{solution}
Let \(z_m^n=u_m^n-U_m^n\).  Subtracting the scheme from the equation
satisfied by the exact solution,
\[
  z_m^{n+1}=(1-2r)z_m^n+r(z_{m+1}^n+z_{m-1}^n)+T_m^n,
  \qquad z_0^n=z_N^n=0,
\]
with \(|T_m^n|\le A(k^2+kh^2)\), where
\(A=\max(\frac12\max|u_{tt}|,\frac1{12}\max|u_{xxxx}|)\).  For
\(r\le\frac12\) all coefficients are non-negative and sum to \(1\), so
with \(Z^{(n)}=\max_m|z_m^n|\),
\[
  Z^{(n+1)}\le Z^{(n)}+A(k^2+kh^2).
\]
Since \(Z^{(0)}=0\),
\(Z^{(n)}\le nA(k^2+kh^2)=nk\,A(k+h^2)=T\,A(k+h^2)\to0\).
\end{solution}

%---------------------------------------------------------------------
\Q{Convergence of the fully implicit scheme for all \(r\)}{Proof}{8}\label{q:implconv}
\begin{qbox}
(Mitchell--Griffiths, Ch.~2, Exercise 5.)  Establish the convergence of
the backward difference implicit formula
\((1-r\delta^2)U_m^{n+1}=U_m^n\) for \(u_t=u_{xx}\) on \([0,1]\) with
Dirichlet conditions, for all \(r>0\).
\end{qbox}

\begin{rem}
\begin{itemize}
  \item Fully implicit:
    \((1+2r)U_m^{n+1}=U_m^n+r\bigl(U_{m+1}^{n+1}+U_{m-1}^{n+1}\bigr)\).
  \item LTE (Q\ref{q:old22}):
    \(T=-\frac12k^2u_{tt}-\frac1{12}kh^2u_{xxxx}+\cdots\), so
    \(|T_m^n|\le A(k^2+kh^2)\).
  \item Same maximum-principle argument as for the explicit scheme
    (Q\ref{q:old26}), but taken at the \emph{new} level.
\end{itemize}
\end{rem}

\begin{solution}
Let \(z_m^n=u_m^n-U_m^n\), with \(z_0^n=z_M^n=0\) and \(z_m^0=0\).  The
exact solution satisfies the scheme up to \(T_m^n\), so
\[
  (1+2r)z_m^{n+1}=z_m^n+r\bigl(z_{m+1}^{n+1}+z_{m-1}^{n+1}\bigr)+T_m^n,
  \qquad |T_m^n|\le A(k^2+kh^2).
\]
Let \(Z^{(n)}=\max_m|z_m^n|\), and choose \(m^*\) with
\(|z_{m^*}^{n+1}|=Z^{(n+1)}\).  Taking moduli at \(m^*\),
\[
  (1+2r)Z^{(n+1)}\le Z^{(n)}+r\bigl(Z^{(n+1)}+Z^{(n+1)}\bigr)+A(k^2+kh^2),
\]
so, for every \(r>0\),
\[
  Z^{(n+1)}\le Z^{(n)}+A(k^2+kh^2).
\]
Since \(Z^{(0)}=0\),
\[
  Z^{(n)}\le nA(k^2+kh^2)=T\,A\,(k+h^2)\to0
  \quad\text{as }h,k\to0,\ nk=T .
\]
So the fully implicit scheme converges for every \(r>0\).  No
restriction on \(r\) appears, because the coefficient \(1+2r\) of the
unknown dominates the \(2r\) of its neighbours for every \(r\).
\end{solution}

%---------------------------------------------------------------------
\Q{Why ``unstable'' even when the solution decays}{Conceptual}{6}\label{q:old19}
\begin{qbox}
Repeat Q\ref{q:old18} with \(r=1\).  The computed solution still
decays.  Explain why the explicit scheme is nevertheless called unstable
for \(r=1\), and what would be seen in practice.
\end{qbox}
\begin{rem}
\begin{itemize}
  \item Stability is a property of the scheme for \emph{all} data,
    i.e.\ all Fourier modes, not of one solution.
  \item Lax equivalence: for a consistent scheme, stable \(\iff\)
    convergent.
\end{itemize}
\end{rem}
\begin{solution}
With \(h=\frac14\) the grid supports the modes \(\sin(s\pi x)\),
\(s=1,2,3\), with amplification factors
\(\xi_s=1-4r\sin^2\frac{s\pi}8\).  For \(r=1\):
\[
  \xi_1=1-4\sin^2\tfrac\pi8=0.414,\qquad
  \xi_2=1-4\sin^2\tfrac\pi4=-1,\qquad
  \xi_3=1-4\sin^2\tfrac{3\pi}8=-2.414 .
\]
The initial data contain only the mode \(s=1\), which decays.  But
rounding errors introduce small amounts of every mode, and the mode
\(s=3\) is multiplied by \(-2.414\) at every step: after \(20\) steps
it has grown by \((2.414)^{20}\approx4.5\times10^7\).  In practice a
saw-tooth oscillation appears and swamps the solution.  Stability is a
property of the scheme for \emph{all} data (all modes), not of one
solution; here \(\max_s|\xi_s|>1\), which violates \(r\le\frac12\).
\end{solution}

%---------------------------------------------------------------------
\Q{The \(\theta\)-method}{Stability analysis}{10}\label{q:old20}
\begin{qbox}
For \(0\le\theta\le1\), consider
\[
  U_m^{n+1}-U_m^n=r\bigl[\theta\,\delta^2U_m^{n+1}+(1-\theta)\,\delta^2U_m^n\bigr].
\]
Find the amplification factor and show that the scheme is stable for
all \(r>0\) if \(\theta\ge\frac12\), and for
\(r\le\frac1{2(1-2\theta)}\) if \(\theta<\frac12\).  Identify the
explicit, Crank--Nicolson and fully implicit schemes.
\end{qbox}
\begin{rem}
\begin{itemize}
  \item Von Neumann: put \(U_m^n=\xi^ne^{i\beta mh}\); stable if
    \(|\xi|\le1\) for all \(\beta\) (\(|\xi|\le1+O(k)\) if the exact
    solution may grow).
  \item \(\delta^2e^{i\beta mh}=-4s^2e^{i\beta mh}\),
    \(s=\sin\frac{\beta h}2\).
  \item \(|\frac{1-a}{1+b}|\le1\) needs \(-1\le\frac{1-a}{1+b}\).
\end{itemize}
\end{rem}
\begin{solution}
Substituting \(U_m^n=\xi^ne^{i\beta mh}\) and using
\(\delta^2e^{i\beta mh}=-4s^2e^{i\beta mh}\),
\[
  \xi-1=-4rs^2\bigl[\theta\xi+(1-\theta)\bigr]
  \ \Longrightarrow\
  \xi=\frac{1-4(1-\theta)rs^2}{1+4\theta rs^2}.
\]
Since \(rs^2\ge0\), \(\xi\le1\) always.  The condition \(\xi\ge-1\) is
\[
  1-4(1-\theta)rs^2\ge-1-4\theta rs^2
  \iff 4rs^2(1-2\theta)\le2 .
\]
If \(\theta\ge\frac12\) the left side is \(\le0\): stable for all \(r\).
If \(\theta<\frac12\), the worst case is \(s^2=1\), giving
\(r\le\frac1{2(1-2\theta)}\).

\(\theta=0\): explicit scheme, \(r\le\frac12\).
\(\theta=\frac12\): Crank--Nicolson, unconditionally stable.
\(\theta=1\): fully implicit, unconditionally stable.
(For example \(\theta=\frac14\) requires \(r\le1\).)
\end{solution}

%---------------------------------------------------------------------
\Q{A growing solution: \(u_t=u_{xx}+cu\)}{Stability analysis}{6}\label{q:old21}
\begin{qbox}
For \(u_t=u_{xx}+cu\) with \(c>0\), the explicit scheme is
\(U_m^{n+1}=U_m^n+r\delta^2U_m^n+ck\,U_m^n\).  Show that \(|\xi|\) can
exceed \(1\), but the scheme is still stable when \(r\le\frac12\).
\end{qbox}
\begin{rem}
\begin{itemize}
  \item Lax--Richtmyer stability: \(\|A^n\|\le M\) (or
    \(|\xi|^n\le Ke^{\beta nk}\)) for \(nk\le T\), independent of
    \(h,k\).
  \item \(1+ck\le e^{ck}\).
\end{itemize}
\end{rem}
\begin{solution}
\(\xi=1-4rs^2+ck\).  For \(s=0\), \(\xi=1+ck>1\).  If \(r\le\frac12\),
then \(1-4rs^2\in[-1,1]\), so
\[
  |\xi|\le1+ck\le e^{ck}
  \quad\Longrightarrow\quad
  |\xi|^n\le e^{cnk}=e^{cT}.
\]
The powers are bounded by a constant independent of \(h\) and \(k\)
(for \(nk\le T\)), which is stability in the Lax--Richtmyer sense
(\(K=1\), \(\beta=c\)).  The condition \(|\xi|\le1\) is only
appropriate when the exact solution does not grow; here the exact modes
grow like \(e^{ct}\), so \(|\xi|\le1+O(k)\) is the correct requirement.
\end{solution}

%---------------------------------------------------------------------
\Q{Crank--Nicolson for \(v_t=v_{xx}-v\)}{Stability analysis}{7}\label{q:cnlower}
\begin{qbox}
Consider \(v_t=v_{xx}-v\), \(t>0\), \(x\in(-\infty,\infty)\).  Write the
Crank--Nicolson scheme and analyse its stability.
\end{qbox}

\begin{rem}
\begin{itemize}
  \item Crank--Nicolson averages the whole right-hand side over levels
    \(n\) and \(n+1\).
  \item \(\delta^2e^{i\beta mh}=-4s^2e^{i\beta mh}\).
  \item \(|1-a|<1+a\) for every \(a>0\).
\end{itemize}
\end{rem}

\begin{solution}
\step{The scheme}
\[
  \frac{U_m^{n+1}-U_m^n}k
  =\frac12\Bigl[\frac{\delta^2U_m^{n+1}}{h^2}-U_m^{n+1}\Bigr]
  +\frac12\Bigl[\frac{\delta^2U_m^n}{h^2}-U_m^n\Bigr],
\]
i.e.\ with \(r=k/h^2\),
\[
  \Bigl(1+\frac k2-\frac r2\delta^2\Bigr)U_m^{n+1}
  =\Bigl(1-\frac k2+\frac r2\delta^2\Bigr)U_m^n .
\]

\step{Amplification factor}
With \(U_m^n=\xi^ne^{i\beta mh}\),
\[
  \xi=\frac{1-\frac k2-2rs^2}{1+\frac k2+2rs^2}
  =\frac{1-a}{1+a},\qquad a=\frac k2+2rs^2>0 .
\]

\step{Stability}
For \(a>0\), \(|1-a|<1+a\), so \(|\xi|<1\) for all \(\beta\), all
\(r>0\) and all \(k>0\).  The scheme is unconditionally stable.  This
is consistent with the PDE: every Fourier mode
\(e^{-(\beta^2+1)t}e^{i\beta x}\) of the exact solution decays, and the
term \(-v\) only adds damping.
\end{solution}

%---------------------------------------------------------------------
\Q{The matrix method: norm, spectral radius and the explicit scheme}{Proof}{8}\label{q:matexp}
\begin{qbox}
\begin{enumerate}[label=(\alph*)]
  \item For \(\mathbf U_{j+1}=A\mathbf U_j+\mathbf f_j\), show that an
    initial perturbation propagates as \(\mathbf e_j=A^j\mathbf e_0\),
    and that \(\|A\|\le1\) implies stability.  Show by the example
    \(A=\begin{pmatrix}-0.8&0\\0.4&0.7\end{pmatrix}\) that
    \(\rho(A)<1\) does not imply \(\|A\|\le1\).
  \item Prove that the explicit scheme with Dirichlet conditions is
    stable in the matrix sense if and only if \(0<r\le\frac12\).
\end{enumerate}
\end{qbox}

\begin{rem}
\begin{itemize}
  \item Lax--Richtmyer: \(\|A^j\|\le M\) for \(jk\le T\).
  \item \(\rho(A)\le\|A\|\) always; for real symmetric \(A\),
    \(\|A\|_2=\rho(A)\).
  \item \(\mathrm{tridiag}(c,a,b)\) of order \(N-1\):
    \(\lambda_s=a+2\sqrt{bc}\cos\frac{s\pi}N\), \(s=1,\dots,N-1\).
\end{itemize}
\end{rem}

\begin{solution}
\step{(a) Propagation of errors}
Repeated substitution gives
\(\mathbf U_j=A^j\mathbf U_0+A^{j-1}\mathbf f_0+\dots+\mathbf f_{j-1}\).
Replacing \(\mathbf U_0\) by \(\mathbf U_0^*\) changes only the first
term, so \(\mathbf e_j=\mathbf U_j^*-\mathbf U_j=A^j\mathbf e_0\) and
\(\|\mathbf e_j\|\le\|A^j\|\|\mathbf e_0\|\le\|A\|^j\|\mathbf e_0\|\).
If \(\|A\|\le1\), then \(\|A^j\|\le1\): Lax--Richtmyer stable with
\(M=1\).

For the example, the matrix is triangular, so its eigenvalues are
\(-0.8\) and \(0.7\), and \(\rho(A)=0.8<1\).  But
\(\|A\|_1=\max(0.8+0.4,\ 0.7)=1.2>1\) and
\(\|A\|_\infty=\max(0.8,\ 1.1)=1.1>1\).  \(A\) is not symmetric (not
even normal), so its norm is not controlled by \(\rho(A)\).

\step{(b) The explicit scheme}
With known boundary values,
\(\mathbf U_{j+1}=A\mathbf U_j+\mathbf f_j\), where
\(A=\mathrm{tridiag}(r,\,1-2r,\,r)\) of order \(N-1\).  Its eigenvalues
are
\[
  \lambda_s=(1-2r)+2r\cos\frac{s\pi}N=1-4r\sin^2\frac{s\pi}{2N},
  \qquad s=1,\dots,N-1 .
\]
\(A\) is real symmetric, so \(\|A^j\|_2=\rho(A)^j\), and stability
holds if and only if \(-1\le1-4r\sin^2\frac{s\pi}{2N}\le1\) for all
\(s\).  The right inequality always holds.  The left one gives
\(r\le\bigl(2\sin^2\frac{s\pi}{2N}\bigr)^{-1}\); since
\(\sin^2\frac{(N-1)\pi}{2N}\to1\) as \(N\to\infty\), stability for all
\(N\) requires
\[
  \boxed{0<r\le\tfrac12},
\]
the same as the von Neumann condition.
\end{solution}

%---------------------------------------------------------------------
\Q{Matrix facts used in stability proofs}{Proof}{8}\label{q:matfacts}
\begin{qbox}
(Mitchell--Griffiths, Ch.~1, Examples 3, 6, 7 and Ch.~2, Example 5.)
Prove:
\begin{enumerate}[label=(\alph*)]
  \item the eigenvalues of \(A\) are preserved under a similarity
    transformation \(H^{-1}AH\) (\(|H|\ne0\));
  \item if \(A\) is symmetric and positive real
    (\((\mathbf x,A\mathbf x)>0\) for all real \(\mathbf x\ne0\)), its
    eigenvalues are all positive;
  \item \(A^TA\) has non-negative eigenvalues;
  \item a calculation \(\mathbf U^{n+1}=C\mathbf U^n\) (\(C\) constant)
    is stable if \(\rho(C)\le1\) and \(C\) is similar to a symmetric
    matrix, \(\widetilde C=P^{-1}CP\).
\end{enumerate}
\end{qbox}

\begin{rem}
\begin{itemize}
  \item \(\det(XY)=\det X\det Y\).
  \item A real symmetric matrix has real eigenvalues and real
    orthonormal eigenvectors, and \(\|S\|_2=\rho(S)\).
  \item Used in Q\ref{q:cnrobin} (\(D^{-1}QD\) symmetric) and for the
    Robin problem in Section~\ref{sec:adi}.
\end{itemize}
\end{rem}

\begin{solution}
\step{(a)}
\(\det(H^{-1}AH-\lambda I)=\det\bigl(H^{-1}(A-\lambda I)H\bigr)
=\det H^{-1}\det(A-\lambda I)\det H=\det(A-\lambda I)\).  The two
characteristic polynomials are equal, so the eigenvalues (with
multiplicities) are the same.  Eigenvectors change: if
\(A\mathbf x=\lambda\mathbf x\), then
\((H^{-1}AH)(H^{-1}\mathbf x)=\lambda H^{-1}\mathbf x\).

\step{(b)}
\(A\) symmetric, so each eigenvalue \(\lambda\) is real with a real
eigenvector \(\mathbf v\ne0\).  Then
\((\mathbf v,A\mathbf v)=\lambda(\mathbf v,\mathbf v)\), so
\(\lambda=\dfrac{(\mathbf v,A\mathbf v)}{(\mathbf v,\mathbf v)}>0\).

\step{(c)}
\(A^TA\) is symmetric.  For an eigenpair
\(A^TA\mathbf v=\mu\mathbf v\) (\(\mathbf v\) real),
\(\mu\|\mathbf v\|^2=\mathbf v^TA^TA\mathbf v=\|A\mathbf v\|^2\ge0\), so
\(\mu\ge0\).  (This is why \(\|A\|_2=\sqrt{\rho(A^TA)}\) makes sense.)

\step{(d)}
\(C=P\widetilde CP^{-1}\), so \(C^n=P\widetilde C^nP^{-1}\).  By (a),
\(\rho(\widetilde C)=\rho(C)\le1\), and \(\widetilde C\) is symmetric, so
\(\|\widetilde C^n\|_2=\rho(\widetilde C)^n\le1\).  Hence
\[
  \|C^n\|_2\le\|P\|_2\,\|\widetilde C^n\|_2\,\|P^{-1}\|_2\le\|P\|_2\|P^{-1}\|_2
\]
for all \(n\): the powers of \(C\) are uniformly bounded, which is
stability (provided \(\|P\|\|P^{-1}\|\) is bounded independently of
\(h\), as for \(P=D=\mathrm{diag}(\sqrt2,1,\dots,1,\sqrt2)\) in
Q\ref{q:cnrobin}).
\end{solution}

%---------------------------------------------------------------------
\Q{Spectral radius versus norm: calculations}{Numerical / proof}{8}\label{q:normcalc}
\begin{qbox}
(Mitchell--Griffiths, Ch.~1, Exercises 4--8.)
\begin{enumerate}[label=(\alph*)]
  \item Show that \(\rho(A^2)=\rho(A)^2\).
  \item Find the spectral radius and spectral norm of
    \(\begin{pmatrix}1&0&\alpha\\0&1&0\\0&0&1\end{pmatrix}\) and show that
    they are equal only if \(\alpha=0\).
  \item Show that \(\|A\|_2=\rho(A)\) for
    \(A=\begin{pmatrix}2&1\\1&2\end{pmatrix}\), but \(\|A\|_2>\rho(A)\)
    for \(A=\begin{pmatrix}2&1\\0&2\end{pmatrix}\).
  \item Show that \(\|A\|_2^2\le\|A\|_\infty\|A\|_1\).
  \item Show by calculation that
    \(\begin{pmatrix}\frac12&0&0\\1&\frac12&0\\0&1&\frac12\end{pmatrix}^r\to0\)
    as \(r\to\infty\).
\end{enumerate}
\end{qbox}

\begin{rem}
\begin{itemize}
  \item \(\rho(A)=\max|\lambda_i|\);
    \(\|A\|_2=\sqrt{\rho(A^TA)}\); \(\|A\|_1\) = largest column sum,
    \(\|A\|_\infty\) = largest row sum (of moduli).
  \item \(\rho(A)\le\|A\|\) for every induced norm; equality in
    \(\|\cdot\|_2\) for symmetric (normal) matrices.
  \item \(A^r\to0\iff\rho(A)<1\).
\end{itemize}
\end{rem}

\begin{solution}
\step{(a)}
If \(A\mathbf x=\lambda\mathbf x\), then \(A^2\mathbf x=\lambda^2\mathbf x\).
More precisely, by the Schur form \(A=QTQ^*\) (\(T\) upper triangular
with diagonal \(\lambda_i\)), \(A^2=QT^2Q^*\) and \(T^2\) has diagonal
\(\lambda_i^2\).  So the eigenvalues of \(A^2\) are exactly
\(\lambda_i^2\), and \(\rho(A^2)=\max|\lambda_i|^2=\rho(A)^2\).

\step{(b)}
The matrix is triangular with unit diagonal, so \(\rho=1\).  Next,
\[
  A^TA=\begin{pmatrix}1&0&\alpha\\0&1&0\\ \alpha&0&1+\alpha^2\end{pmatrix},
\]
with eigenvalue \(1\) and the eigenvalues of
\(\begin{pmatrix}1&\alpha\\ \alpha&1+\alpha^2\end{pmatrix}\) (trace
\(2+\alpha^2\), determinant \(1\)):
\(\mu=\frac12\bigl(2+\alpha^2\pm|\alpha|\sqrt{\alpha^2+4}\bigr)\).  Hence
\[
  \|A\|_2=\sqrt{\mu_{\max}}=\frac{|\alpha|+\sqrt{\alpha^2+4}}2,
\]
(check by squaring), which equals \(1=\rho(A)\) only if \(\alpha=0\).

\step{(c)}
\(\begin{pmatrix}2&1\\1&2\end{pmatrix}\) is symmetric with eigenvalues
\(1,3\): \(\|A\|_2=\rho(A)=3\).  For
\(\begin{pmatrix}2&1\\0&2\end{pmatrix}\), \(\rho=2\), but
\(A^TA=\begin{pmatrix}4&2\\2&5\end{pmatrix}\) has eigenvalues
\(\frac{9\pm\sqrt{17}}2\), so
\(\|A\|_2=\sqrt{(9+\sqrt{17})/2}=2.5616>2\).

\step{(d)}
\(\|A\|_2^2=\rho(A^TA)\le\|A^TA\|_\infty\le\|A^T\|_\infty\|A\|_\infty
=\|A\|_1\|A\|_\infty\), since the row sums of \(A^T\) are the column
sums of \(A\).

\step{(e)}
Write the matrix as \(\frac12I+N\), with \(N\) having ones on the
subdiagonal; \(N^2\) has a single \(1\) in position \((3,1)\), and
\(N^3=0\).  Since \(I\) and \(N\) commute, the binomial theorem gives
\[
  \Bigl(\tfrac12I+N\Bigr)^r=2^{-r}I+r\,2^{1-r}N+r(r-1)\,2^{1-r}N^2
  =\begin{pmatrix}2^{-r}&0&0\\ r2^{1-r}&2^{-r}&0\\ r(r-1)2^{1-r}&r2^{1-r}&2^{-r}\end{pmatrix}.
\]
Every entry tends to zero because \(r^22^{-r}\to0\).  (Here
\(\rho=\frac12<1\), although \(\|\cdot\|_1=\|\cdot\|_\infty=\frac32>1\).)
\end{solution}

%---------------------------------------------------------------------
\Q{Eigenvalues of a tridiagonal matrix}{Proof}{8}\label{q:tridproof}
\begin{qbox}
Let \(A\) be the \(N\times N\) tridiagonal matrix with \(a\) on the
diagonal, \(b\) above it and \(c\) below it, \(bc\ne0\).  Prove that
its eigenvalues are
\[
  \lambda_s=a+2b\Bigl(\frac cb\Bigr)^{1/2}\cos\frac{s\pi}{N+1},
  \qquad s=1,\dots,N,
\]
and find the eigenvectors.
\end{qbox}

\begin{rem}
\begin{itemize}
  \item Write \(A\mathbf v=\lambda\mathbf v\) as a difference equation
    with \(v_0=v_{N+1}=0\).
  \item \(v_j=m^j\) gives \(bm^2+(a-\lambda)m+c=0\): product of roots
    \(m_1m_2=c/b\), sum \(m_1+m_2=(\lambda-a)/b\).
  \item For \(b,c>0\):
    \(\lambda_s=a+2\sqrt{bc}\cos\frac{s\pi}{N+1}\).
\end{itemize}
\end{rem}

\begin{solution}
\step{Step 1: difference equation}
\(A\mathbf v=\lambda\mathbf v\), row by row, with \(v_0=v_{N+1}=0\):
\[
  cv_{j-1}+(a-\lambda)v_j+bv_{j+1}=0,\qquad j=1,\dots,N .
\]
\step{Step 2: general solution}
\(v_j=m^j\) gives \(bm^2+(a-\lambda)m+c=0\), with roots \(m_1,m_2\).
If \(m_1\ne m_2\), \(v_j=Bm_1^j+Cm_2^j\).

\step{Step 3: equal roots are impossible}
If \(m_1=m_2\), \(v_j=(B+Cj)m_1^j\); \(v_0=0\Rightarrow B=0\), and
\(v_{N+1}=0\Rightarrow C=0\) (\(m_1\ne0\) since \(m_1m_2=c/b\ne0\)).
Then \(\mathbf v=0\), impossible for an eigenvector.

\step{Step 4: boundary conditions}
\(v_0=0\Rightarrow C=-B\) (\(B\ne0\)); \(v_{N+1}=0\Rightarrow
m_1^{N+1}=m_2^{N+1}\), so
\[
  \Bigl(\frac{m_1}{m_2}\Bigr)^{N+1}=1
  \ \Longrightarrow\
  \frac{m_1}{m_2}=e^{2is\pi/(N+1)},\qquad s=1,\dots,N
\]
(\(s=0\) is excluded because \(m_1\ne m_2\)).

\step{Step 5: the roots}
With \(m_1m_2=c/b\):
\(m_1=\bigl(\frac cb\bigr)^{1/2}e^{is\pi/(N+1)}\),
\(m_2=\bigl(\frac cb\bigr)^{1/2}e^{-is\pi/(N+1)}\).

\step{Step 6: eigenvalues}
\(m_1+m_2=\frac{\lambda-a}b\), so
\[
  \lambda_s=a+b\Bigl(\frac cb\Bigr)^{1/2}\bigl(e^{is\pi/(N+1)}+e^{-is\pi/(N+1)}\bigr)
  =a+2b\Bigl(\frac cb\Bigr)^{1/2}\cos\frac{s\pi}{N+1}.
\]
\step{Step 7: eigenvectors}
\(v_j=B(m_1^j-m_2^j)=2iB\bigl(\frac cb\bigr)^{j/2}\sin\frac{js\pi}{N+1}\);
with \(2iB=1\),
\[
  \mathbf v_s=\Bigl[\Bigl(\tfrac cb\Bigr)^{1/2}\sin\tfrac{s\pi}{N+1},\
  \Bigl(\tfrac cb\Bigr)\sin\tfrac{2s\pi}{N+1},\ \dots,\
  \Bigl(\tfrac cb\Bigr)^{N/2}\sin\tfrac{Ns\pi}{N+1}\Bigr]^T .
\]
The \(N\) values of \(\lambda_s\) are distinct, so these are all the
eigenvalues.
\end{solution}

%---------------------------------------------------------------------
\Q{Eigenvectors of \(\mathrm{tridiag}(1,-2,1)\) of order 3}{Numerical}{4}\label{q:eig3}
\begin{qbox}
(Mitchell--Griffiths, Ch.~1, Exercise 1.)  Find the eigenvalues and
eigenvectors of
\(\begin{pmatrix}-2&1&0\\1&-2&1\\0&1&-2\end{pmatrix}\) and show that the
eigenvectors are mutually orthogonal.
\end{qbox}

\begin{rem}
\begin{itemize}
  \item Q\ref{q:tridproof} with \(a=-2\), \(b=c=1\), \(N=3\):
    \(\lambda_s=-2+2\cos\frac{s\pi}4\),
    \((\mathbf v_s)_j=\sin\frac{js\pi}4\).
  \item Eigenvectors of a symmetric matrix for distinct eigenvalues
    are orthogonal.
\end{itemize}
\end{rem}

\begin{solution}
\(\lambda_s=-2+2\cos\frac{s\pi}4\), \(s=1,2,3\):
\[
  \lambda_1=-2+\sqrt2,\qquad\lambda_2=-2,\qquad\lambda_3=-2-\sqrt2 .
\]
Eigenvectors \((\sin\frac{s\pi}4,\sin\frac{2s\pi}4,\sin\frac{3s\pi}4)\),
scaled:
\[
  \mathbf v_1=(1,\sqrt2,1)^T,\qquad
  \mathbf v_2=(1,0,-1)^T,\qquad
  \mathbf v_3=(1,-\sqrt2,1)^T .
\]
Check for \(\mathbf v_1\): \(A\mathbf v_1=(-2+\sqrt2,\ 2-2\sqrt2,\ -2+\sqrt2)^T
=(-2+\sqrt2)(1,\sqrt2,1)^T\).  Orthogonality:
\[
  \mathbf v_1\cdot\mathbf v_2=1-1=0,\qquad
  \mathbf v_1\cdot\mathbf v_3=1-2+1=0,\qquad
  \mathbf v_2\cdot\mathbf v_3=1-1=0 .
\]
This is expected: the matrix is symmetric with distinct eigenvalues.
\end{solution}

%---------------------------------------------------------------------
\Q{Matrix stability of the fully implicit scheme}{Proof}{6}\label{q:old24}
\begin{qbox}
Show by the matrix method that the fully implicit scheme with Dirichlet
conditions is unconditionally stable.
\end{qbox}
\begin{rem}
\begin{itemize}
  \item Fully implicit: \(A\mathbf U^{n+1}=\mathbf U^n+\mathbf b\),
    \(A=\mathrm{tridiag}(-r,1+2r,-r)\).
  \item Eigenvalues of \(A^{-1}\) are \(1/\lambda_s(A)\); symmetric
    matrices have \(\|\cdot\|_2=\rho(\cdot)\).
\end{itemize}
\end{rem}
\begin{solution}
The scheme is \(A\mathbf U^{n+1}=\mathbf U^n+\mathbf b\) with
\(A=\mathrm{tridiag}(-r,\,1+2r,\,-r)\) of size \(N-1\).  By the
tridiagonal eigenvalue formula,
\[
  \lambda_s(A)=1+2r-2r\cos\frac{s\pi}N=1+4r\sin^2\frac{s\pi}{2N}\ \ge1 .
\]
\(A\) is real symmetric, so \(A^{-1}\) is symmetric with eigenvalues
\(1/\lambda_s(A)\in(0,1)\).  Hence
\(\|A^{-1}\|_2=\rho(A^{-1})<1\) for every \(r>0\), and
\(\|(A^{-1})^n\|_2\le1\): unconditionally stable.
\end{solution}

%---------------------------------------------------------------------
\Q{Matrix stability of Crank--Nicolson}{Proof}{6}\label{q:cnmat}
\begin{qbox}
Show by the matrix method that the Crank--Nicolson scheme with
Dirichlet boundary conditions is unconditionally stable, and hence
convergent for every \(r>0\).  Justify carefully why the eigenvalues of
\(A^{-1}B\) are ratios of eigenvalues.
\end{qbox}

\begin{rem}
\begin{itemize}
  \item Crank--Nicolson: \(A\mathbf U^{n+1}=B\mathbf U^n\), with
    \(A=\mathrm{tridiag}(-\frac r2,1+r,-\frac r2)\),
    \(B=\mathrm{tridiag}(\frac r2,1-r,\frac r2)\).
  \item In general the eigenvalues of \(A^{-1}B\) are \emph{not}
    ratios of eigenvalues; here \(A\) and \(B\) are polynomials in the
    same symmetric matrix.
  \item Lax: consistent \(+\) stable \(\Rightarrow\) convergent.
\end{itemize}
\end{rem}

\begin{solution}
Let \(T_0=\mathrm{tridiag}(1,-2,1)\) of order \(N-1\).  Then
\[
  A=I-\tfrac r2T_0,\qquad B=I+\tfrac r2T_0 .
\]
Both are polynomials in the same symmetric matrix \(T_0\), so they
commute and share its eigenvectors \(\mathbf v_s\).  The eigenvalues of
\(T_0\) are \(-2+2\cos\frac{s\pi}N=-4\sin^2\theta_s\),
\(\theta_s=\frac{s\pi}{2N}\), so
\[
  A\mathbf v_s=(1+2r\sin^2\theta_s)\mathbf v_s,\qquad
  B\mathbf v_s=(1-2r\sin^2\theta_s)\mathbf v_s,
\]
and therefore
\[
  A^{-1}B\,\mathbf v_s=\frac{1-2r\sin^2\theta_s}{1+2r\sin^2\theta_s}\,\mathbf v_s,
  \qquad s=1,\dots,N-1 .
\]
Since \(2r\sin^2\theta_s>0\), each ratio has modulus less than \(1\).
\(A^{-1}B\) is symmetric (a product of commuting symmetric matrices),
so \(\|A^{-1}B\|_2=\rho(A^{-1}B)<1\) for every \(r>0\): the scheme is
unconditionally stable.  It is also consistent
(Q\ref{q:cnlte}), so by the Lax equivalence theorem it converges for
every \(r>0\).
\end{solution}

%---------------------------------------------------------------------
\Q{Stability of Crank--Nicolson by the energy method}{Proof}{6}\label{q:energy}
\begin{qbox}
(Mitchell--Griffiths, Ch.~2, Exercise 9.)  Crank--Nicolson is used for
\(u_t=u_{xx}\) with homogeneous Dirichlet conditions,
\(Z_0^n=Z_M^n=0\).  With \(V_m=Z_m^{n+1}+Z_m^n\) and
\(\|Z\|^2=\sum_mZ_m^2\), show that
\[
  \|Z^{n+1}\|^2-\|Z^n\|^2=-\frac r2\sum_{m=0}^{M-1}\bigl(V_{m+1}-V_m\bigr)^2,
\]
and deduce directly that the method is stable for all \(r>0\).
\end{qbox}

\begin{rem}
\begin{itemize}
  \item Crank--Nicolson: \(Z^{n+1}-Z^n=\frac r2\delta^2(Z^{n+1}+Z^n)=\frac r2\delta^2V\).
  \item \((a-b)(a+b)=a^2-b^2\).
  \item Summation by parts, with \(V_0=V_M=0\):
    \(\sum_{m=1}^{M-1}V_m\,\delta^2V_m=-\sum_{m=0}^{M-1}(V_{m+1}-V_m)^2\).
\end{itemize}
\end{rem}

\begin{solution}
\step{Multiply by \(V_m\) and sum}
\(Z_m^{n+1}-Z_m^n=\frac r2\delta^2V_m\).  Multiplying by
\(V_m=Z_m^{n+1}+Z_m^n\) and summing over \(m=1,\dots,M-1\),
\[
  \sum_m\bigl[(Z_m^{n+1})^2-(Z_m^n)^2\bigr]=\|Z^{n+1}\|^2-\|Z^n\|^2
  =\frac r2\sum_{m=1}^{M-1}V_m\bigl(V_{m+1}-2V_m+V_{m-1}\bigr).
\]
\step{Summation by parts}
\[
  \sum_{m=1}^{M-1}V_m(V_{m+1}-V_m)-\sum_{m=1}^{M-1}V_m(V_m-V_{m-1}).
\]
Since \(V_0=V_M=0\), the first sum equals
\(\sum_{m=0}^{M-1}V_m(V_{m+1}-V_m)\), and the second, with \(j=m-1\),
equals \(\sum_{j=0}^{M-1}V_{j+1}(V_{j+1}-V_j)\).  Their difference is
\(\sum_{m=0}^{M-1}(V_m-V_{m+1})(V_{m+1}-V_m)=-\sum_{m=0}^{M-1}(V_{m+1}-V_m)^2\).
Hence
\[
  \|Z^{n+1}\|^2-\|Z^n\|^2=-\frac r2\sum_{m=0}^{M-1}(V_{m+1}-V_m)^2\le0 .
\]
\step{Stability}
So \(\|Z^{n+1}\|\le\|Z^n\|\le\dots\le\|Z^0\|\) for every \(r>0\): any
perturbation (or error) is never amplified, which is unconditional
stability.  No eigenvalues are needed.

(The exercise as extracted shows the factor \(r\); with
\(V_m=Z_m^{n+1}+Z_m^n\) the factor is \(\frac r2\), as derived above and
checked numerically.  The conclusion is the same.)
\end{solution}

%---------------------------------------------------------------------
\Q{A Neumann condition by a fictitious point}{Derivation}{5}\label{q:old25}
\begin{qbox}
For the explicit scheme with \(u_x(0,t)=0\), derive the equation at
\(m=0\) using a fictitious point and central differences.  Compare its
accuracy with the condition \(U_0^n=U_1^n\).
\end{qbox}
\begin{rem}
\begin{itemize}
  \item Fictitious point: use the central difference in the BC to
    eliminate \(U_{-1}\), then apply the scheme at \(m=0\).
  \item Central \(O(h^2)\) versus one-sided \(O(h)\).
\end{itemize}
\end{rem}
\begin{solution}
Central difference: \(\frac{U_1^n-U_{-1}^n}{2h}=0\Rightarrow U_{-1}^n=U_1^n\).
The scheme at \(m=0\),
\(U_0^{n+1}=(1-2r)U_0^n+r(U_1^n+U_{-1}^n)\), becomes
\[
  U_0^{n+1}=(1-2r)U_0^n+2rU_1^n .
\]
The central approximation of \(u_x\) has error \(O(h^2)\), while
\(U_0=U_1\), i.e.\ \(\frac{U_1-U_0}h=0\), has error
\(\frac h2u_{xx}(0,t)=O(h)\).  So the fictitious-point treatment keeps
the second-order spatial accuracy of the interior scheme.  The
coefficients \(1-2r\) and \(2r\) are non-negative for \(r\le\frac12\),
so the maximum-principle stability argument still applies.
\end{solution}

%---------------------------------------------------------------------
\Q{Crank--Nicolson with derivative boundary conditions}{Proof}{10}\label{q:cnrobin}
\begin{qbox}
For \(u_t=u_{xx}\) on \(0\le x\le1\) with
\(u_x-pu=f_0(t)\) at \(x=0\) and \(u_x+qu=f_1(t)\) at \(x=1\)
(\(p,q\ge0\)), apply Crank--Nicolson at all points \(m=0,\dots,M\),
eliminating fictitious values by central differences.  Show that the
scheme is \(A\mathbf U^{n+1}=B\mathbf U^n+\mathbf k^n\) with
\(A=I-\frac r2Q\), \(B=I+\frac r2Q\), find \(Q\), and prove that the
scheme is unconditionally stable.
\end{qbox}

\begin{rem}
\begin{itemize}
  \item Fictitious points: \(U_{-1}=U_1-2hpU_0-2hf_0\),
    \(U_{M+1}=U_{M-1}-2hqU_M+2hf_1\).
  \item \(Q\) is not symmetric, but \(D^{-1}QD\) is, with
    \(D=\mathrm{diag}(\sqrt2,1,\dots,1,\sqrt2)\): similar matrices have
    the same eigenvalues.
  \item \(\mu=\frac{1+x}{1-x}\) satisfies \(|\mu|\le1\iff x\le0\)
    (\(x\) real).
\end{itemize}
\end{rem}

\begin{solution}
\step{The scheme}
At the boundary points,
\(\delta^2U_0=U_1-2U_0+U_{-1}=2U_1-2(1+ph)U_0-2hf_0\) and
\(\delta^2U_M=2U_{M-1}-2(1+qh)U_M+2hf_1\).  So Crank--Nicolson at all
\(M+1\) points gives \(A\mathbf U^{n+1}=B\mathbf U^n+\mathbf k^n\), with
\(A=I-\frac r2Q\), \(B=I+\frac r2Q\),
\[
  Q=\begin{pmatrix}
    -2(1+ph)&2&&&\\
    1&-2&1&&\\
    &\ddots&\ddots&\ddots&\\
    &&1&-2&1\\
    &&&2&-2(1+qh)
  \end{pmatrix}_{(M+1)\times(M+1)}
\]
and \(\mathbf k^n=rh\bigl(-(f_0^n+f_0^{n+1}),0,\dots,0,f_1^n+f_1^{n+1}\bigr)^T\).

\step{Step 1: symmetrise}
With \(D=\mathrm{diag}(\sqrt2,1,\dots,1,\sqrt2)\),
\((D^{-1}QD)_{ij}=Q_{ij}d_j/d_i\): the corner pairs \((2,1)\) both
become \(\sqrt2\), so \(\widetilde Q=D^{-1}QD\) is symmetric.  Then
\[
  D^{-1}(A^{-1}B)D=\bigl(I-\tfrac r2\widetilde Q\bigr)^{-1}\bigl(I+\tfrac r2\widetilde Q\bigr),
\]
which is symmetric (both factors are functions of \(\widetilde Q\)).
So \(A^{-1}B\) has real eigenvalues, and in the norm
\(\|\mathbf x\|_D=\|D^{-1}\mathbf x\|_2\) its norm equals its spectral
radius: stability \(\iff\rho(A^{-1}B)\le1\).

\step{Step 2: eigenvalues}
If \(\lambda_j\) are the (real) eigenvalues of \(Q\), the eigenvalues
of \(A^{-1}B\) are
\(\mu_j=\frac{1+\frac12r\lambda_j}{1-\frac12r\lambda_j}\), and
\(|\mu_j|\le1\) for all \(r>0\) if and only if \(\lambda_j\le0\).

\step{Step 3: \(\lambda_j\le0\)}
Let \(W=D^{-2}=\mathrm{diag}(\frac12,1,\dots,1,\frac12)\); \(WQ\) is
symmetric and
\[
  -\mathbf x^TWQ\,\mathbf x=\sum_{j=0}^{M-1}(x_{j+1}-x_j)^2+ph\,x_0^2+qh\,x_M^2\ \ge0 .
\]
If \(Q\mathbf v=\lambda\mathbf v\) with \(\mathbf v\) real, then
\(\lambda=\dfrac{\mathbf v^TWQ\mathbf v}{\mathbf v^TW\mathbf v}\le0\).
Hence \(|\mu_j|\le1\) for all \(j\) and all \(r>0\): the scheme is
unconditionally stable.
\end{solution}

%---------------------------------------------------------------------
\Q{Stability of the two-dimensional explicit schemes}{Stability analysis}{7}\label{q:vn2d}
\begin{qbox}
For \(u_t=u_{x_1x_1}+u_{x_2x_2}\), show that the five-point explicit
scheme \(U^{n+1}=[1+r(\delta_{x_1}^2+\delta_{x_2}^2)]U^n\) is stable if
and only if \(r\le\frac14\), and that the nine-point scheme
\(U^{n+1}=(1+r\delta_{x_1}^2)(1+r\delta_{x_2}^2)U^n\) is stable for
\(r\le\frac12\).  What is the condition in \(d\) dimensions?
\end{qbox}

\begin{rem}
\begin{itemize}
  \item Fourier mode \(e^{i(\beta lh+\gamma mh)}\);
    \(s_1=\sin\frac{\beta h}2\), \(s_2=\sin\frac{\gamma h}2\).
  \item Both schemes have LTE \(O(k^2+kh^2)\); they come from
    truncating \(\exp(kD_1^2)\exp(kD_2^2)\).
\end{itemize}
\end{rem}

\begin{solution}
\step{Five-point scheme}
\(U_{l,m}^{n+1}=(1-4r)U_{l,m}^n+r\bigl(U_{l\pm1,m}^n+U_{l,m\pm1}^n\bigr)\).
Substituting \(U_{l,m}^n=\xi^ne^{i(\beta l+\gamma m)h}\),
\[
  \xi=1-4r+2r(\cos\beta h+\cos\gamma h)=1-4r\bigl(s_1^2+s_2^2\bigr).
\]
\(\xi\le1\) always; \(\xi\ge-1\iff r\le\frac1{2(s_1^2+s_2^2)}\).  Since
\(s_1^2+s_2^2\) can be as large as \(2\), stability holds for all modes
if and only if \(\boxed{r\le\frac14}\).

\step{Nine-point scheme}
\(\xi=(1-4rs_1^2)(1-4rs_2^2)\).  Each factor lies in \([-1,1]\) when
\(r\le\frac12\), so \(|\xi|\le1\): stable for \(r\le\frac12\), twice the
five-point time step.

\step{\(d\) dimensions}
\(\xi=1-4r\sum_{i=1}^ds_i^2\), and the sum can reach \(d\), so the
condition is \(r\le\frac1{2d}\).
\end{solution}

%---------------------------------------------------------------------
\Q{Truncation errors of the two-dimensional explicit formulas}{Derivation}{6}\label{q:te2d}
\begin{qbox}
(Mitchell--Griffiths, Ch.~2, Exercise 11.)  Obtain the truncation
errors of the two-dimensional explicit formulas
\[
  U^{n+1}=\bigl[1+r(\delta_{x_1}^2+\delta_{x_2}^2)\bigr]U^n
  \qquad\text{and}\qquad
  U^{n+1}=(1+r\delta_{x_1}^2)(1+r\delta_{x_2}^2)U^n,
\]
and show that when \(r=\frac16\) the second formula has truncation error
\(O(k^3+kh^4)\).
\end{qbox}

\begin{rem}
\begin{itemize}
  \item Exact: \(u^{n+1}=\exp(kD_1^2)\exp(kD_2^2)u^n=E_1E_2u^n\), with
    \(E_i=1+r\delta_{x_i}^2+\frac12r(r-\frac16)\delta_{x_i}^4+d\,\delta_{x_i}^6+\cdots\).
  \item \(\delta_{x_1}^2\delta_{x_2}^2u\approx h^4u_{x_1x_1x_2x_2}\);
    \(r^2h^4=k^2\).
  \item \(u_{tt}=(\nabla^2)^2u=u_{x_1^4}+2u_{x_1^2x_2^2}+u_{x_2^4}\).
\end{itemize}
\end{rem}

\begin{solution}
\step{Expansion of the exact operator}
\[
  E_1E_2=1+r(\delta_{x_1}^2+\delta_{x_2}^2)+r^2\delta_{x_1}^2\delta_{x_2}^2
  +\tfrac12r\bigl(r-\tfrac16\bigr)(\delta_{x_1}^4+\delta_{x_2}^4)+\cdots,
\]
where the dots are terms of order \(\delta^6\) (\(\sim kh^4\) and
\(k^3\) for fixed \(r\)).

\step{Five-point formula}
It omits \(r^2\delta_{x_1}^2\delta_{x_2}^2\) and the \(\delta^4\) terms, so
\[
  T=k^2u_{x_1x_1x_2x_2}+\tfrac12k^2\Bigl(1-\frac1{6r}\Bigr)\bigl(u_{x_1^4}+u_{x_2^4}\bigr)+\cdots
  =\tfrac12k^2\Bigl[u_{tt}-\frac1{6r}\bigl(u_{x_1^4}+u_{x_2^4}\bigr)\Bigr]+\cdots,
\]
using \(r\,h^4=kh^2\), \(\frac12r(r-\frac16)h^4=\frac12k^2(1-\frac1{6r})\).
Since \(u_{tt}\) contains \(2u_{x_1^2x_2^2}\), no value of \(r\) removes
the \(k^2\) term: \(T=O(k^2+kh^2)\).

\step{Nine-point (factored) formula}
\((1+r\delta_{x_1}^2)(1+r\delta_{x_2}^2)=1+r(\delta_{x_1}^2+\delta_{x_2}^2)+r^2\delta_{x_1}^2\delta_{x_2}^2\)
contains the cross term, so
\[
  T=\tfrac12k^2\Bigl(1-\frac1{6r}\Bigr)\bigl(u_{x_1^4}+u_{x_2^4}\bigr)+\cdots
  =O(k^2+kh^2).
\]
For \(r=\frac16\) this term vanishes; so do the products
\(r\delta_{x_1}^2\cdot\frac12r(r-\frac16)\delta_{x_2}^4\).  What remains
starts with \(d(\delta_{x_1}^6+\delta_{x_2}^6)\), \(d=\frac1{3240}\),
i.e.\ \(T=O(h^6)=O(k^3+kh^4)\) for fixed \(r\).
\end{solution}

%---------------------------------------------------------------------
\Q{Time-step limits in one, two and three dimensions}{Short answer}{5}\label{q:old27}
\begin{qbox}
With \(h=0.1\), what is the largest stable time step for the explicit
scheme for the heat equation in 1D, 2D (five-point) and 3D?  What is it
for the 2D nine-point scheme, and for Crank--Nicolson or ADI?
\end{qbox}
\begin{rem}
\begin{itemize}
  \item Explicit heat equation in \(d\) dimensions: \(r\le\frac1{2d}\),
    i.e.\ \(k\le\frac{h^2}{2d}\).
  \item Crank--Nicolson and ADI: unconditionally stable.
\end{itemize}
\end{rem}
\begin{solution}
In \(d\) dimensions the explicit scheme needs \(r\le\frac1{2d}\), i.e.
\(k\le\frac{h^2}{2d}\):
\[
  \text{1D: }k\le0.005,\qquad
  \text{2D: }k\le0.0025,\qquad
  \text{3D: }k\le0.00\overline{16}.
\]
The nine-point scheme has
\(\xi=(1-4rs_1^2)(1-4rs_2^2)\), stable for \(r\le\frac12\), so
\(k\le0.005\), twice the five-point limit.  Crank--Nicolson and the
Peaceman--Rachford ADI scheme are unconditionally stable: no limit (the
step is chosen for accuracy).
\end{solution}

%#####################################################################
\Needspace{32\baselineskip}
\section{Lecture 35: ADI, Three-Level and Du Fort--Frankel Schemes}
%#####################################################################

%---------------------------------------------------------------------
\Q{Peaceman--Rachford reproduces the factored scheme}{Proof}{6}\label{q:old28}
\begin{qbox}
Show that eliminating \(U^{n+1/2}\) from the Peaceman--Rachford
equations
\[
  \bigl(1-\tfrac r2\delta_{x_1}^2\bigr)U^{n+1/2}=\bigl(1+\tfrac r2\delta_{x_2}^2\bigr)U^n,
  \qquad
  \bigl(1-\tfrac r2\delta_{x_2}^2\bigr)U^{n+1}=\bigl(1+\tfrac r2\delta_{x_1}^2\bigr)U^{n+1/2}
\]
gives the two-dimensional Crank--Nicolson-type scheme.  What assumption
is needed?  Why is the split form preferred?
\end{qbox}
\begin{rem}
\begin{itemize}
  \item 2D Crank--Nicolson (factored):
    \((1-\frac r2\delta_{x_1}^2)(1-\frac r2\delta_{x_2}^2)U^{n+1}
    =(1+\frac r2\delta_{x_1}^2)(1+\frac r2\delta_{x_2}^2)U^n\).
  \item Peaceman--Rachford: first half-step implicit in \(x_1\),
    explicit in \(x_2\); second half-step the other way round.
  \item Needs \(\delta_{x_1}^2\delta_{x_2}^2=\delta_{x_2}^2\delta_{x_1}^2\)
    (constant coefficients, rectangle).
\end{itemize}
\end{rem}
\begin{solution}
Apply \(\bigl(1-\frac r2\delta_{x_1}^2\bigr)\) to the second equation.
If \(\delta_{x_1}^2\) and \(\delta_{x_2}^2\) commute,
\[
  \bigl(1-\tfrac r2\delta_{x_1}^2\bigr)\bigl(1-\tfrac r2\delta_{x_2}^2\bigr)U^{n+1}
  =\bigl(1+\tfrac r2\delta_{x_1}^2\bigr)\bigl(1-\tfrac r2\delta_{x_1}^2\bigr)U^{n+1/2}
  =\bigl(1+\tfrac r2\delta_{x_1}^2\bigr)\bigl(1+\tfrac r2\delta_{x_2}^2\bigr)U^n .
\]
(The commutation holds for constant coefficients on a rectangle.)

The factored scheme couples all grid points at the new level through a
large banded system.  The split form needs only tridiagonal solves: in
the first half-step one tridiagonal system along each line in the
\(x_1\)-direction, in the second along each line in the
\(x_2\)-direction.  Each costs \(O(N)\) by the Thomas algorithm, so a
full step costs \(O(N^2)\) for \(N^2\) unknowns.
\end{solution}

%---------------------------------------------------------------------
\Q{Consistency of the Peaceman--Rachford half-steps}{Analysis}{5}\label{q:prhalf}
\begin{qbox}
(Mitchell--Griffiths, Ch.~2, Exercise 14, first part.)  Show that if
\(U^{n+1*}=U^{n+1/2}\), the individual formulas of the
Peaceman--Rachford scheme
\[
  \bigl(1-\tfrac12r\delta_{x_1}^2\bigr)U^{n+1*}=\bigl(1+\tfrac12r\delta_{x_2}^2\bigr)U^n,
  \qquad
  \bigl(1-\tfrac12r\delta_{x_2}^2\bigr)U^{n+1}=\bigl(1+\tfrac12r\delta_{x_1}^2\bigr)U^{n+1*}
\]
are each consistent with \(u_t=u_{x_1x_1}+u_{x_2x_2}\).
\end{qbox}

\begin{rem}
\begin{itemize}
  \item Each half-step advances by \(\frac k2\); consistency means
    \(T/(\frac k2)\to0\).
  \item \(\frac r2\delta_{x_i}^2v=\frac k2v_{x_ix_i}+\frac{kh^2}{24}v_{x_i^4}+\cdots\).
\end{itemize}
\end{rem}

\begin{solution}
\step{First half-step}
Substitute the exact solution, with \(u^{n+1/2}\) for \(U^{n+1*}\), and
expand about \((lh,mh,nk)\):
\[
  T_1=u^{n+\frac12}-u^n-\tfrac r2\delta_{x_1}^2u^{n+\frac12}-\tfrac r2\delta_{x_2}^2u^n .
\]
Here \(u^{n+\frac12}-u^n=\frac k2u_t+\frac{k^2}8u_{tt}+\cdots\),
\(\frac r2\delta_{x_1}^2u^{n+\frac12}=\frac k2u_{x_1x_1}+\frac{k^2}4u_{x_1x_1t}+\frac{kh^2}{24}u_{x_1^4}+\cdots\)
and \(\frac r2\delta_{x_2}^2u^n=\frac k2u_{x_2x_2}+\frac{kh^2}{24}u_{x_2^4}+\cdots\).
So
\[
  T_1=\frac k2\bigl(u_t-u_{x_1x_1}-u_{x_2x_2}\bigr)+O(k^2+kh^2)=O(k^2+kh^2),
\]
and \(T_1/(\frac k2)=O(k+h^2)\to0\): consistent.

\step{Second half-step}
\(T_2=u^{n+1}-u^{n+\frac12}-\frac r2\delta_{x_2}^2u^{n+1}-\frac r2\delta_{x_1}^2u^{n+\frac12}\).
Expand about \((lh,mh,(n+\frac12)k)\), writing \(u\) for the value
there: \(u^{n+1}-u^{n+\frac12}=\frac k2u_t+\frac{k^2}8u_{tt}+\cdots\),
\(\frac r2\delta_{x_2}^2u^{n+1}=\frac k2u_{x_2x_2}+\frac{k^2}4u_{x_2x_2t}+\frac{kh^2}{24}u_{x_2^4}+\cdots\),
and \(\frac r2\delta_{x_1}^2u^{n+\frac12}=\frac k2u_{x_1x_1}+\frac{kh^2}{24}u_{x_1^4}+\cdots\).
Hence
\[
  T_2=\frac k2\bigl(u_t-u_{x_1x_1}-u_{x_2x_2}\bigr)+O(k^2+kh^2)=O(k^2+kh^2),
\]
and \(T_2/(\frac k2)=O(k+h^2)\to0\): consistent.

Each half-step is a first-order (in time) approximation; combined, the
errors cancel to give the second-order Crank--Nicolson accuracy.
(Compare the high-order Douglas split in Q\ref{q:mfcons}, whose
half-steps are \emph{not} consistent.)
\end{solution}

%---------------------------------------------------------------------
\Q{Stability of the Peaceman--Rachford method}{Stability analysis}{8}\label{q:old29}
\begin{qbox}
Find the amplification factor of one full Peaceman--Rachford step and
show that the method is unconditionally stable.  Is each half-step
stable by itself?
\end{qbox}
\begin{rem}
\begin{itemize}
  \item Fourier mode \(\xi^ne^{i(\beta lh+\gamma mh)}\), with
    \(\delta_{x_1}^2\to-4s_1^2\), \(\delta_{x_2}^2\to-4s_2^2\).
  \item \(\bigl|\frac{1-a}{1+a}\bigr|\le1\) for \(a\ge0\).
\end{itemize}
\end{rem}
\begin{solution}
Let \(U^n=\xi^ne^{i(\beta lh+\gamma mh)}\), \(s_1=\sin\frac{\beta h}2\),
\(s_2=\sin\frac{\gamma h}2\), and \(\xi^*\) the half-step factor.  Then
\[
  (1+2rs_1^2)\,\xi^*=1-2rs_2^2,
  \qquad
  (1+2rs_2^2)\,\xi=(1-2rs_1^2)\,\xi^*,
\]
so
\[
  \xi=\frac{1-2rs_1^2}{1+2rs_1^2}\cdot\frac{1-2rs_2^2}{1+2rs_2^2}.
\]
Each factor has the Crank--Nicolson form \(\frac{1-a}{1+a}\) with
\(a\ge0\), so \(|\xi|\le1\) for all \(r\), \(\beta\), \(\gamma\): the
method is unconditionally stable.

A half-step alone is not: \(\xi^*=\frac{1-2rs_2^2}{1+2rs_1^2}\) equals
\(1-2r\) when \(s_1=0\), \(s_2=1\), which has modulus greater than \(1\)
for \(r>1\).  Only the pair of half-steps is stable.
\end{solution}

%---------------------------------------------------------------------
\Q{One Peaceman--Rachford step}{Numerical}{6}\label{q:prnum}
\begin{qbox}
Solve \(u_t=u_{x_1x_1}+u_{x_2x_2}\) on the unit square, with
\(u(x_1,x_2,0)=\sin\pi x_1\sin\pi x_2\) and \(u=0\) on the boundary, by
one Peaceman--Rachford step with \(h=\frac13\) and \(r=1\).  Compare
with the exact solution \(e^{-2\pi^2t}\sin\pi x_1\sin\pi x_2\).
\end{qbox}

\begin{rem}
\begin{itemize}
  \item \(\bigl(1-\frac r2\delta_{x_1}^2\bigr)U^{*}=\bigl(1+\frac r2\delta_{x_2}^2\bigr)U^n\),
    \(\bigl(1-\frac r2\delta_{x_2}^2\bigr)U^{n+1}=\bigl(1+\frac r2\delta_{x_1}^2\bigr)U^{*}\).
  \item Written out: \((1+r)U_{l,m}^*-\frac r2(U_{l+1,m}^*+U_{l-1,m}^*)
    =(1-r)U_{l,m}^n+\frac r2(U_{l,m+1}^n+U_{l,m-1}^n)\); a tridiagonal
    system along each line \(x_2=\)const.
  \item Amplification factor
    \(\xi=\frac{1-2rs_1^2}{1+2rs_1^2}\cdot\frac{1-2rs_2^2}{1+2rs_2^2}\).
\end{itemize}
\end{rem}

\begin{solution}
\(k=rh^2=\frac19\).  The interior points are \((l,m)\), \(l,m=1,2\), all
with initial value \(\sin^2\frac\pi3=0.75\).  By symmetry all four
values stay equal at every stage; call them \(U^*\) and \(U^1\).

\step{First half-step (implicit in \(x_1\))}
Along the line \(m=1\), with \(r=1\):
\(2U_{l,1}^*-\frac12(U_{l+1,1}^*+U_{l-1,1}^*)=0\cdot U_{l,1}^0+\frac12(U_{l,2}^0+U_{l,0}^0)\).
For \(l=1\): \(U_{0,1}^*=0\), \(U_{2,1}^*=U^*\), \(U_{1,2}^0=0.75\),
\(U_{1,0}^0=0\):
\[
  2U^*-\tfrac12U^*=\tfrac12(0.75)\ \Longrightarrow\ 1.5\,U^*=0.375,
  \qquad U^*=0.25 .
\]
(The \(2\times2\) system on each line is
\(\begin{pmatrix}2&-\frac12\\-\frac12&2\end{pmatrix}\begin{pmatrix}U_{1}^*\\U_{2}^*\end{pmatrix}
=\begin{pmatrix}0.375\\0.375\end{pmatrix}\).)

\step{Second half-step (implicit in \(x_2\))}
Along each line \(l=\)const:
\(2U_{l,m}^1-\frac12(U_{l,m+1}^1+U_{l,m-1}^1)=\frac12(U_{l+1,m}^*+U_{l-1,m}^*)\),
so \(1.5\,U^1=\frac12(0.25)=0.125\):
\[
  U^1=0.083333 .
\]

\step{Check and comparison}
With \(s_1^2=s_2^2=\sin^2\frac\pi6=\frac14\),
\(\xi=\bigl(\frac{1-\frac12}{1+\frac12}\bigr)^2=\frac19\), and
\(0.75\times\frac19=0.083333\).  The exact value at \(t=\frac19\) is
\(0.75\,e^{-2\pi^2/9}=0.083666\): an error of \(3.3\times10^{-4}\), although
\(r=1\) is four times the explicit limit \(r\le\frac14\).
\end{solution}

%---------------------------------------------------------------------
\Q{Richardson's scheme is unstable: a determinant proof}{Proof}{6}\label{q:old30}
\begin{qbox}
The amplification matrix of the three-level scheme
\(U^{n+1}=2r\delta^2U^n+U^{n-1}\) is
\(G=\begin{pmatrix}-8rs^2&1\\1&0\end{pmatrix}\).  Without solving for the
eigenvalues, prove that \(G\) has an eigenvalue of modulus greater than
\(1\) whenever \(s\ne0\).
\end{qbox}
\begin{rem}
\begin{itemize}
  \item Richardson: \(U_m^{n+1}=U_m^{n-1}+2r\delta^2U_m^n\); write as a
    two-level system for \(W=(U,V)^T\), \(V^n=U^{n-1}\).
  \item Amplification matrix
    \(G=\begin{pmatrix}-8rs^2&1\\1&0\end{pmatrix}\); von Neumann:
    \(\max|\lambda_i(G)|\le1\).
  \item \(\lambda_1\lambda_2=\det G\), \(\lambda_1+\lambda_2=\mathrm{tr}\,G\).
\end{itemize}
\end{rem}
\begin{solution}
\(\det G=\lambda_1\lambda_2=-1\) and
\(\mathrm{trace}\,G=\lambda_1+\lambda_2=-8rs^2\).  The discriminant of
\(\lambda^2+8rs^2\lambda-1\) is \(64r^2s^4+4>0\), so the eigenvalues are
real and distinct.  Suppose both had modulus \(\le1\).  Since
\(|\lambda_1||\lambda_2|=1\), both have modulus exactly \(1\); being
real, they are \(\pm1\), and with product \(-1\) they are \(1\) and
\(-1\).  Then the trace is \(0\), contradicting
\(-8rs^2\ne0\).  Hence \(\max|\lambda_i|>1\) for every \(r>0\) and every
mode with \(s\ne0\): the scheme is unconditionally unstable.
\end{solution}

%---------------------------------------------------------------------
\Q{Stability of the Du Fort--Frankel scheme}{Stability analysis}{10}\label{q:old31}
\begin{qbox}
Prove that the Du Fort--Frankel scheme
\((1+2r)U_m^{n+1}=2r(U_{m+1}^n+U_{m-1}^n)+(1-2r)U_m^{n-1}\) is
unconditionally stable.
\end{qbox}
\begin{rem}
\begin{itemize}
  \item Du Fort--Frankel:
    \((1+2r)U_m^{n+1}=2r(U_{m+1}^n+U_{m-1}^n)+(1-2r)U_m^{n-1}\);
    replace \(U_m^n\) in Richardson by \(\frac12(U_m^{n+1}+U_m^{n-1})\).
  \item Characteristic equation
    \((1+2r)\xi^2-4r\cos\theta\,\xi-(1-2r)=0\); product of roots
    \(-\frac{1-2r}{1+2r}\).
  \item The matrix-method proof is in Q\ref{q:dffmat}.
\end{itemize}
\end{rem}
\begin{solution}
Substitute \(U_m^n=\xi^ne^{i\beta mh}\), \(\theta=\beta h\):
\[
  (1+2r)\xi^2-4r\cos\theta\,\xi-(1-2r)=0,
\]
whose roots are the eigenvalues of the \(2\times2\) amplification
matrix of the equivalent two-level system.  They are
\[
  \xi_{1,2}=\frac{2r\cos\theta\pm\sqrt{1-4r^2\sin^2\theta}}{1+2r}.
\]
\step{Case 1: \(4r^2\sin^2\theta\le1\) (real roots)}
\[
  |\xi_{1,2}|\le\frac{2r|\cos\theta|+\sqrt{1-4r^2\sin^2\theta}}{1+2r}
  \le\frac{2r+1}{1+2r}=1 .
\]
\step{Case 2: \(4r^2\sin^2\theta>1\) (complex conjugate roots)}
Then \(r>\frac12\), and
\(|\xi|^2=|\xi_1\xi_2|=\bigl|\frac{-(1-2r)}{1+2r}\bigr|=\frac{2r-1}{2r+1}<1\).

So \(|\xi_i|\le1\) for all \(r>0\) and all \(\theta\).  Roots of
modulus one occur only for \(\theta=0\) or \(\pi\), where they are
simple, so the powers of the amplification matrix stay bounded: the
scheme is unconditionally stable.
\end{solution}

%---------------------------------------------------------------------
\Q{Consistency of the Du Fort--Frankel scheme}{Derivation}{8}\label{q:old32}
\begin{qbox}
Show that the truncation error of the Du Fort--Frankel scheme (written
as an approximation to \(u_t=u_{xx}\)) is
\[
  T=\Bigl(\frac kh\Bigr)^2u_{tt}+\frac{k^2}6u_{ttt}-\frac{h^2}{12}u_{xxxx}+\cdots .
\]
When is the scheme consistent?  What happens if \(k/h=c\) is held fixed?
\end{qbox}
\begin{rem}
\begin{itemize}
  \item \(u^{n+1}-u^{n-1}=2ku_t+\frac{k^3}3u_{ttt}+\cdots\),
    \(u^{n+1}+u^{n-1}=2u+k^2u_{tt}+\cdots\).
  \item Consistent with \(u_t=u_{xx}\) only if \(k/h\to0\); with
    \(k/h=c\) fixed, consistent with \(u_t+c^2u_{tt}=u_{xx}\).
\end{itemize}
\end{rem}
\begin{solution}
Write the scheme as
\[
  \frac{U_m^{n+1}-U_m^{n-1}}{2k}
  =\frac{U_{m+1}^n-U_m^{n+1}-U_m^{n-1}+U_{m-1}^n}{h^2}.
\]
Using
\(u^{n+1}-u^{n-1}=2ku_t+\frac{k^3}3u_{ttt}+\cdots\),
\(u^{n+1}+u^{n-1}=2u+k^2u_{tt}+\cdots\) and
\(u_{m+1}+u_{m-1}=2u+h^2u_{xx}+\frac{h^4}{12}u_{xxxx}+\cdots\),
\[
  T=u_t+\frac{k^2}6u_{ttt}-\Bigl[u_{xx}+\frac{h^2}{12}u_{xxxx}-\frac{k^2}{h^2}u_{tt}\Bigr]+\cdots
   =\Bigl(\frac kh\Bigr)^2u_{tt}+\frac{k^2}6u_{ttt}-\frac{h^2}{12}u_{xxxx}+\cdots,
\]
since \(u_t=u_{xx}\).  So \(T\to0\) if and only if \(k/h\to0\).  With
\(r=k/h^2\) fixed, \(k/h=rh\to0\), so the scheme is consistent.  But if
\(k/h=c\) is fixed, \(T\to c^2u_{tt}\ne0\): the scheme is then
consistent with the hyperbolic equation \(u_t+c^2u_{tt}=u_{xx}\), not
with the heat equation.  Unconditional stability does not mean that any
\(k\) gives a good answer.
\end{solution}

%---------------------------------------------------------------------
\Q{Du Fort--Frankel with a large time step}{Numerical}{6}\label{q:old33}
\begin{qbox}
For \(u_t=u_{xx}\), \(u(x,0)=\sin\pi x\), \(u(0,t)=u(1,t)=0\), take
\(h=\frac14\), \(r=1\), and the exact value at \(t=k\) as the second
starting level.  Compute \(U\) at \(x=\frac12\), \(t=2k\) by Du
Fort--Frankel and compare with the exact value.  Explain the size of
the error.
\end{qbox}
\begin{rem}
\begin{itemize}
  \item A three-level scheme needs \(U^0\) and \(U^1\).
  \item For \(\sin\pi x\) data, \(U_{m+1}+U_{m-1}=2\cos(\pi h)U_m\).
  \item Unconditional stability does not mean accuracy for every \(k\).
\end{itemize}
\end{rem}
\begin{solution}
\(k=\frac1{16}\).  The data stay proportional to \(\sin\pi x\):
\(U_m^n=a_n\sin\pi x_m\), with \(U_{m+1}+U_{m-1}=2\cos\frac\pi4\,U_m\).
The scheme gives
\[
  3a_{n+1}=2\cdot2\cos\tfrac\pi4\,a_n-a_{n-1}=2.828427\,a_n-a_{n-1}.
\]
With \(a_0=1\), \(a_1=e^{-\pi^2/16}=0.539641\):
\[
  a_2=\frac{2.828427\times0.539641-1}{3}=0.175445,
\]
so \(U(\tfrac12,2k)=0.175445\), against the exact
\(e^{-2\pi^2/16}=0.291213\): an error of about \(40\%\).

The reason is the consistency error \((k/h)^2u_{tt}\).  Here
\(k/h=\frac14\), and for this mode \(u_{tt}=\pi^4u\approx97.4u\), so the
extra term is about \(6.1u\), comparable with \(u_t=-\pi^2u\approx-9.9u\).
The scheme is stable, but at this step size it is not accurate.
\end{solution}

%---------------------------------------------------------------------
\Q{The three-level formula of maximum accuracy}{Derivation / stability analysis}{10}\label{q:maxacc}
\begin{qbox}
(Mitchell--Griffiths, Ch.~2, Exercise 30.)  Show that the formula of
maximum accuracy for \(u_t=u_{xx}\) involving the grid points
\((mh,(n+1)k)\), \((mh,nk)\), \((mh,(n-1)k)\), \(((m\pm1)h,nk)\) (the
points of the Du Fort--Frankel scheme) is
\[
  (1+6r)U_m^{n+1}=2(1-12r^2)U_m^n+12r^2\bigl(U_{m+1}^n+U_{m-1}^n\bigr)-(1-6r)U_m^{n-1}.
\]
Find its truncation error and show that it is stable provided
\(r\le\frac1{2\sqrt3}\).
\end{qbox}

\begin{rem}
\begin{itemize}
  \item Undetermined coefficients: expand every point about
    \((mh,nk)\) and use \(u_t=u_{xx}\), \(u_{tt}=u_{xxxx}\),
    \(u_{ttt}=u_{xxxxxx}\) to remove time derivatives.
  \item For a real quadratic \(z^2+a_1z+a_0\): both roots in
    \(|z|\le1\) iff \(|a_0|\le1\) and \(|a_1|\le1+a_0\).
\end{itemize}
\end{rem}

\begin{solution}
\step{Derivation}
Seek \(\alpha_1U_m^{n+1}+\alpha_0U_m^n+\beta(U_{m+1}^n+U_{m-1}^n)+\alpha_{-1}U_m^{n-1}=0\).
Expanding,
\begin{align*}
  u:&\quad \alpha_1+\alpha_0+2\beta+\alpha_{-1},\\
  u_t=u_{xx}:&\quad (\alpha_1-\alpha_{-1})k+\beta h^2,\\
  u_{tt}=u_{xxxx}:&\quad \tfrac12(\alpha_1+\alpha_{-1})k^2+\tfrac1{12}\beta h^4,\\
  u_{ttt}=u_{xxxxxx}:&\quad \tfrac16(\alpha_1-\alpha_{-1})k^3+\tfrac1{360}\beta h^6 .
\end{align*}
Four coefficients (one fixed by scaling) allow three conditions: set
the first three coefficients to zero.  Take \(\beta=-12r^2\).  Then
\(\alpha_1-\alpha_{-1}=-\beta h^2/k=12r\) and
\(\alpha_1+\alpha_{-1}=-\beta h^4/(6k^2)=2\), so
\(\alpha_1=1+6r\), \(\alpha_{-1}=1-6r\), and
\(\alpha_0=-(\alpha_1+\alpha_{-1}+2\beta)=-2(1-12r^2)\).  This is the
stated formula.

\step{Truncation error}
The first non-zero term is the \(u_{xxxxxx}\) one:
\[
  T=\Bigl(\tfrac16\cdot12r\,k^3-\tfrac{12r^2}{360}h^6\Bigr)\frac{\partial^6u}{\partial x^6}+\cdots
  =\Bigl(2r^2-\frac1{30}\Bigr)k^2h^2\frac{\partial^6u}{\partial x^6}+\cdots,
\]
using \(k=rh^2\).  The formula is \(12rk\,(u_t-u_{xx})+T\), so the
normalised error is \(T/(12rk)=O(h^4)\) for fixed \(r\) (and it is of
even higher order when \(r^2=\frac1{60}\)).

\step{Stability}
With \(U_m^n=\xi^ne^{i\beta mh}\) and \(S=\sin^2\frac{\beta h}2\),
\(U_{m+1}+U_{m-1}=(2-4S)U_m\), so
\[
  (1+6r)\xi^2-2\bigl(1-24r^2S\bigr)\xi+(1-6r)=0 .
\]
Divide by \(1+6r\): \(a_0=\frac{1-6r}{1+6r}\), \(|a_0|<1\) for \(r>0\);
and \(|a_1|\le1+a_0\) reads
\[
  \frac{2|1-24r^2S|}{1+6r}\le\frac2{1+6r}\iff|1-24r^2S|\le1\iff24r^2S\le2 .
\]
This must hold for all \(S\in[0,1]\): \(r^2\le\frac1{12}\), i.e.\
\[
  \boxed{r\le\frac1{2\sqrt3}}.
\]
(Equality \(|a_1|=1+a_0\) gives a root \(-1\), simple, which is
allowed.)
\end{solution}

%---------------------------------------------------------------------
\Q{Short answers on Lecture 35}{Short answer}{6}\label{q:old34}
\begin{qbox}
\begin{enumerate}[label=(\alph*)]
  \item Du Fort--Frankel contains \(U_m^{n+1}\) on both sides.  Is it
    explicit or implicit?
  \item What extra information does a three-level scheme need to start?
  \item Richardson's scheme has truncation error \(O(k^2+h^2)\), better
    than the explicit scheme.  Why is it useless?
  \item In ADI, what does ``alternating direction implicit'' refer to?
\end{enumerate}
\end{qbox}
\begin{rem}
\begin{itemize}
  \item Explicit: the only unknown at the new level in each equation.
  \item Lax equivalence: consistent \(+\) stable \(\iff\) convergent.
\end{itemize}
\end{rem}
\begin{solution}
\step{(a)}
Explicit: \(U_m^{n+1}\) is the only unknown at the new level in each
equation, so it is obtained directly by dividing by \(1+2r\); no system
has to be solved.

\step{(b)}
Two initial levels, \(U^0\) and \(U^1\).  \(U^1\) must be computed by a
two-level method (for example one explicit or Crank--Nicolson step, or
a Taylor expansion) of adequate accuracy.

\step{(c)}
It is unconditionally unstable (Q\ref{q:old30}): for every \(r>0\) some mode grows
like \(|\lambda_2|^n\) with \(|\lambda_2|>1\).  By the Lax equivalence
theorem it does not converge, whatever its truncation error.

\step{(d)}
Each half-step is implicit in one space direction and explicit in the
other, and the implicit direction alternates between \(x_1\) and
\(x_2\) from one half-step to the next.
\end{solution}

%#####################################################################
\Needspace{32\baselineskip}
\section{Block Matrices and the Du Fort--Frankel Scheme by the Matrix Method}
%#####################################################################

%---------------------------------------------------------------------
\Q{Eigenvalues of block matrices}{Proof / numerical}{8}\label{q:block}
\begin{qbox}
\begin{enumerate}[label=(\alph*)]
  \item Let \(A=\begin{bmatrix}A_{11}&A_{12}\\A_{21}&A_{22}\end{bmatrix}\),
    where the \(n\times n\) blocks \(A_{ij}\) have a common set of
    \(n\) linearly independent eigenvectors \(\mathbf v_k\), with
    \(A_{ij}\mathbf v_k=\lambda_{ij}^{(k)}\mathbf v_k\).  Prove that the
    eigenvalues of \(A\) are the eigenvalues of the \(2\times2\)
    matrices \(\bigl[\lambda_{ij}^{(k)}\bigr]\), \(k=1,\dots,n\).
  \item Use this to find the eigenvalues of
    \(\begin{bmatrix}T&I\\I&T\end{bmatrix}\), where
    \(T=\begin{bmatrix}0&1\\1&0\end{bmatrix}\).
\end{enumerate}
\end{qbox}

\begin{rem}
\begin{itemize}
  \item Theorem (Lecture 36): if all blocks \(A_{ij}\) (\(m\times m\)
    blocks of size \(n\)) share \(n\) independent eigenvectors, the
    eigenvalues of \(A\) are those of the \(m\times m\) matrices
    \(\Lambda^{(k)}=\bigl[\lambda_{ij}^{(k)}\bigr]\): \(mn\) in total.
  \item Look for eigenvectors of the form
    \((\alpha_1\mathbf v,\alpha_2\mathbf v)^T\).
  \item \(I\) has every vector as eigenvector (eigenvalue \(1\)); the
    zero block has eigenvalue \(0\).
\end{itemize}
\end{rem}

\begin{solution}
\step{(a) Proof}
Fix \(k\) and write \(\mathbf v=\mathbf v_k\),
\(\lambda_{ij}=\lambda_{ij}^{(k)}\).  For constants \(\alpha_1,\alpha_2\),
\[
  \begin{bmatrix}A_{11}&A_{12}\\A_{21}&A_{22}\end{bmatrix}
  \begin{bmatrix}\alpha_1\mathbf v\\\alpha_2\mathbf v\end{bmatrix}
  =\begin{bmatrix}(\lambda_{11}\alpha_1+\lambda_{12}\alpha_2)\mathbf v\\
  (\lambda_{21}\alpha_1+\lambda_{22}\alpha_2)\mathbf v\end{bmatrix}.
\]
This equals \(\mu\,(\alpha_1\mathbf v,\alpha_2\mathbf v)^T\) if and only if
\[
  (\lambda_{11}-\mu)\alpha_1+\lambda_{12}\alpha_2=0,\qquad
  \lambda_{21}\alpha_1+(\lambda_{22}-\mu)\alpha_2=0,
\]
which has a non-trivial solution \((\alpha_1,\alpha_2)\) if and only if
\[
  \det\begin{bmatrix}\lambda_{11}-\mu&\lambda_{12}\\\lambda_{21}&\lambda_{22}-\mu\end{bmatrix}=0,
\]
i.e.\ \(\mu\) is an eigenvalue of \(\bigl[\lambda_{ij}\bigr]\).  Each
\(k\) gives \(2\) eigenvalues, with eigenvectors
\((\alpha_1\mathbf v_k,\alpha_2\mathbf v_k)^T\); these \(2n\) vectors are
linearly independent (the \(\mathbf v_k\) are), so all \(2n\)
eigenvalues of \(A\) are obtained.

\step{(b) Example}
\(T\) and \(I\) share the eigenvectors \(\mathbf v_1=(1,1)^T\) (\(T\):
\(1\)) and \(\mathbf v_2=(1,-1)^T\) (\(T\): \(-1\)); \(I\) has eigenvalue
\(1\) for both.
\[
  k=1:\ \begin{bmatrix}1&1\\1&1\end{bmatrix}\to\mu=2,\,0;
  \qquad
  k=2:\ \begin{bmatrix}-1&1\\1&-1\end{bmatrix}\to\mu=0,\,-2 .
\]
So the eigenvalues are \(2,0,0,-2\).  (Check: the \(4\times4\) matrix has
trace \(0\) and rank \(2\), and \((1,1,1,1)^T\) gives eigenvalue \(2\).)
\end{solution}

%---------------------------------------------------------------------
\Q{Du Fort--Frankel: stability by the matrix method}{Stability analysis}{10}\label{q:dffmat}
\begin{qbox}
Investigate by the matrix method the stability of the Du Fort--Frankel
scheme
\[
  U_m^{n+1}=\frac{2r}{1+2r}\bigl(U_{m+1}^n+U_{m-1}^n\bigr)+\frac{1-2r}{1+2r}U_m^{n-1},
  \qquad r=\frac k{h^2},
\]
for \(u_t=u_{xx}\), \(0<x<1\), \(t>0\), for fixed mesh length (\(Nh=1\)),
with \(u\) known on \(x=0\) and \(x=1\).
\end{qbox}

\begin{rem}
\begin{itemize}
  \item Reduce three levels to two: \(V^n=U^{n-1}\), \(W=(U,V)^T\),
    \(W^{n+1}=BW^n+g^n\); stable if every eigenvalue of \(B\) has
    modulus \(\le1\).
  \item \(\mathrm{tridiag}(1,0,1)\) of order \(N-1\) has the distinct
    eigenvalues \(2\cos\frac{k\pi}N\).
  \item Block theorem (Q\ref{q:block}).
\end{itemize}
\end{rem}

\begin{solution}
\step{Matrix form}
With \(U^n=(U_1^n,\dots,U_{N-1}^n)^T\),
\[
  U^{n+1}=\frac{2r}{1+2r}AU^n+\frac{1-2r}{1+2r}U^{n-1}+f_n,
  \qquad A=\mathrm{tridiag}(1,0,1)_{(N-1)\times(N-1)},
\]
\(f_n=\frac{2r}{1+2r}(U_0^n,0,\dots,0,U_N^n)^T\).

\step{Two-level form}
With \(V^n=U^{n-1}\) and \(W=(U,V)^T\):
\[
  W^{n+1}=BW^n+g^n,\qquad
  B=\begin{bmatrix}\dfrac{2r}{1+2r}A&\dfrac{1-2r}{1+2r}I\\[6pt] I&0\end{bmatrix},
  \qquad g^n=\begin{bmatrix}f_n\\0\end{bmatrix}.
\]
The scheme is stable if every eigenvalue of \(B\) has modulus \(\le1\).

\step{Eigenvalues of \(B\)}
\(A\) has the \(N-1\) distinct eigenvalues
\(\lambda_k=2\cos\frac{k\pi}N\), so it has \(N-1\) independent
eigenvectors \(\mathbf v_k\), which are also eigenvectors of \(I\)
(eigenvalue \(1\)) and \(0\) (eigenvalue \(0\)).  By the block theorem
the eigenvalues \(\lambda\) of \(B\) are those of
\(\begin{bmatrix}\frac{2r}{1+2r}\lambda_k&\frac{1-2r}{1+2r}\\1&0\end{bmatrix}\):
\[
  \lambda^2-\frac{2r\lambda_k}{1+2r}\lambda-\frac{1-2r}{1+2r}=0
  \ \Longrightarrow\
  \lambda=\frac{r\lambda_k\pm\bigl(1+(\lambda_k^2-4)r^2\bigr)^{1/2}}{1+2r},
\]
that is,
\[
  \lambda=\frac{2r\cos\frac{k\pi}N\pm\bigl(1-4r^2\sin^2\frac{k\pi}N\bigr)^{1/2}}{1+2r}.
\]

\step{Case I: \(0\le1-4r^2\sin^2\frac{k\pi}N<1\) (real roots)}
Since \(|\cos\frac{k\pi}N|<1\),
\(|\lambda|<\frac{2r+1}{1+2r}=1\).

\step{Case II: \(1-4r^2\sin^2\frac{k\pi}N<0\) (complex roots)}
\[
  |\lambda|=\frac{\bigl(4r^2\cos^2\frac{k\pi}N+4r^2\sin^2\frac{k\pi}N-1\bigr)^{1/2}}{1+2r}
  =\frac{(4r^2-1)^{1/2}}{1+2r}=\Bigl(\frac{2r-1}{2r+1}\Bigr)^{1/2}<1 .
\]
Hence all eigenvalues of \(B\) have modulus less than \(1\) for every
\(r>0\): for a fixed mesh size the scheme is unconditionally stable.
(It is consistent only if \(k/h\to0\), Q\ref{q:old32}.)
\end{solution}

%#####################################################################
\Needspace{32\baselineskip}
\section{Elliptic Equations and the Five-Point Formula}
%#####################################################################

%---------------------------------------------------------------------
\Q{Boundary value problems and the five-point formula}{Derivation}{8}\label{q:fivept}
\begin{qbox}
\begin{enumerate}[label=(\alph*)]
  \item State the Dirichlet, Neumann and Robin problems for an
    elliptic equation in a bounded region \(R\).
  \item Derive the five-point replacement of Laplace's equation on a
    square grid of spacing \(h\), and show that its local truncation
    error is \(\frac{h^2}{12}\bigl(u_{x_1^4}+u_{x_2^4}\bigr)+\cdots\).
\end{enumerate}
\end{qbox}

\begin{rem}
\begin{itemize}
  \item \(au_{x_1x_1}+2bu_{x_1x_2}+cu_{x_2x_2}=d\) is elliptic if
    \(b^2-ac<0\); Laplace \(\nabla^2u=0\), Poisson \(\nabla^2u=f\).
  \item \(u_{l+1,m}+u_{l-1,m}-2u_{l,m}=h^2u_{x_1x_1}+\frac{h^4}{12}u_{x_1^4}+\cdots\).
  \item Five-point: \(U_{l+1,m}+U_{l-1,m}+U_{l,m+1}+U_{l,m-1}-4U_{l,m}=h^2f_{l,m}\).
\end{itemize}
\end{rem}

\begin{solution}
\step{(a) The three problems}
\begin{itemize}
  \item \emph{Dirichlet (first BVP):} \(u\) is prescribed on
    \(\partial R\).
  \item \emph{Neumann (second BVP):} \(\partial u/\partial n=f\) on
    \(\partial R\), \(n\) the outward normal.
  \item \emph{Robin (third BVP):} \(\alpha u+\beta\,\partial u/\partial
    n=\gamma\) on \(\partial R\), with \(\alpha,\beta>0\).
\end{itemize}

\step{(b) Derivation}
By Taylor's theorem about \((lh,mh)\),
\[
  u_{l\pm1,m}=u\pm hu_{x_1}+\frac{h^2}2u_{x_1x_1}\pm\frac{h^3}6u_{x_1^3}+\frac{h^4}{24}u_{x_1^4}\pm\cdots,
\]
and similarly for \(u_{l,m\pm1}\).  Adding all four, the odd terms
cancel:
\[
  u_{l+1,m}+u_{l-1,m}+u_{l,m+1}+u_{l,m-1}-4u_{l,m}
  =h^2\nabla^2u+\frac{h^4}{12}\bigl(u_{x_1^4}+u_{x_2^4}\bigr)+\cdots .
\]
Dividing by \(h^2\),
\(\nabla^2u\approx\frac1{h^2}\bigl(u_{l+1,m}+u_{l-1,m}+u_{l,m+1}+u_{l,m-1}-4u_{l,m}\bigr)\),
with local truncation error
\(\frac{h^2}{12}\bigl(u_{x_1^4}+u_{x_2^4}\bigr)+\cdots=O(h^2)\).  The
five-point replacement of Laplace's equation is
\[
  U_{l+1,m}+U_{l-1,m}+U_{l,m+1}+U_{l,m-1}-4U_{l,m}=0 .
\]
\end{solution}

%---------------------------------------------------------------------
\Q{Difference coefficients from the Taylor-series system}{Derivation}{6}\label{q:taylorcoef}
\begin{qbox}
(Mitchell--Griffiths, Ch.~3, Example 3 and Exercise 4.)  Laplace's
equation is replaced at \(P\) by
\(L_hU=\alpha_0U_P-\sum_{i=1}^4\alpha_iU_{Q_i}=0\), where \(L_h\)
approximates \(-\nabla^2\).
\begin{enumerate}[label=(\alph*)]
  \item For the neighbours \((h,0),(0,h),(-h,0),(0,-h)\) relative to
    \(P\), show that the Taylor-series conditions give
    \(\alpha_0=\frac4{h^2}\), \(\alpha_1=\dots=\alpha_4=\frac1{h^2}\).
  \item Find \(\alpha_0,\dots,\alpha_4\) for the neighbours
    \((2h,0),(0,h),(-2h,0),(0,-h)\).
\end{enumerate}
\end{qbox}

\begin{rem}
\begin{itemize}
  \item For \(Q_i=P+(a_i,b_i)\):
    \(u_{Q_i}=u+a_iu_{x_1}+b_iu_{x_2}+\frac12(a_i^2u_{x_1x_1}+2a_ib_iu_{x_1x_2}+b_i^2u_{x_2x_2})+\cdots\).
  \item Require \(L_hu=-\nabla^2u+\)(higher order): match the
    coefficients of \(u,u_{x_1},u_{x_2},u_{x_1x_1},u_{x_1x_2},u_{x_2x_2}\).
\end{itemize}
\end{rem}

\begin{solution}
\step{Conditions}
\begin{align*}
  L_hu&=\Bigl(\alpha_0-\sum\alpha_i\Bigr)u-\Bigl(\sum\alpha_ia_i\Bigr)u_{x_1}-\Bigl(\sum\alpha_ib_i\Bigr)u_{x_2}\\
  &\quad-\tfrac12\Bigl(\sum\alpha_ia_i^2\Bigr)u_{x_1x_1}-\Bigl(\sum\alpha_ia_ib_i\Bigr)u_{x_1x_2}
  -\tfrac12\Bigl(\sum\alpha_ib_i^2\Bigr)u_{x_2x_2}-\cdots .
\end{align*}
For \(L_hu\approx-(u_{x_1x_1}+u_{x_2x_2})\):
\[
  \alpha_0=\sum\alpha_i,\quad \sum\alpha_ia_i=\sum\alpha_ib_i=0,\quad
  \sum\alpha_ia_ib_i=0,\quad \tfrac12\sum\alpha_ia_i^2=\tfrac12\sum\alpha_ib_i^2=1 .
\]

\step{(a) Square stencil}
\((a_i,b_i)=(h,0),(0,h),(-h,0),(0,-h)\).  First-order conditions:
\(h(\alpha_1-\alpha_3)=0\), \(h(\alpha_2-\alpha_4)=0\); the mixed one
holds automatically (\(a_ib_i=0\)).  Second order:
\(\frac12h^2(\alpha_1+\alpha_3)=1\Rightarrow\alpha_1=\alpha_3=\frac1{h^2}\),
and likewise \(\alpha_2=\alpha_4=\frac1{h^2}\).  Then
\(\alpha_0=\frac4{h^2}\): the five-point formula
\(\frac1{h^2}(4U_P-\sum U_{Q_i})=0\).  Four neighbours are just enough
for these five conditions (one is automatic).

\step{(b) Stretched stencil}
\((a_i,b_i)=(2h,0),(0,h),(-2h,0),(0,-h)\).  First order gives
\(\alpha_1=\alpha_3\), \(\alpha_2=\alpha_4\).  Second order:
\(\frac12(4h^2)(\alpha_1+\alpha_3)=1\Rightarrow\alpha_1=\alpha_3=\frac1{4h^2}\),
and \(\frac12h^2(\alpha_2+\alpha_4)=1\Rightarrow\alpha_2=\alpha_4=\frac1{h^2}\).
Hence
\[
  \alpha_0=\frac2{4h^2}+\frac2{h^2}=\frac5{2h^2},\qquad
  \alpha_1=\alpha_3=\frac1{4h^2},\qquad \alpha_2=\alpha_4=\frac1{h^2}.
\]
(This agrees with the rectangular-mesh formula of Q\ref{q:rect} with
\(h_1=2h\), \(h_2=h\).)
\end{solution}

%---------------------------------------------------------------------
\Q{The five-point matrix: structure and eigenvalues}{Proof}{8}\label{q:fivemat}
\begin{qbox}
For Laplace's equation on the unit square with Dirichlet conditions
and \(h=1/N\):
\begin{enumerate}[label=(\alph*)]
  \item show that the five-point equations form \(AU=b\), with \(A\) of
    order \((N-1)^2\), block tridiagonal, and with band width \(2N-1\)
    for the natural ordering (Mitchell--Griffiths, Ch.~3, Exercise 12);
  \item show that the eigenvalues of \(A\) are
    \(\lambda_{ij}=-4\bigl(\sin^2\frac{i\pi}{2N}+\sin^2\frac{j\pi}{2N}\bigr)\),
    \(1\le i,j\le N-1\), and hence \(|A|\ne0\).
\end{enumerate}
\end{qbox}

\begin{rem}
\begin{itemize}
  \item \(A=\mathrm{tridiag}(I,B,I)\), \(B=\mathrm{tridiag}(1,-4,1)\),
    blocks of order \(N-1\).
  \item \(\mathrm{tridiag}(1,0,1)\) of order \(N-1\): eigenvalues
    \(2\cos\frac{j\pi}N\), eigenvectors \((\sin\frac{lj\pi}N)_l\).
  \item \(2\cos\theta-2=-4\sin^2\frac\theta2\).
\end{itemize}
\end{rem}

\begin{solution}
\step{(a) Structure}
There are \((N-1)^2\) interior points, one unknown each, so \(A\) has
order \((N-1)^2\).  Order the unknowns line by line:
\(U=(U_{1,1},\dots,U_{1,N-1},U_{2,1},\dots,U_{N-1,N-1})^T\).  The
equation at \((l,m)\) couples \(U_{l,m\pm1}\) (same block, positions
\(\pm1\)) and \(U_{l\pm1,m}\) (neighbouring blocks, positions
\(\pm(N-1)\)).  So
\[
  A=\begin{bmatrix}B&I&&\\I&B&I&\\&\ddots&\ddots&\ddots\\&&I&B\end{bmatrix},
  \qquad B=\mathrm{tridiag}(1,-4,1),
\]
and \(b\) holds the boundary values.  The non-zero entries of each row
lie within \(N-1\) positions of the diagonal, so the band width is
\(2(N-1)+1=2N-1\).

\step{(b) Eigenvalues}
\(B=-4I+\mathrm{tridiag}(1,0,1)\) and \(I\) share the eigenvectors
\(\mathbf v_i\), with eigenvalues \(\beta_i=-4+2\cos\frac{i\pi}N\) and
\(1\).  By the block theorem (Q\ref{q:block}: if all the blocks share
a complete set of eigenvectors, the eigenvalues of \(A\) are those of
the matrices formed from the block eigenvalues), with \(N-1\) blocks
per row, the eigenvalues of \(A\) are those of
\(\mathrm{tridiag}(1,\beta_i,1)=\beta_iI+\mathrm{tridiag}(1,0,1)\),
namely \(\beta_i+2\cos\frac{j\pi}N\):
\[
  \lambda_{ij}=-4+2\cos\frac{i\pi}N+2\cos\frac{j\pi}N
  =-4\Bigl(\sin^2\frac{i\pi}{2N}+\sin^2\frac{j\pi}{2N}\Bigr).
\]
All \(\lambda_{ij}<0\), so none is zero: \(|A|\ne0\) and the discrete
Dirichlet problem has a unique solution.
\end{solution}

%---------------------------------------------------------------------
\Q{Rectangular mesh}{Derivation}{6}\label{q:rect}
\begin{qbox}
(Mitchell--Griffiths, Ch.~3, Exercise 1; exercise of Lecture 37.)
Derive the five-point replacement of Laplace's equation at
\((lh_1,mh_2)\):
\[
  \theta^{-1}(U_{l+1,m}+U_{l-1,m})+\theta(U_{l,m+1}+U_{l,m-1})-2(\theta+\theta^{-1})U_{l,m}=0,
  \qquad\theta=\frac{h_1}{h_2},
\]
and show that its local truncation error is
\(\frac1{12}\bigl(h_1^2u_{x_1^4}+h_2^2u_{x_2^4}\bigr)+\cdots\).
\end{qbox}

\begin{rem}
\begin{itemize}
  \item \(u_{l+1,m}+u_{l-1,m}-2u_{l,m}=h_1^2u_{x_1x_1}+\frac{h_1^4}{12}u_{x_1^4}+\cdots\);
    similarly in \(x_2\) with \(h_2\).
  \item Multiply the central-difference equation by \(h_1h_2\);
    divide by \(h_1h_2\) to get the LTE of \(\nabla^2\).
\end{itemize}
\end{rem}

\begin{solution}
\step{Derivation}
Central differences give
\[
  \frac{U_{l+1,m}-2U_{l,m}+U_{l-1,m}}{h_1^2}+\frac{U_{l,m+1}-2U_{l,m}+U_{l,m-1}}{h_2^2}=0 .
\]
Multiplying by \(h_1h_2\), with \(h_2/h_1=\theta^{-1}\), \(h_1/h_2=\theta\):
\(\theta^{-1}(U_{l+1,m}-2U_{l,m}+U_{l-1,m})+\theta(U_{l,m+1}-2U_{l,m}+U_{l,m-1})=0\),
which is the stated formula.

\step{Truncation error}
Substituting the exact solution and dividing by \(h_1h_2\):
\[
  \frac1{h_1^2}\Bigl(h_1^2u_{x_1x_1}+\frac{h_1^4}{12}u_{x_1^4}\Bigr)
  +\frac1{h_2^2}\Bigl(h_2^2u_{x_2x_2}+\frac{h_2^4}{12}u_{x_2^4}\Bigr)+\cdots
  =\nabla^2u+\frac1{12}\bigl(h_1^2u_{x_1^4}+h_2^2u_{x_2^4}\bigr)+\cdots .
\]
Since \(\nabla^2u=0\), the LTE is
\(\frac1{12}\bigl(h_1^2u_{x_1^4}+h_2^2u_{x_2^4}\bigr)=O(h_1^2+h_2^2)\).
For \(\theta=1\) it reduces to the square five-point formula.
\end{solution}

%---------------------------------------------------------------------
\Q{Laplace and Poisson problems by the five-point formula}{Numerical}{8}\label{q:lapnum}
\begin{qbox}
On the unit square, with the five-point formula:
\begin{enumerate}[label=(\alph*)]
  \item solve \(\nabla^2u=0\) with \(u=\sin\pi x_1\) on \(x_2=1\) and
    \(u=0\) on the other three sides, for \(h=\frac13\); compare with
    \(u=\sin\pi x_1\,\sinh\pi x_2/\sinh\pi\);
  \item solve \(\nabla^2u=-1\) with \(u=0\) on the boundary, for
    \(h=\frac14\).
\end{enumerate}
\end{qbox}

\begin{rem}
\begin{itemize}
  \item \(U_{l+1,m}+U_{l-1,m}+U_{l,m+1}+U_{l,m-1}-4U_{l,m}=h^2f_{l,m}\).
  \item Use the symmetry of the data to reduce the number of
    unknowns.
\end{itemize}
\end{rem}

\begin{solution}
\step{(a) Laplace, \(h=\frac13\)}
Interior points \((l,m)\), \(l,m=1,2\).  The data are symmetric about
\(x_1=\frac12\), so \(U_{1,1}=U_{2,1}=a\), \(U_{1,2}=U_{2,2}=b\).  The
top boundary values are \(\sin\frac\pi3=\sin\frac{2\pi}3=0.866025\).
\begin{align*}
  (1,1):&\quad a+0+b+0-4a=0\ \Rightarrow\ b=3a,\\
  (1,2):&\quad b+0+0.866025+a-4b=0\ \Rightarrow\ a-3b=-0.866025 .
\end{align*}
So \(-8a=-0.866025\):
\[
  a=0.108253,\qquad b=0.324760 .
\]
Exact: \(u(\frac13,\frac13)=0.093688\), \(u(\frac13,\frac23)=0.299857\);
errors \(0.015\) and \(0.025\), as expected for \(O(h^2)\) with
\(h=\frac13\).

\step{(b) Poisson, \(h=\frac14\)}
Nine interior points; by symmetry about both centre lines and the
diagonals there are three unknowns: corners \(q\) (e.g.\ \((1,1)\)),
edge midpoints \(e\) (e.g.\ \((2,1)\)) and the centre \(c\).  With
\(h^2f=-\frac1{16}\):
\begin{align*}
  (1,1):&\quad 2e-4q=-\tfrac1{16},\\
  (2,1):&\quad 2q+c-4e=-\tfrac1{16},\\
  (2,2):&\quad 4e-4c=-\tfrac1{16}.
\end{align*}
From the first, \(q=\frac e2+\frac1{64}\); from the third,
\(c=e+\frac1{64}\).  Substituting in the second:
\(e+\frac1{32}+e+\frac1{64}-4e=-\frac1{16}\Rightarrow2e=\frac7{64}\).
Hence
\[
  e=\frac7{128}=0.0546875,\qquad q=\frac{11}{256}=0.0429688,\qquad
  c=\frac9{128}=0.0703125 .
\]
(The exact centre value of this problem is about \(0.0737\).)
\end{solution}

%#####################################################################
\Needspace{32\baselineskip}
\section{ADI Iterations: Dirichlet, Neumann and Robin Problems}\label{sec:adi}
%#####################################################################

Notation: \(Nh=1\) (or \(Mh=1\)); \(H\) and \(V\) are the matrices of
\(-\delta_{x_1}^2\) and \(-\delta_{x_2}^2\); \(\lam\) is the iteration
parameter.

%---------------------------------------------------------------------
\Q{Peaceman--Rachford ADI as an iteration for Laplace's equation}{Derivation / proof}{10}\label{q:priter}
\begin{qbox}
\begin{enumerate}[label=(\alph*)]
  \item Show that if the Peaceman--Rachford scheme for
    \(u_t=\nabla^2u\) (time-independent Dirichlet data) converges as
    \(n\to\infty\), the limit satisfies the five-point Laplace
    equation.
  \item Write the iteration in matrix form
    \(\bigl(\frac2{\lam}I+H\bigr)U^{n+1*}=\bigl(\frac2{\lam}I-V\bigr)U^n+K\),
    \(\bigl(\frac2{\lam}I+V\bigr)U^{n+1}=\bigl(\frac2{\lam}I-H\bigr)U^{n+1*}+K\),
    and show by eliminating \(U^{n+1*}\) that \(U^{n+1}=U^n=U\) gives
    \((H+V)U=K\) (Mitchell--Griffiths, Ch.~3, Exercise 16).
  \item Prove that the iteration converges for every choice of
    \(\lam>0\).
\end{enumerate}
\end{qbox}

\begin{rem}
\begin{itemize}
  \item For the unit square with \((N-1)^2\) interior points:
    \(H=\mathrm{diag}(C,\dots,C)\), \(V=\mathrm{tridiag}(-J,2J,-J)\),
    \(C=\mathrm{tridiag}(-1,2,-1)\) of order \(N-1\); \(HV=VH\).
  \item Common eigenvectors \(\alpha_{p,q}\):
    \(H\alpha_{p,q}=4\sin^2\frac{p\pi}{2N}\alpha_{p,q}\),
    \(V\alpha_{p,q}=4\sin^2\frac{q\pi}{2N}\alpha_{p,q}\).
  \item Iteration matrix
    \(T_{n+1}=\bigl(\frac2{\lam}I+V\bigr)^{-1}\bigl(\frac2{\lam}I+H\bigr)^{-1}\bigl(\frac2{\lam}I-H\bigr)\bigl(\frac2{\lam}I-V\bigr)\);
    converges iff \(\rho(T)<1\).
  \item \(\lambda\) is now an \emph{iteration parameter}, not
    \(k/h^2\).
\end{itemize}
\end{rem}

\begin{solution}
\step{(a) Limit of the scheme}
Eliminating \(U^{n+1*}\) from Peaceman--Rachford,
\(\bigl(1-\frac\lambda2\delta_{x_1}^2\bigr)\bigl(1-\frac\lambda2\delta_{x_2}^2\bigr)U^{n+1}
=\bigl(1+\frac\lambda2\delta_{x_1}^2\bigr)\bigl(1+\frac\lambda2\delta_{x_2}^2\bigr)U^n\).
If \(U^{n+1}=U^n=U\), expanding both sides,
\[
  \Bigl(1-\tfrac\lambda2(\delta_{x_1}^2+\delta_{x_2}^2)+\tfrac{\lambda^2}4\delta_{x_1}^2\delta_{x_2}^2\Bigr)U
  =\Bigl(1+\tfrac\lambda2(\delta_{x_1}^2+\delta_{x_2}^2)+\tfrac{\lambda^2}4\delta_{x_1}^2\delta_{x_2}^2\Bigr)U,
\]
so \(\lambda(\delta_{x_1}^2+\delta_{x_2}^2)U=0\), i.e.\
\((\delta_{x_1}^2+\delta_{x_2}^2)U_{l,m}=0\): the five-point formula.

\step{(b) Matrix form and its limit}
Dividing each half-step by \(\lambda/2\) and moving boundary values to
\(K\) gives the stated equations.  From the first,
\(U^{n+1*}=\bigl(\frac2{\lam}I+H\bigr)^{-1}\bigl[\bigl(\frac2{\lam}I-V\bigr)U^n+K\bigr]\).
Substituting in the second and multiplying by
\(\bigl(\frac2{\lam}I+H\bigr)\), which commutes with
\(\bigl(\frac2{\lam}I-H\bigr)\),
\[
  \Bigl(\tfrac2{\lam}I+H\Bigr)\Bigl(\tfrac2{\lam}I+V\Bigr)U^{n+1}
  =\Bigl(\tfrac2{\lam}I-H\Bigr)\Bigl(\tfrac2{\lam}I-V\Bigr)U^n
  +\Bigl(\tfrac2{\lam}I-H\Bigr)K+\Bigl(\tfrac2{\lam}I+H\Bigr)K .
\]
With \(U^{n+1}=U^n=U\), the \(\frac4{\lam^2}U\) and \(HVU\) terms
cancel, leaving \(\frac4{\lam}(H+V)U=\frac4{\lam}K\), i.e.\
\((H+V)U=K\).

\step{(c) Convergence}
\(U^{n+1}=T_{n+1}U^n+g_{n+1}\).  Since \(HV=VH\), they have common
orthonormal eigenvectors \(\alpha_{p,q}\), \(1\le p,q\le N-1\), and
\[
  T_{n+1}\alpha_{p,q}
  =\frac{\bigl(\frac2{\lam}-4\sin^2\frac{p\pi}{2N}\bigr)\bigl(\frac2{\lam}-4\sin^2\frac{q\pi}{2N}\bigr)}
        {\bigl(\frac2{\lam}+4\sin^2\frac{p\pi}{2N}\bigr)\bigl(\frac2{\lam}+4\sin^2\frac{q\pi}{2N}\bigr)}\alpha_{p,q}
  =\frac{(1-a)(1-b)}{(1+a)(1+b)}\alpha_{p,q},
\]
with \(a=2\lam\sin^2\frac{p\pi}{2N}>0\), \(b=2\lam\sin^2\frac{q\pi}{2N}>0\).
Now
\[
  \frac{(1-a)(1-b)}{(1+a)(1+b)}=1-\frac{2(a+b)}{(1+a)(1+b)}\in(-1,1),
\]
since \(a+b>0\) and \(a+b<1+a+b+ab\).  So \(\rho(T_{n+1})<1\) for every
\(\lam>0\).  The error \(e^n=U^n-U\) satisfies \(e^{n+1}=T_{n+1}e^n\),
and all \(T_{n+1}\) share the eigenvectors \(\alpha_{p,q}\), so each
component of the error is multiplied by a factor of modulus less than
\(1\) at every iteration: the iteration converges.
\end{solution}

%---------------------------------------------------------------------
\Q{The optimal fixed iteration parameter}{Derivation}{6}\label{q:optpar}
\begin{qbox}
(Mitchell--Griffiths, Ch.~3, Example 11.)  For the Peaceman--Rachford
iteration on the unit square with \(h=1/M\), show that the fixed
parameter which minimises the spectral radius of the iteration matrix
is \(\lambda=\csc(\pi/M)\).  Find the resulting spectral radius for
\(M=11\).
\end{qbox}

\begin{rem}
\begin{itemize}
  \item Eigenvalues of \(T\):
    \(\frac{(c-\mu_p)(c-\mu_q)}{(c+\mu_p)(c+\mu_q)}\), \(c=\frac2\lambda\),
    \(\mu_p=4\sin^2\frac{p\pi}{2M}\in[\mu_{\min},\mu_{\max}]\).
  \item \(\bigl|\frac{c-\mu}{c+\mu}\bigr|\) is largest at an end point of
    \([\mu_{\min},\mu_{\max}]\).
  \item Optimum: \(c=\sqrt{\mu_{\min}\mu_{\max}}\).
\end{itemize}
\end{rem}

\begin{solution}
\step{Minimising the spectral radius}
\(\rho(T)=\bigl[\max_p\bigl|\frac{c-\mu_p}{c+\mu_p}\bigr|\bigr]^2\).  The
function \(\mu\mapsto\frac{c-\mu}{c+\mu}\) is decreasing, so the maximum
of its modulus over \([\mu_{\min},\mu_{\max}]\) is
\(\max\bigl(\frac{c-\mu_{\min}}{c+\mu_{\min}},\frac{\mu_{\max}-c}{\mu_{\max}+c}\bigr)\).
The first decreases and the second increases in \(c\), so the
maximum is smallest when they are equal:
\[
  (c-\mu_{\min})(\mu_{\max}+c)=(\mu_{\max}-c)(c+\mu_{\min})
  \iff c^2=\mu_{\min}\mu_{\max} .
\]
\step{The value}
\(\mu_{\min}=4\sin^2\frac\pi{2M}\), \(\mu_{\max}=4\sin^2\frac{(M-1)\pi}{2M}=4\cos^2\frac\pi{2M}\), so
\[
  c=4\sin\frac\pi{2M}\cos\frac\pi{2M}=2\sin\frac\pi M,
  \qquad
  \lambda=\frac2c=\csc\frac\pi M .
\]
\step{\(M=11\)}
\(\lambda=\csc\frac\pi{11}=3.5495\).  Then
\(\frac{c-\mu_{\min}}{c+\mu_{\min}}=\frac{1-\tan(\pi/22)}{1+\tan(\pi/22)}=0.7486\),
and \(\rho(T)=0.7486^2=0.5604\): each iteration reduces the error by
a factor of about \(0.56\).  (Wachspress' varying parameters
\(\lam\) do even better.)
\end{solution}

%---------------------------------------------------------------------
\Q{The Douglas--Rachford ADI method}{Derivation / analysis}{8}\label{q:dr}
\begin{qbox}
(Lecture 39; Mitchell--Griffiths, Ch.~2, Exercise 15.)  For
\(u_t=\nabla^2u\), consider
\((1-r\delta_{x_1}^2)(1-r\delta_{x_2}^2)U^{n+1}=(1+r^2\delta_{x_1}^2\delta_{x_2}^2)U^n\).
\begin{enumerate}[label=(\alph*)]
  \item Show that the split form
    \((1-r\delta_{x_1}^2)U^{n+1*}=(1+r\delta_{x_2}^2)U^n\),
    \((1-r\delta_{x_2}^2)U^{n+1}=U^{n+1*}-r\delta_{x_2}^2U^n\) gives the
    combined formula.
  \item Show that, used as an iteration, its limit is the five-point
    Laplace equation.
  \item Show that its truncation error is \(O(k^2+kh^2)\), and that,
    with \(U^{n+1*}=U^{n+1}\), the split formulas are consistent.
  \item Show that it is unconditionally stable.
\end{enumerate}
\end{qbox}

\begin{rem}
\begin{itemize}
  \item Douglas--Rachford \(=\) fully implicit
    \(\bigl(1-r(\delta_{x_1}^2+\delta_{x_2}^2)\bigr)U^{n+1}=U^n\) plus the
    small term \(r^2\delta_{x_1}^2\delta_{x_2}^2(U^{n+1}-U^n)\).
  \item Fully implicit LTE: \(-\frac{k^2}2u_{tt}-\frac{kh^2}{12}(u_{x_1^4}+u_{x_2^4})\).
  \item On a Fourier mode \(\delta_{x_i}^2\to-4S_i\), with
    \(S_i=\sin^2\frac{\beta_ih}2\).
\end{itemize}
\end{rem}

\begin{solution}
\step{(a) Split \(\Rightarrow\) combined}
From the second equation
\(U^{n+1*}=(1-r\delta_{x_2}^2)U^{n+1}+r\delta_{x_2}^2U^n\).  Applying
\((1-r\delta_{x_1}^2)\) and using the first,
\[
  (1-r\delta_{x_1}^2)(1-r\delta_{x_2}^2)U^{n+1}
  =(1+r\delta_{x_2}^2)U^n-r\delta_{x_2}^2(1-r\delta_{x_1}^2)U^n
  =(1+r^2\delta_{x_1}^2\delta_{x_2}^2)U^n .
\]
\step{(b) Limit}
With \(U^{n+1}=U^n=U\):
\(\bigl(1-r(\delta_{x_1}^2+\delta_{x_2}^2)+r^2\delta_{x_1}^2\delta_{x_2}^2\bigr)U
=(1+r^2\delta_{x_1}^2\delta_{x_2}^2)U\), so
\((\delta_{x_1}^2+\delta_{x_2}^2)U=0\).

\step{(c) Truncation error}
Write the formula as
\(\bigl[1-r(\delta_{x_1}^2+\delta_{x_2}^2)\bigr]U^{n+1}-U^n
+r^2\delta_{x_1}^2\delta_{x_2}^2(U^{n+1}-U^n)=0\).  The first part is
the fully implicit scheme, with \(T=-\frac{k^2}2u_{tt}-\frac{kh^2}{12}(u_{x_1^4}+u_{x_2^4})+\cdots\).
The extra term is
\(r^2h^4\,k\,u_{x_1x_1x_2x_2t}+\cdots=k^3u_{x_1x_1x_2x_2t}+\cdots\).
So \(T=O(k^2+kh^2)\), and \(T/k=O(k+h^2)\): consistent.

With \(U^{n+1*}=U^{n+1}\): the first split formula becomes
\(u^{n+1}-u^n-r\delta_{x_1}^2u^{n+1}-r\delta_{x_2}^2u^n=k(u_t-\nabla^2u)+O(k^2+kh^2)\),
consistent; the second becomes
\(-r\delta_{x_2}^2(u^{n+1}-u^n)=-k^2u_{x_2x_2t}+\cdots\), which divided by
\(k\) tends to zero: also consistent.

\step{(d) Stability}
With \(S_i=\sin^2\frac{\beta_ih}2\),
\[
  \xi=\frac{1+16r^2S_1S_2}{(1+4rS_1)(1+4rS_2)}
  =\frac{1+16r^2S_1S_2}{1+4r(S_1+S_2)+16r^2S_1S_2}\in(0,1]
\]
for every \(r>0\): unconditionally stable.
\end{solution}

%---------------------------------------------------------------------
\Q{The Neumann problem: discretization and ADI convergence}{Derivation / proof}{10}\label{q:neumann}
\begin{qbox}
Consider \(\nabla^2u=0\) on \(0\le x_1,x_2\le1\) with
\(\partial u/\partial n=g(x_1,x_2)\) (outward normal) on the boundary,
and \(Mh=1\).
\begin{enumerate}[label=(\alph*)]
  \item Using fictitious points, derive the difference equations at a
    corner and at a side point, and show that the system is
    \(AU=2hG\) with \(A\) of order \((M+1)^2\).
  \item Show that \(A\) is singular and explain why.
  \item Prove that the Peaceman--Rachford ADI iteration for this
    problem converges for any fixed value of the iteration parameter
    (in the sense explained).
\end{enumerate}
\end{qbox}

\begin{rem}
\begin{itemize}
  \item All \((M+1)^2\) points are unknown.  Along \(x_1=0\) the
    central difference of \(\partial u/\partial n=-u_{x_1}\) gives
    \(u_{-1,m}=u_{1,m}+2hg_{0,m}\).
  \item \(A=X+Y\), \(X=\mathrm{diag}(\widetilde C,\dots,\widetilde C)\),
    \(\widetilde C=\) tridiag\((-1,2,-1)\) with first row \((2,-2)\) and
    last row \((-2,2)\); \(Y\) the same along \(x_2\).
  \item Eigenvalues of \(\widetilde C\): \(4\sin^2\frac{p\pi}{2M}\),
    \(p=0,\dots,M\), eigenvectors \(\bigl(\cos\frac{lp\pi}M\bigr)_l\).
\end{itemize}
\end{rem}

\begin{solution}
\step{(a) Boundary equations}
At the corner \((0,0)\), the five-point formula is
\(u_{1,0}+u_{-1,0}+u_{0,1}+u_{0,-1}-4u_{0,0}=0\).  Central differences
for the outward normal derivatives give
\(u_{-1,0}=u_{1,0}+2hg_{0,0}\) and \(u_{0,-1}=u_{0,1}+2hg_{0,0}\), so
\[
  4u_{0,0}-2u_{1,0}-2u_{0,1}=2h\cdot2g_{0,0}.
\]
At a side point \((0,m)\):
\(4u_{0,m}-2u_{1,m}-u_{0,m+1}-u_{0,m-1}=2hg_{0,m}\); similarly on the
other sides.  With the \(x_1\) index running fastest,
\[
  A=\begin{bmatrix}B&-2I&&\\-I&B&-I&\\&\ddots&\ddots&\ddots\\&&-2I&B\end{bmatrix},
  \qquad
  B=\begin{bmatrix}4&-2&&\\-1&4&-1&\\&\ddots&\ddots&\ddots\\&&-2&4\end{bmatrix}_{(M+1)\times(M+1)}
\]
and \(G\) has entries \(2g\) at the corners, \(g\) at other boundary
points and \(0\) inside.

\step{(b) Singularity}
Every row of \(A\) sums to zero (corner \(4-2-2\), side \(4-2-1-1\),
interior \(4-1-1-1-1\)), so \(A\mathbf1=0\): \(A\) is singular.  This
reflects the continuous problem: a solution of the Neumann problem is
determined only up to an additive constant, and it exists only if
\(\oint_{\partial R}g\,ds=0\) (a discrete version is derived in
Q\ref{q:neu9}).

\step{(c) Convergence}
The iteration is
\(\bigl(\frac2{\lam}I+X\bigr)U^{n+1*}=\bigl(\frac2{\lam}I-Y\bigr)U^n+2hG\),
\(\bigl(\frac2{\lam}I+Y\bigr)U^{n+1}=\bigl(\frac2{\lam}I-X\bigr)U^{n+1*}+2hG\).

\emph{Limit.} From the first equation,
\(U^{n+1*}=\bigl(\frac2{\lam}I+X\bigr)^{-1}\bigl[\bigl(\frac2{\lam}I-Y\bigr)U^n+2hG\bigr]\).
Substituting in the second and multiplying by
\(\bigl(\frac2{\lam}I+X\bigr)\), which commutes with
\(\bigl(\frac2{\lam}I-X\bigr)\),
\begin{align*}
  \Bigl(\tfrac2{\lam}I+X\Bigr)\Bigl(\tfrac2{\lam}I+Y\Bigr)U^{n+1}
  &=\Bigl(\tfrac2{\lam}I-X\Bigr)\Bigl(\tfrac2{\lam}I-Y\Bigr)U^n\\
  &\quad+\Bigl[\Bigl(\tfrac2{\lam}I-X\Bigr)+\Bigl(\tfrac2{\lam}I+X\Bigr)\Bigr]2hG .
\end{align*}
With \(U^{n+1}=U^n=U\), expanding both products, the
\(\frac4{\lam^2}U\) and \(XYU\) terms cancel and
\(\frac4{\lam}(X+Y)U=\frac4{\lam}\,2hG\), i.e.\ \(AU=2hG\): the limit is a
solution of the difference equations.

\emph{Error propagation.}  The iteration matrix is
\(S=\bigl(\frac2{\lam}I+Y\bigr)^{-1}\bigl(\frac2{\lam}I+X\bigr)^{-1}\bigl(\frac2{\lam}I-X\bigr)\bigl(\frac2{\lam}I-Y\bigr)\).
\(X\) and \(Y\) act along different directions, so \(XY=YX\), with
common eigenvectors \(\beta_{p,q}\) (\(0\le p,q\le M\)):
\(X\beta_{p,q}=4\sin^2\frac{p\pi}{2M}\beta_{p,q}\),
\(Y\beta_{p,q}=4\sin^2\frac{q\pi}{2M}\beta_{p,q}\).  Hence
\[
  S\beta_{p,q}=\frac{(1-a)(1-b)}{(1+a)(1+b)}\beta_{p,q},
  \qquad a=2\lam\sin^2\frac{p\pi}{2M}\ge0,\ b=2\lam\sin^2\frac{q\pi}{2M}\ge0 .
\]
For \((p,q)\ne(0,0)\), \(a+b>0\) and the modulus is less than \(1\).
For \(p=q=0\) the eigenvalue is \(1\), with eigenvector \(\mathbf1\), the
null vector of \(A\).  So every error component except the constant one
decays: when the data are compatible, the iteration converges, for any
fixed \(\lambda>0\), to a solution of \(AU=2hG\), the additive constant
being fixed by the initial guess.
\end{solution}

%---------------------------------------------------------------------
\Q{A Neumann problem with \(h=\frac12\)}{Numerical / proof}{8}\label{q:neu9}
\begin{qbox}
(Mitchell--Griffiths, Ch.~3, Example 1.)  For the Neumann problem of
Q\ref{q:neumann} with \(h=\frac12\), assemble the \(9\times9\) system
\(AU=2hG\).  Show that scaling the nine equations by
\(\frac14,\frac12,\frac14,\frac12,1,\frac12,\frac14,\frac12,\frac14\)
makes the coefficient matrix symmetric, verify that \(A\) is singular,
and deduce the discrete compatibility condition
\(\sum_{i=1}^9g_i=0\) (with \(g_5=0\)).
\end{qbox}

\begin{rem}
\begin{itemize}
  \item Number the points \(1,\dots,9\) row by row from \((0,0)\);
    point 5 is the centre.
  \item \(AU=\mathbf b\) with \(A\) singular is solvable iff
    \(\mathbf w^T\mathbf b=0\) for every \(\mathbf w\) with
    \(\mathbf w^TA=0\).
\end{itemize}
\end{rem}

\begin{solution}
\step{The system}
With points \(1\)--\(3\) on \(x_2=0\), \(4\)--\(6\) on \(x_2=\frac12\),
\(7\)--\(9\) on \(x_2=1\), the five-point formula with fictitious points
(\(u_{-1,m}=u_{1,m}+2hg_{0,m}\), etc.) gives, at a corner,
\(4u_c-2u-2u=2h\cdot2g_c\), at a side point
\(4u_s-2u_{\text{inner}}-u-u=2hg_s\), and at the centre the ordinary
five-point equation.  Hence
\[
  A=\begin{bmatrix}
   4&-2&0&-2&0&0&0&0&0\\
   -1&4&-1&0&-2&0&0&0&0\\
   0&-2&4&0&0&-2&0&0&0\\
   -1&0&0&4&-2&0&-1&0&0\\
   0&-1&0&-1&4&-1&0&-1&0\\
   0&0&-1&0&-2&4&0&0&-1\\
   0&0&0&-2&0&0&4&-2&0\\
   0&0&0&0&-2&0&-1&4&-1\\
   0&0&0&0&0&-2&0&-2&4
  \end{bmatrix},
\]
and \(2hG=G=(2g_1,g_2,2g_3,g_4,0,g_6,2g_7,g_8,2g_9)^T\) (\(2h=1\)).

\step{Symmetrisation}
Let \(W=\mathrm{diag}(\frac14,\frac12,\frac14,\frac12,1,\frac12,\frac14,\frac12,\frac14)\).
Row \(i\) of \(WA\) is row \(i\) of \(A\) times \(w_i\).  For example
\((WA)_{12}=\frac14(-2)=-\frac12=\frac12(-1)=(WA)_{21}\), and
\((WA)_{25}=\frac12(-2)=-1=1\cdot(-1)=(WA)_{52}\); the same holds for
every pair (corner--side: \(-\frac12\); side--centre: \(-1\)).  So
\(WA\) is symmetric.

\step{Singularity}
Each row of \(A\) sums to zero, so \(A\mathbf1=0\) and \(A\) is
singular (rank \(8\)).

\step{Compatibility}
Since \(WA\) is symmetric, \(\mathbf1^TWA=(WA\mathbf1)^T=0\): the vector
\(\mathbf w=W\mathbf1=(\frac14,\frac12,\frac14,\frac12,1,\frac12,\frac14,\frac12,\frac14)^T\)
satisfies \(\mathbf w^TA=0\).  So \(AU=G\) is solvable only if
\(\mathbf w^TG=0\):
\[
  \tfrac14(2g_1+2g_3+2g_7+2g_9)+\tfrac12(g_2+g_4+g_6+g_8)
  =\tfrac12\sum_{i\ne5}g_i=0
  \iff\sum_{i=1}^9g_i=0\quad(g_5=0).
\]
This is the discrete form of \(\oint_{\partial R}\frac{\partial u}{\partial n}ds=0\).
\end{solution}

%---------------------------------------------------------------------
\Q{The Robin problem: matrices and convergence}{Derivation / proof}{10}\label{q:robin}
\begin{qbox}
For \(\nabla^2u=0\) on \(0\le x_1,x_2\le1\) with
\(u_{x_1}-p_1u=f_0(x_2)\) on \(x_1=0\), \(u_{x_1}+q_1u=f_1(x_2)\) on
\(x_1=1\), and similarly with \(p_2,q_2,g_0,g_1\) on \(x_2=0,1\)
(\(p_i,q_i>0\)), \(Nh=1\):
\begin{enumerate}[label=(\alph*)]
  \item derive the matrices \(E\) and \(V_1\) of the system
    \(AU=hG\), \(A=H_1+V_1\), \(H_1=\mathrm{diag}(E,\dots,E)\);
  \item prove that the Peaceman--Rachford iteration converges for
    every \(\lam>0\).
\end{enumerate}
\end{qbox}

\begin{rem}
\begin{itemize}
  \item Fictitious points:
    \(u_{-1,m}=u_{1,m}-2hp_1u_{0,m}-2hf_{0,m}\),
    \(u_{N+1,m}=u_{N-1,m}-2hq_1u_{N,m}+2hf_{1,m}\).
  \item \(W=\mathrm{diag}(\frac12,1,\dots,1,\frac12)\) makes \(WE\)
    symmetric positive definite, so \(E\) has real positive
    eigenvalues (Q\ref{q:matfacts}).
  \item \(\bigl|\frac{1-a}{1+a}\bigr|<1\) for \(a>0\).
\end{itemize}
\end{rem}

\begin{solution}
\step{(a) The matrices}
The \(x_1\)-part of the five-point formula (times \(-1\)) at \((0,m)\)
is \(2u_{0,m}-u_{1,m}-u_{-1,m}=2(1+hp_1)u_{0,m}-2u_{1,m}+2hf_{0,m}\),
and at \((N,m)\) it is \(-2u_{N-1,m}+2(1+hq_1)u_{N,m}-2hf_{1,m}\).  So,
ordering with the \(x_1\) index fastest,
\[
  E=\begin{bmatrix}2(1+hp_1)&-2&&\\-1&2&-1&\\&\ddots&\ddots&\ddots\\&&-2&2(1+hq_1)\end{bmatrix}_{(N+1)\times(N+1)}
\]
\[
  V_1=\begin{bmatrix}2(1+hp_2)I&-2I&&\\-I&2I&-I&\\&\ddots&\ddots&\ddots\\&&-2I&2(1+hq_2)I\end{bmatrix},
\]
and \(G_{l,m}=2\bigl(-f_{0,m}[l=0]+f_{1,m}[l=N]-g_{0,l}[m=0]+g_{1,l}[m=N]\bigr)\).

\step{(b) Eigenvalues of \(E\) are positive}
\(WE\) has rows \((1+hp_1,-1)\), \((-1,2,-1)\), \((-1,1+hq_1)\): it is
symmetric, and
\[
  \mathbf x^T(WE)\mathbf x=hp_1x_0^2+hq_1x_N^2+\sum_{l=0}^{N-1}(x_{l+1}-x_l)^2>0
  \quad(\mathbf x\ne0),
\]
since a zero sum forces \(\mathbf x\) constant and then
\(hp_1x_0^2>0\).  \(E=W^{-1}(WE)\) is similar to
\(W^{-1/2}(WE)W^{-1/2}\), symmetric positive definite, so its
eigenvalues \(\mu_p\) are real and positive and it has a complete set
of eigenvectors.  Likewise for the \(x_2\) matrix (eigenvalues
\(\nu_q>0\)).

\step{Convergence}
\(H_1\) acts along \(x_1\) and \(V_1\) along \(x_2\), so they commute,
with common eigenvectors.  The iteration matrix has eigenvalues
\[
  \frac{\bigl(\frac2{\lam}-\mu_p\bigr)\bigl(\frac2{\lam}-\nu_q\bigr)}{\bigl(\frac2{\lam}+\mu_p\bigr)\bigl(\frac2{\lam}+\nu_q\bigr)}
  =\frac{1-a}{1+a}\cdot\frac{1-b}{1+b},
  \qquad a=\tfrac12\lam\mu_p>0,\ b=\tfrac12\lam\nu_q>0,
\]
each factor of modulus less than \(1\).  So \(\rho<1\) and the
iteration converges to the unique solution of \((H_1+V_1)U=hG\)
(unique because all eigenvalues \(\mu_p+\nu_q\) of \(A\) are positive,
unlike the Neumann case).
\end{solution}

%#####################################################################
\Needspace{32\baselineskip}
\section{Hyperbolic Equations: The Wave Equation}
%#####################################################################

Throughout, \(u_{tt}=u_{xx}\) (or \(u_{tt}=\nabla^2u\)) and \(r=k/h\)
(written \(p\) in Mitchell--Griffiths).

%---------------------------------------------------------------------
\Q{The explicit scheme: derivation, starting values and truncation error}{Derivation}{8}\label{q:waveexp}
\begin{qbox}
For \(u_{tt}=u_{xx}\), \(0<x<1\), with \(u(x,0)=f(x)\),
\(u_t(x,0)=g(x)\), \(u(0,t)=u(1,t)=0\):
\begin{enumerate}[label=(\alph*)]
  \item derive the explicit scheme
    \(U_m^{n+1}=2(1-r^2)U_m^n+r^2(U_{m+1}^n+U_{m-1}^n)-U_m^{n-1}\);
  \item derive the formula for the first time level
    \(U_m^1=(1-r^2)f_m+\frac{r^2}2(f_{m+1}+f_{m-1})+kg_m\);
  \item show that the truncation error is \(O(k^2+h^2)\), and that the
    leading term vanishes when \(r=1\).
\end{enumerate}
\end{qbox}

\begin{rem}
\begin{itemize}
  \item \(\frac1{k^2}\delta_t^2u=u_{tt}+\frac{k^2}{12}u_{tttt}+\cdots\),
    \(\frac1{h^2}\delta_x^2u=u_{xx}+\frac{h^2}{12}u_{xxxx}+\cdots\).
  \item Three levels are the minimum (second order in \(t\)); the start
    needs two levels.
  \item Central difference for \(u_t(x,0)\):
    \(U_m^{-1}=U_m^1-2kg_m\).
  \item Consistent (hyperbolic) if \(T/k^2\to0\) for the scheme
    multiplied through by \(k^2\).
\end{itemize}
\end{rem}

\begin{solution}
\step{(a) The scheme}
At \((mh,nk)\),
\(\frac1{k^2}\bigl(U_m^{n+1}-2U_m^n+U_m^{n-1}\bigr)=\frac1{h^2}\bigl(U_{m+1}^n-2U_m^n+U_{m-1}^n\bigr)\).
Multiplying by \(k^2\) (\(r^2=k^2/h^2\)) and solving for \(U_m^{n+1}\)
gives the stated explicit three-level formula.

\step{(b) First level}
\(U_m^0=f_m\).  From
\(\frac{U_m^1-U_m^{-1}}{2k}=g_m+O(k^2)\), \(U_m^{-1}=U_m^1-2kg_m\).
Putting \(n=0\) in the scheme:
\(U_m^1=2(1-r^2)f_m+r^2(f_{m+1}+f_{m-1})-U_m^1+2kg_m\), so
\[
  U_m^1=(1-r^2)f_m+\frac{r^2}2(f_{m+1}+f_{m-1})+kg_m .
\]

\step{(c) Truncation error}
For the exact solution,
\[
  \frac{\delta_t^2u}{k^2}-\frac{\delta_x^2u}{h^2}
  =(u_{tt}-u_{xx})+\frac1{12}\bigl(k^2u_{tttt}-h^2u_{xxxx}\bigr)+\cdots
  =\frac1{12}\bigl(k^2u_{tttt}-h^2u_{xxxx}\bigr)+\cdots=O(k^2+h^2).
\]
Since \(u_{tttt}=u_{xxxx}\), the leading term is
\(\frac{h^2}{12}(r^2-1)u_{xxxx}\), which vanishes for \(r=1\) (see
Q\ref{q:dalembert}).  So the scheme is consistent.
\end{solution}

%---------------------------------------------------------------------
\Q{Stability of the explicit scheme by the amplification matrix: the CFL condition}{Stability analysis}{6}\label{q:cfl}
\begin{qbox}
(Mitchell--Griffiths, Ch.~4, Exercise 10.)  Write the explicit scheme
as the two-level system
\(U_m^{n+1}=2(1-r^2)U_m^n+r^2(U_{m+1}^n+U_{m-1}^n)-V_m^n\),
\(V_m^{n+1}=U_m^n\).  Obtain the amplification matrix and show that the
scheme is stable if \(r\le1\).  What is the CFL condition?
\end{qbox}

\begin{rem}
\begin{itemize}
  \item Von Neumann for a system: all eigenvalues of the amplification
    matrix have modulus \(\le1\).
  \item \(U_{m+1}+U_{m-1}=2\cos\beta h\,U_m=(2-4S)U_m\) on a mode,
    \(S=\sin^2\frac{\beta h}2\).
  \item CFL: the numerical domain of dependence must contain the
    exact one.
\end{itemize}
\end{rem}

\begin{solution}
\step{Amplification matrix}
With \((U,V)_m^n=(U^n,V^n)e^{i\beta mh}\):
\[
  \begin{bmatrix}U\\V\end{bmatrix}^{n+1}
  =\underbrace{\begin{bmatrix}2-4r^2S&-1\\1&0\end{bmatrix}}_{G}
  \begin{bmatrix}U\\V\end{bmatrix}^{n}.
\]
\step{Eigenvalues}
\(\det(G-\lambda I)=\lambda^2-(2-4r^2S)\lambda+1=0\).  The product of
the roots is \(1\).  If the roots are real and distinct, one exceeds
\(1\) in modulus.  So we need complex (or coincident) roots:
\(|2-4r^2S|\le2\iff r^2S\le1\); then \(|\lambda_1|=|\lambda_2|=1\).  For
all modes (\(S\le1\)) this gives \(\boxed{r\le1}\).  (Coincident roots
\(\lambda=-1\) occur only for \(r=1\), \(S=1\), i.e.\ \(\beta h=\pi\), a mode
not present on the grid with zero boundary values.)

\step{CFL}
The exact solution at \((x,t_n)\) depends on the data in
\([x-t_n,x+t_n]\); the scheme uses \([x-nh,x+nh]\).  The CFL
(Courant--Friedrichs--Lewy) condition
\([x-nk,x+nk]\subseteq[x-nh,x+nh]\) is \(k\le h\), i.e.\ \(r\le1\): here
it coincides with the stability condition.
\end{solution}

%---------------------------------------------------------------------
\Q{Two steps of the explicit scheme}{Numerical}{6}\label{q:wavenum}
\begin{qbox}
Solve \(u_{tt}=u_{xx}\), \(u(x,0)=\sin\pi x\), \(u_t(x,0)=0\),
\(u(0,t)=u(1,t)=0\), by the explicit scheme with \(h=\frac14\) and
\(r=\frac12\).  Find \(U\) at \(t=k\) and \(t=2k\), and compare with
\(u=\cos\pi t\sin\pi x\).
\end{qbox}

\begin{rem}
\begin{itemize}
  \item \(U_m^1=(1-r^2)f_m+\frac{r^2}2(f_{m+1}+f_{m-1})+kg_m\).
  \item \(U_m^{n+1}=2(1-r^2)U_m^n+r^2(U_{m+1}^n+U_{m-1}^n)-U_m^{n-1}\);
    stable for \(r\le1\).
\end{itemize}
\end{rem}

\begin{solution}
\(k=rh=\frac18\), \(r^2=\frac14\).  \(U^0=(0,\ 0.707107,\ 1,\ 0.707107,\ 0)\);
by symmetry \(U_1=U_3\).

\step{\(t=k\)}
\(U_m^1=\frac34f_m+\frac18(f_{m+1}+f_{m-1})\):
\[
  U_2^1=\tfrac34(1)+\tfrac18(1.414214)=0.926777,\qquad
  U_1^1=\tfrac34(0.707107)+\tfrac18(1)=0.655330 .
\]
\step{\(t=2k\)}
\(U_m^2=\frac32U_m^1+\frac14(U_{m+1}^1+U_{m-1}^1)-U_m^0\):
\[
  U_2^2=1.390165+0.327665-1=0.717830,\qquad
  U_1^2=0.982995+0.231694-0.707107=0.507583 .
\]
\step{Comparison}
\begin{center}
\begin{tabular}{@{}ccccc@{}}
\toprule
\(t\) & \(U(\frac12)\) & exact & \(U(\frac14)\) & exact\\
\midrule
\(\frac18\) & 0.926777 & 0.923880 & 0.655330 & 0.653281\\
\(\frac14\) & 0.717830 & 0.707107 & 0.507583 & 0.500000\\
\bottomrule
\end{tabular}
\end{center}
Errors of about \(1\%\) at \(t=\frac14\), consistent with
\(O(k^2+h^2)\).
\end{solution}

%---------------------------------------------------------------------
\Q{The case \(r=1\) and d'Alembert's solution}{Proof}{6}\label{q:dalembert}
\begin{qbox}
(Mitchell--Griffiths, Ch.~4, Example 12.)  On a square grid (\(k=h\),
\(r=1\)) the explicit scheme is
\(U_m^{n+1}=U_{m+1}^n+U_{m-1}^n-U_m^{n-1}\).  Show that every function
\(F(x+t)+G(x-t)\) satisfies it exactly, and deduce that the scheme
reproduces d'Alembert's solution
\(u=\frac12\bigl[f(x+t)+f(x-t)+\int_{x-t}^{x+t}g(\xi)\,d\xi\bigr]\) up to the
error of the starting values.
\end{qbox}

\begin{rem}
\begin{itemize}
  \item General solution of \(u_{tt}=u_{xx}\): \(F(x+t)+G(x-t)\);
    characteristics \(x\pm t=\)const.
  \item For \(k=h\), the grid points lie on characteristics.
\end{itemize}
\end{rem}

\begin{solution}
\step{Exactness}
Let \(u=F(x+t)+G(x-t)\), \(x=mh\), \(t=nk=nh\).  Then
\begin{align*}
  u_{m+1}^n+u_{m-1}^n-u_m^{n-1}
  &=F(x+t+h)+G(x-t+h)+F(x+t-h)+G(x-t-h)\\
  &\quad-F(x+t-h)-G(x-t+h)\\
  &=F(x+t+h)+G(x-t-h)=u_m^{n+1}.
\end{align*}
So the scheme has zero truncation error for every solution of the wave
equation (consistent with Q\ref{q:waveexp}(c)).

\step{Starting values}
The scheme is therefore exact provided \(U^0\) and \(U^1\) are exact.
\(U^0=f\) is exact; \(U_m^1=\frac12(f_{m+1}+f_{m-1})+kg_m\) (the formula
of Q\ref{q:waveexp} with \(r=1\)), while d'Alembert gives
\(u(x,k)=\frac12[f(x+h)+f(x-h)]+\frac12\int_{x-h}^{x+h}g\,d\xi\).  These
agree except that the integral is replaced by the mid-point value
\(\frac12\cdot2h\,g(x)=kg_m\), an error \(O(h^3)\).  Telescoping along the
characteristics, the numerical solution equals d'Alembert's formula
with the integral replaced by a mid-point sum, which tends to the exact
solution as \(h\to0\).
\end{solution}

%---------------------------------------------------------------------
\Q{Implicit \(\theta\)-schemes: stability}{Stability analysis}{8}\label{q:thetawave}
\begin{qbox}
For
\(\delta_t^2U_m^n=r^2\delta_x^2\bigl[\theta U_m^{n+1}+(1-2\theta)U_m^n+\theta U_m^{n-1}\bigr]\),
\(0\le\theta\le1\), show that the scheme is unconditionally stable for
\(\theta\ge\frac14\), and stable for \(r^2\le\frac1{1-4\theta}\) when
\(0\le\theta<\frac14\).  Show that \(\theta=\frac14\) gives
\(\bigl(1-\frac14r^2\delta_x^2\bigr)\delta_t^2U_m^n=r^2\delta_x^2U_m^n\).
\end{qbox}

\begin{rem}
\begin{itemize}
  \item \(\theta U^{n+1}+(1-2\theta)U^n+\theta U^{n-1}=U^n+\theta\delta_t^2U^n\).
  \item Characteristic equation \(a\xi^2-2b\xi+a=0\): product of roots
    \(1\); stable iff \(|b|\le a\).
  \item \(\theta=0\) is the explicit scheme (\(r\le1\)).
\end{itemize}
\end{rem}

\begin{solution}
With \(U_m^n=\xi^ne^{i\beta mh}\), \(S=\sin^2\frac{\beta h}2\):
\[
  \xi-2+\xi^{-1}=-4r^2S\bigl[\theta\xi+(1-2\theta)+\theta\xi^{-1}\bigr],
\]
\[
  \bigl(1+4r^2S\theta\bigr)\xi^2-2\bigl(1-2r^2S(1-2\theta)\bigr)\xi+\bigl(1+4r^2S\theta\bigr)=0 .
\]
The product of the roots is \(1\), so stability requires complex (or
equal) roots: \(|1-2r^2S(1-2\theta)|\le1+4r^2S\theta\).  The inequality
\(1-2r^2S(1-2\theta)\le1+4r^2S\theta\) always holds; the other,
\(-1+2r^2S(1-2\theta)\le1+4r^2S\theta\), is
\[
  r^2S(1-4\theta)\le1 .
\]
For \(\theta\ge\frac14\) it holds for all \(r\): unconditional stability.
For \(\theta<\frac14\) the worst case is \(S=1\):
\(r^2\le\frac1{1-4\theta}\).

\step{\(\theta=\frac14\)}
\(\frac14U^{n+1}+\frac12U^n+\frac14U^{n-1}=U^n+\frac14\delta_t^2U^n\), so the
scheme is \(\delta_t^2U=r^2\delta_x^2\bigl(U^n+\frac14\delta_t^2U^n\bigr)\),
i.e.\ \(\bigl(1-\frac14r^2\delta_x^2\bigr)\delta_t^2U_m^n=r^2\delta_x^2U_m^n\)
(the von Neumann scheme).
\end{solution}

%---------------------------------------------------------------------
\Q{Consistency of the \(\theta=\frac14\) scheme written out}{Analysis}{6}\label{q:quartercons}
\begin{qbox}
(Mitchell--Griffiths, Ch.~4, Exercise 11.)  Show that
\begin{align*}
  -(U_{m+1}^{n+1}+U_{m-1}^{n+1})+2\Bigl(1+\frac2{r^2}\Bigr)U_m^{n+1}
  &=2(U_{m+1}^n+U_{m-1}^n)-4\Bigl(1-\frac2{r^2}\Bigr)U_m^n\\
  &\quad+(U_{m+1}^{n-1}+U_{m-1}^{n-1})-2\Bigl(1+\frac2{r^2}\Bigr)U_m^{n-1}
\end{align*}
is a consistent approximation to \(u_{tt}=u_{xx}\).
\end{qbox}

\begin{rem}
\begin{itemize}
  \item \(\theta=\frac14\) scheme:
    \(\delta_t^2U=\frac{r^2}4\delta_x^2\bigl(U^{n+1}+2U^n+U^{n-1}\bigr)\).
  \item \(\frac12(u^{n+1}+u^{n-1})=u+\frac{k^2}2u_{tt}+\cdots\).
\end{itemize}
\end{rem}

\begin{solution}
\step{It is the \(\theta=\frac14\) scheme}
Multiply \(\delta_t^2U=\frac{r^2}4\delta_x^2(U^{n+1}+2U^n+U^{n-1})\) by
\(\frac4{r^2}\) and collect each time level:
\[
  -\delta_x^2U^{n+1}+\tfrac4{r^2}U^{n+1}=2\delta_x^2U^n+\tfrac8{r^2}U^n+\delta_x^2U^{n-1}-\tfrac4{r^2}U^{n-1}.
\]
Writing \(\delta_x^2U_m=U_{m+1}+U_{m-1}-2U_m\): the left side is
\(-(U_{m+1}^{n+1}+U_{m-1}^{n+1})+2(1+\frac2{r^2})U_m^{n+1}\); the right side
is \(2(U_{m+1}^n+U_{m-1}^n)-4(1-\frac2{r^2})U_m^n+(U_{m+1}^{n-1}+U_{m-1}^{n-1})-2(1+\frac2{r^2})U_m^{n-1}\).
This is the given formula.

\step{Truncation error}
\(\delta_t^2u=k^2u_{tt}+\frac{k^4}{12}u_{tttt}+\cdots\), and
\(\frac{r^2}4\delta_x^2(u^{n+1}+2u^n+u^{n-1})=r^2\delta_x^2\bigl(u+\frac{k^2}4u_{tt}+\cdots\bigr)
=k^2u_{xx}+\frac{k^2h^2}{12}u_{xxxx}+\frac{k^4}4u_{xxtt}+\cdots\).  With
\(u_{tt}=u_{xx}\) (so \(u_{tttt}=u_{xxtt}=u_{xxxx}\)):
\[
  T=k^2(u_{tt}-u_{xx})+\Bigl(\frac{k^4}{12}-\frac{k^4}4-\frac{k^2h^2}{12}\Bigr)u_{xxxx}+\cdots
  =-\frac{k^2}{12}\bigl(2k^2+h^2\bigr)u_{xxxx}+\cdots .
\]
\(T/k^2=O(k^2+h^2)\to0\): consistent (and unconditionally stable, by
Q\ref{q:thetawave}).
\end{solution}

%---------------------------------------------------------------------
\Q{One step of the implicit scheme with \(\theta=\frac12\)}{Numerical}{6}\label{q:wavetheta}
\begin{qbox}
Solve \(u_{tt}=u_{xx}\), \(u(x,0)=\sin\pi x\), \(u_t(x,0)=0\),
\(u(0,t)=u(1,t)=0\), by the implicit scheme with \(\theta=\frac12\),
\(h=\frac14\), \(r=k/h=1\), for one time step.  Compare with the exact
solution.
\end{qbox}

\begin{rem}
\begin{itemize}
  \item \(\theta=\frac12\): \(\delta_t^2U=\frac{r^2}2\delta_x^2(U^{n+1}+U^{n-1})\);
    unconditionally stable.
  \item \(u_t(x,0)=0\Rightarrow U^{-1}=U^1\); then
    \(\bigl(1-\frac{r^2}2\delta_x^2\bigr)U^1=U^0\).
\end{itemize}
\end{rem}

\begin{solution}
\step{The first step}
With \(n=0\) and \(U_m^{-1}=U_m^1\):
\(2U_m^1-2U_m^0=\frac{r^2}2\delta_x^2(2U_m^1)\), i.e.
\[
  (1+r^2)U_m^1-\frac{r^2}2\bigl(U_{m+1}^1+U_{m-1}^1\bigr)=U_m^0 .
\]
With \(r=1\): \(2U_m^1-\frac12(U_{m+1}^1+U_{m-1}^1)=U_m^0\), \(m=1,2,3\).

\step{Solving}
By symmetry \(U_1^1=U_3^1=a\), \(U_2^1=b\):
\[
  2a-\tfrac12b=0.707107,\qquad 2b-a=1 .
\]
So \(b=\frac{1+a}2\), \(\frac74a=0.957107\):
\[
  a=0.546918,\qquad b=0.773459 .
\]
(Check: \(U^1=\xi\,U^0\) with \(\xi=\frac1{1+2r^2\sin^2(\pi/8)}=0.773459\).)

\step{Comparison}
\(t=k=\frac14\): exact \(\cos\frac\pi4\sin\pi x\), i.e.\ \(0.5\) and
\(0.707107\).  The errors (\(0.047\), \(0.066\)) are large because
\(k=\frac14\) is a large step for a solution with period \(2\); the
scheme is stable but its phase error grows with \(k\).
\end{solution}

%---------------------------------------------------------------------
\Q{The wave equation in two space dimensions}{Stability analysis}{8}\label{q:wave2d}
\begin{qbox}
For \(u_{tt}=u_{x_1x_1}+u_{x_2x_2}\) with \(h_1=h_2=h\), \(r=k/h\):
\begin{enumerate}[label=(\alph*)]
  \item show that the explicit scheme
    \(\delta_t^2U=r^2(\delta_{x_1}^2+\delta_{x_2}^2)U\) is stable for
    \(r\le\frac1{\sqrt2}\);
  \item state and prove the stability conditions of the implicit
    scheme
    \(\bigl[1-\theta r^2(\delta_{x_1}^2+\delta_{x_2}^2)\bigr]\delta_t^2U=r^2(\delta_{x_1}^2+\delta_{x_2}^2)U\).
\end{enumerate}
\end{qbox}

\begin{rem}
\begin{itemize}
  \item Explicit:
    \(U^{n+1}_{l,m}=2(1-2r^2)U_{l,m}^n+r^2(U_{l\pm1,m}^n+U_{l,m\pm1}^n)-U_{l,m}^{n-1}\),
    order \(O(k^2+h^2)\).
  \item On a mode, \(\delta_{x_1}^2+\delta_{x_2}^2\to-4(S_1+S_2)\),
    \(S_1+S_2\le2\).
  \item Implicit: \(\theta\ge\frac14\) unconditional;
    \(r^2\le\frac1{2(1-4\theta)}\) for \(0\le\theta<\frac14\).
\end{itemize}
\end{rem}

\begin{solution}
Let \(U_{l,m}^n=\xi^ne^{i(\beta_1l+\beta_2m)h}\),
\(S_i=\sin^2\frac{\beta_ih}2\) and \(\Sigma=S_1+S_2\in[0,2]\), so that
\((\delta_{x_1}^2+\delta_{x_2}^2)\to-4\Sigma\) and \(\delta_t^2\to\xi-2+\xi^{-1}\).

\step{(a) Explicit}
\(\xi-2+\xi^{-1}=-4r^2\Sigma\), i.e.\
\[
  \xi^2-2\bigl(1-2r^2\Sigma\bigr)\xi+1=0 .
\]
The product of the roots is \(1\).  If the roots were real and distinct,
one would have modulus greater than \(1\); so stability requires
complex (or equal) roots, \(|1-2r^2\Sigma|\le1\iff r^2\Sigma\le1\), and then
\(|\xi_1|=|\xi_2|=1\).  This must hold for all modes; since \(\Sigma\) can be
\(2\), we need \(2r^2\le1\):
\[
  r\le\frac1{\sqrt2}.
\]

\step{(b) Implicit}
\(\bigl(1+4\theta r^2\Sigma\bigr)(\xi-2+\xi^{-1})=-4r^2\Sigma\), i.e.\
\[
  \bigl(1+4\theta r^2\Sigma\bigr)\xi^2-2\bigl(1-2r^2\Sigma(1-2\theta)\bigr)\xi+\bigl(1+4\theta r^2\Sigma\bigr)=0 .
\]
Again the product of the roots is \(1\), so stability requires
\(|1-2r^2\Sigma(1-2\theta)|\le1+4\theta r^2\Sigma\).  The upper inequality
always holds; the lower one,
\(-1+2r^2\Sigma(1-2\theta)\le1+4\theta r^2\Sigma\), is
\[
  r^2\Sigma(1-4\theta)\le1 .
\]
For \(\theta\ge\frac14\) it holds for all \(r\): unconditionally stable.
For \(0\le\theta<\frac14\) the worst case \(\Sigma=2\) gives
\(r^2\le\frac1{2(1-4\theta)}\).  An ADI method can be developed for this
scheme by factorising the operator on the left.
\end{solution}

%---------------------------------------------------------------------
\Q{A nine-point explicit scheme for the two-dimensional wave equation}{Stability analysis}{8}\label{q:wave9}
\begin{qbox}
(Mitchell--Griffiths, Ch.~4, Exercise 12.)  Using the nine-point
replacement
\(\nabla^2u\approx\frac1{h^2}\bigl(\delta_{x_1}^2+\delta_{x_2}^2+\frac16\delta_{x_1}^2\delta_{x_2}^2\bigr)u\),
derive an explicit scheme for \(u_{tt}=\nabla^2u\), and determine its
stability and CFL conditions.
\end{qbox}

\begin{rem}
\begin{itemize}
  \item \(\delta_{x_1}^2\delta_{x_2}^2\) uses the four diagonal
    neighbours: the stencil is the full \(3\times3\) square.
  \item On a mode: \(\delta_{x_i}^2\to-4S_i\),
    \(\delta_{x_1}^2\delta_{x_2}^2\to16S_1S_2\).
  \item CFL: the numerical domain of dependence (here a square) must
    contain the exact one (a disc of radius \(t\)).
\end{itemize}
\end{rem}

\begin{solution}
\step{The scheme}
\(\delta_t^2U_{l,m}^n=r^2\bigl(\delta_{x_1}^2+\delta_{x_2}^2+\frac16\delta_{x_1}^2\delta_{x_2}^2\bigr)U_{l,m}^n\),
i.e.\ \(U^{n+1}=2U^n-U^{n-1}+r^2\bigl(\delta_{x_1}^2+\delta_{x_2}^2+\frac16\delta_{x_1}^2\delta_{x_2}^2\bigr)U^n\),
which involves the nine points \((l+i,m+j)\), \(i,j\in\{-1,0,1\}\), at
level \(n\).

\step{Stability}
\(\xi-2+\xi^{-1}=-r^2\Phi\), with
\(\Phi=4S_1+4S_2-\frac{16}6S_1S_2=4S_1+4S_2-\frac83S_1S_2\ge0\), i.e.\
\(\xi^2-(2-r^2\Phi)\xi+1=0\).  The product of the roots is \(1\), so
stability requires complex (or equal) roots:
\(|2-r^2\Phi|\le2\iff r^2\Phi\le4\), for all \(S_1,S_2\in[0,1]\).  Since \(\partial\Phi/\partial S_1=4-\frac83S_2>0\)
(and similarly for \(S_2\)), the maximum is at \(S_1=S_2=1\):
\(\Phi_{\max}=8-\frac83=\frac{16}3\).  Hence
\[
  r^2\le\frac34,\qquad \boxed{r\le\frac{\sqrt3}2\approx0.866}.
\]
\step{CFL}
After \(n\) steps the scheme uses data in the square
\(|x_1-X_1|,|x_2-X_2|\le nh\); the exact solution depends on the disc of
radius \(nk\).  The disc lies in the square iff \(k\le h\): CFL
condition \(r\le1\).  Stability (\(r\le\frac{\sqrt3}2\)) is the stricter
requirement.  (For the five-point scheme the numerical domain is the
diamond \(|x_1|+|x_2|\le nh\), so CFL gives \(r\le\frac1{\sqrt2}\), the same as
its stability condition; the nine-point scheme allows a larger time
step.)
\end{solution}

%#####################################################################
\Needspace{32\baselineskip}
\section{The Heat Equation: Comparison of Methods}
%#####################################################################

%---------------------------------------------------------------------
\Q{Four methods for the heat equation}{Numerical}{10}\label{q:heat4}
\begin{qbox}
Solve \(u_t=u_{xx}\), \(u(x,0)=\sin\pi x\), \(u(0,t)=u(1,t)=0\), with
\(h=\frac14\), \(k=\frac1{64}\) (\(r=\frac14\)), up to \(t=2k\), by
(i) Crank--Nicolson, (ii) the explicit (Schmidt) method, (iii) Du
Fort--Frankel (with \(U^1\) from the explicit method), and (iv) the
Douglas method.  Compare the values at \(x=\frac12\) with the exact
\(e^{-\pi^2t}\sin\pi x\).
\end{qbox}

\begin{rem}
\begin{itemize}
  \item Explicit \(U^{n+1}=(1+r\delta^2)U^n\), \(O(k+h^2)\), \(r\le\frac12\).
  \item Crank--Nicolson
    \((1-\frac r2\delta^2)U^{n+1}=(1+\frac r2\delta^2)U^n\), \(O(k^2+h^2)\),
    unconditionally stable.
  \item Du Fort--Frankel
    \((1+2r)U_m^{n+1}=2r(U_{m+1}^n+U_{m-1}^n)+(1-2r)U_m^{n-1}\),
    \(O(k^2+h^2+k^2/h^2)\), unconditionally stable.
  \item Douglas
    \(\bigl[1-\frac12(r-\frac16)\delta^2\bigr]U^{n+1}=\bigl[1+\frac12(r+\frac16)\delta^2\bigr]U^n\),
    \(O(k^2+h^4)\), unconditionally stable.
  \item For \(\sin\pi x\) data: \(\delta^2\to-4s\), \(s=\sin^2\frac{\pi h}2\);
    two-level schemes give \(U_m^n=\xi^n\sin\pi x_m\).
\end{itemize}
\end{rem}

\begin{solution}
Interior points \(x=\frac14,\frac12,\frac34\);
\(U^0=(0.707107,1,0.707107)\); \(s=\sin^2\frac\pi8=0.146447\), so
\(4rs=0.146447\).

\step{(ii) Explicit}
\(\xi=1-4rs=0.853553\) (Q\ref{q:old18}):
\[
  U^1=(0.603553,0.853553,0.603553),\qquad
  U^2=(0.515165,0.728553,0.515165).
\]

\step{(i) Crank--Nicolson}
With \(U_1=U_3=a\), \(U_2=b\), the first step
\((1+r)U_m^1-\frac r2(U_{m+1}^1+U_{m-1}^1)=(1-r)U_m^0+\frac r2(U_{m+1}^0+U_{m-1}^0)\)
reads
\[
  1.25a-0.125b=0.655330,\qquad -0.25a+1.25b=0.926777,
\]
so \(a=0.610619\), \(b=0.863545\) (determinant \(1.53125\)).  This agrees
with \(\xi=\frac{1-2rs}{1+2rs}=0.863545\).  The second step gives
\(U^2=(0.527297,0.745710,0.527297)\).

\step{(iv) Douglas}
\(\frac12(r-\frac16)=\frac1{24}\), \(\frac12(r+\frac16)=\frac5{24}\):
\[
  \xi=\frac{1-\frac{20}{24}s}{1+\frac4{24}s}=\frac{0.8779612}{1.0244078}=0.857043,
\]
\(U^1=(0.606021,0.857043,0.606021)\),
\(U^2=(0.519386,0.734522,0.519386)\).

\step{(iii) Du Fort--Frankel}
With \(r=\frac14\): \(1.5U_m^2=0.5(U_{m+1}^1+U_{m-1}^1)+0.5U_m^0\), using
the explicit \(U^1\):
\[
  U_2^2=\frac{0.5(1.207107)+0.5}{1.5}=0.735702,\qquad
  U_1^2=\frac{0.5(0.853553)+0.5(0.707107)}{1.5}=0.520220 .
\]

\step{Comparison at \(x=\frac12\), \(t=\frac1{32}\)} (exact \(0.734603\))
\begin{center}
\begin{tabular}{@{}lcc@{}}
\toprule
Method & \(U\) & error\\
\midrule
Explicit & 0.728553 & \(6.0\times10^{-3}\)\\
Crank--Nicolson & 0.745710 & \(1.1\times10^{-2}\)\\
Du Fort--Frankel & 0.735702 & \(1.1\times10^{-3}\)\\
Douglas & 0.734522 & \(8.1\times10^{-5}\)\\
\bottomrule
\end{tabular}
\end{center}
Douglas (\(O(k^2+h^4)\)) is by far the most accurate.  Crank--Nicolson
is worse than the explicit scheme here: both carry the spatial error of
\(\delta^2\) (which slows the decay), but for the explicit scheme with
\(r=\frac14\) the time error (which speeds it up) partly cancels it.
\end{solution}

%---------------------------------------------------------------------
\Q{Crank--Nicolson versus Douglas at the optimal \(r\)}{Numerical / analysis}{10}\label{q:mgdouglas}
\begin{qbox}
(Mitchell--Griffiths, Ch.~2, Example 2 and Exercise 7.)  Solve
\(u_t=u_{xx}\), \(u(x,0)=\sin x\) (\(0\le x\le\pi\)), \(u(0,t)=u(\pi,t)=0\),
by Crank--Nicolson and by the Douglas formula, with \(h=\pi/20\) and
\(r=1/\sqrt{20}\).
\begin{enumerate}[label=(\alph*)]
  \item Show that the truncation error of the Douglas formula is
    \(\frac{kh^4}{240}(1-20r^2)u_{xxxxxx}+\cdots\), so that
    \(r=1/\sqrt{20}\) minimises it.
  \item Show that both formulas are unconditionally stable.
  \item Compute \(U\) at \(x=\frac\pi2\) after \(10\) steps and compare
    with the exact solution \(e^{-t}\sin x\).
\end{enumerate}
\end{qbox}

\begin{rem}
\begin{itemize}
  \item Douglas:
    \(\bigl[1-\frac12(r-\frac16)\delta^2\bigr]U^{n+1}=\bigl[1+\frac12(r+\frac16)\delta^2\bigr]U^n\);
    \(\xi=\dfrac{1-2(r+\frac16)S}{1+2(r-\frac16)S}\).
  \item \(\sin x_m\) is an eigenvector of \(\delta^2\):
    \(\delta^2\sin x_m=-4\sin^2\frac h2\,\sin x_m\).
  \item \(u_{tt}=u_{xxxx}\), \(u_{ttt}=u_{xxxxxx}\).
\end{itemize}
\end{rem}

\begin{solution}
\step{(a) Truncation error}
Expanding
\(T=u^{n+1}-u^n-\frac12(r-\frac16)\delta^2u^{n+1}-\frac12(r+\frac16)\delta^2u^n\)
about \((mh,nk)\) and using the heat equation to replace time
derivatives, the terms up to order \(kh^2\) and \(k^2\) cancel (this is
the purpose of the \(\pm\frac16\)); the first surviving term is
\[
  T=\frac{kh^4}{240}\bigl(1-20r^2\bigr)\frac{\partial^6u}{\partial x^6}+\cdots .
\]
(Checked by a symbolic expansion.)  So \(T/k=O(h^4)\) in general, and
for \(r^2=\frac1{20}\) the leading term vanishes and the error is of
still higher order: \(r=1/\sqrt{20}\) gives the minimum truncation error.

\step{(b) Stability}
Crank--Nicolson: \(\xi=\frac{1-2rS}{1+2rS}\in(-1,1]\) for \(r>0\).
Douglas: the denominator \(1+2(r-\frac16)S\ge1-\frac S3>0\).  Then
\(\xi\le1\iff-2(r+\frac16)S\le2(r-\frac16)S\iff r\ge0\), and
\(\xi\ge-1\iff2\ge\bigl[2(r+\frac16)-2(r-\frac16)\bigr]S=\frac23S\), true.  So
both are stable for all \(r>0\).

\step{(c) Numbers}
\(h=\frac\pi{20}\), \(k=rh^2=0.0055173\), \(S=\sin^2\frac\pi{40}=0.0061558\).
\[
  \xi_{CN}=\frac{1-2rS}{1+2rS}=0.9945092,\qquad
  \xi_{D}=\frac{1-2(r+\frac16)S}{1+2(r-\frac16)S}=0.9944979,
  \qquad e^{-k}=0.9944979 .
\]
After \(10\) steps (\(t=0.055173\)), at \(x=\frac\pi2\):
\begin{center}
\begin{tabular}{@{}lcc@{}}
\toprule
 & \(U(\frac\pi2)\) & error\\
\midrule
Crank--Nicolson & 0.9464288 & \(1.1\times10^{-4}\)\\
Douglas & 0.9463216 & \(1.3\times10^{-10}\)\\
exact \(e^{-t}\) & 0.9463216 & --\\
\bottomrule
\end{tabular}
\end{center}
At the optimal \(r\) the Douglas formula is essentially exact for this
smooth solution.
\end{solution}

%---------------------------------------------------------------------
\Q{Two-dimensional ADI forms of the Douglas scheme}{Analysis}{8}\label{q:mfcons}
\begin{qbox}
(Lecture notes; Mitchell--Griffiths, Ch.~2, Exercise 14, second part.)
For \(u_t=u_{x_1x_1}+u_{x_2x_2}\), let \(a=\frac12(r-\frac16)\),
\(b=\frac12(r+\frac16)\), and consider the split (Mitchell--Fairweather,
or high-order Douglas, ADI) scheme
\[
  (1-a\delta_{x_1}^2)U^{n+1*}=(1+b\delta_{x_2}^2)U^n,\qquad
  (1-a\delta_{x_2}^2)U^{n+1}=(1+b\delta_{x_1}^2)U^{n+1*}.
\]
\begin{enumerate}[label=(\alph*)]
  \item Show that eliminating \(U^{n+1*}\) gives the factored
    Douglas scheme
    \((1-a\delta_{x_1}^2)(1-a\delta_{x_2}^2)U^{n+1}=(1+b\delta_{x_1}^2)(1+b\delta_{x_2}^2)U^n\).
  \item Show that, with \(U^{n+1*}=U^{n+1/2}\), the individual
    formulas are \emph{not} consistent (contrast Q\ref{q:prhalf}).
\end{enumerate}
\end{qbox}

\begin{rem}
\begin{itemize}
  \item Eliminating the intermediate value needs
    \(\delta_{x_1}^2\delta_{x_2}^2=\delta_{x_2}^2\delta_{x_1}^2\).
  \item \(a\,\delta^2u=\frac k2u_{xx}-\frac{h^2}{12}u_{xx}+\cdots\) and
    \(b\,\delta^2u=\frac k2u_{xx}+\frac{h^2}{12}u_{xx}+\cdots\).
  \item In the notes this split is labelled ``Douglas--Rachford ADI'';
    the genuine Douglas--Rachford split is that of Q\ref{q:dr}.
\end{itemize}
\end{rem}

\begin{solution}
\step{(a) Elimination}
\(U^{n+1*}=(1-a\delta_{x_1}^2)^{-1}(1+b\delta_{x_2}^2)U^n\), so
\((1-a\delta_{x_2}^2)U^{n+1}=(1+b\delta_{x_1}^2)(1-a\delta_{x_1}^2)^{-1}(1+b\delta_{x_2}^2)U^n\).
Applying \((1-a\delta_{x_1}^2)\), which commutes with the other
operators,
\((1-a\delta_{x_1}^2)(1-a\delta_{x_2}^2)U^{n+1}=(1+b\delta_{x_1}^2)(1+b\delta_{x_2}^2)U^n\).

\step{(b) The half-steps}
First formula with \(u^{n+1/2}\), expanded about \((lh,mh,nk)\):
\[
  T_1=u^{n+\frac12}-u^n-a\delta_{x_1}^2u^{n+\frac12}-b\delta_{x_2}^2u^n
  =\frac k2\bigl(u_t-\nabla^2u\bigr)+\frac{h^2}{12}\bigl(u_{x_1x_1}-u_{x_2x_2}\bigr)+\cdots .
\]
Dividing by the time step \(\frac k2\):
\(\frac{T_1}{k/2}\to\frac{h^2}{6k}\bigl(u_{x_1x_1}-u_{x_2x_2}\bigr)=\frac1{6r}\bigl(u_{x_1x_1}-u_{x_2x_2}\bigr)\ne0\)
as \(h,k\to0\) with \(r\) fixed: not consistent.  The second formula
gives the same term with the opposite sign,
\(-\frac{h^2}{12}(u_{x_1x_1}-u_{x_2x_2})\); the two errors cancel when the
half-steps are combined, so the full step is accurate even though
neither half is a consistent approximation on its own.
\end{solution}

%---------------------------------------------------------------------
\Q{Short answers on Note 6}{Short answer}{8}\label{q:short6}
\begin{qbox}
Answer briefly, with reasons.
\begin{enumerate}[label=(\alph*)]
  \item In the ADI iteration for Laplace's equation, what is the role
    of \(\lambda\)?  How should it be chosen?
  \item Why is the discrete Neumann problem singular while the Robin
    problem is not?
  \item Is a scheme that satisfies the CFL condition necessarily
    stable?
  \item Why does the Du Fort--Frankel scheme need care although it is
    unconditionally stable?
  \item Why is a three-level scheme natural for the wave equation, and
    what extra data does it need?
  \item Which of Peaceman--Rachford, Douglas--Rachford and the Douglas
    scheme are second order in time?
\end{enumerate}
\end{qbox}

\begin{rem}
\begin{itemize}
  \item ADI iteration converges for every \(\lambda>0\); optimum
    \(\frac2\lambda=\sqrt{\mu_{\min}\mu_{\max}}\).
  \item CFL is necessary, not sufficient, for convergence.
  \item Orders: PR \(O(k^2+h^2)\), DR \(O(k+h^2)\), Douglas
    \(O(k^2+h^4)\).
\end{itemize}
\end{rem}

\begin{solution}
\step{(a)}
\(\lambda\) is an iteration (acceleration) parameter, no longer
\(k/h^2\).  The iteration converges for every \(\lambda>0\)
(Q\ref{q:priter}); the best fixed value is
\(\frac2\lambda=\sqrt{\mu_{\min}\mu_{\max}}\), i.e.\ \(\lambda=\csc\frac\pi M\) on
the unit square (Q\ref{q:optpar}); varying \(\lam\) (Wachspress) is
better still.

\step{(b)}
For Neumann data every row of \(A\) sums to zero, so \(A\mathbf1=0\):
the solution is fixed only up to a constant, and the data must satisfy
a compatibility condition.  The Robin terms \(hp_i,hq_i>0\) add to the
boundary diagonal entries, making \(A\) positive definite in the
weighted sense, hence non-singular (Q\ref{q:robin}).

\step{(c)}
No.  CFL is necessary but not sufficient.  For the nine-point wave
scheme, CFL allows \(r\le1\), but stability needs \(r\le\frac{\sqrt3}2\)
(Q\ref{q:wave9}).

\step{(d)}
It is consistent with \(u_t=u_{xx}\) only if \(k/h\to0\); if \(k/h\) is
fixed it approximates \(u_t+c^2u_{tt}=u_{xx}\) (Q\ref{q:old32}).  So
large steps give stable but wrong answers.

\step{(e)}
The equation is second order in \(t\), and \(\delta_t^2\) needs three
levels.  The two initial conditions \(u(x,0)=f\) and \(u_t(x,0)=g\)
provide \(U^0\) and, through \(U^{-1}=U^1-2kg\), the level \(U^1\).

\step{(f)}
Peaceman--Rachford (\(O(k^2+h^2)\)) and Douglas (\(O(k^2+h^4)\)) are
second order in time; Douglas--Rachford is a perturbation of the fully
implicit scheme and is only first order (\(O(k^2+kh^2)\) truncation
error, i.e.\ \(O(k+h^2)\) accuracy).
\end{solution}


\end{document}
